% Môn: Vật lý
% Lớp: 12
% Chương: Chương III. Từ trường
% Bài: L12C3 Bài 14. Từ trường
% Số câu: 284
% App đọc file này trực tiếp — sửa rồi Commit trên GitHub, bấm Đồng bộ GitHub .tex

% L12C3 Bài 14. Từ trường

% ===== Câu 1 =====
% ID: L12C3B14-102-TN
% Mức: VD
\dangbt{Từ trường của dòng điện thẳng, dòng điện tròn và ống dây}
\begin{ex}
% ID: L12C3B14-102-TN
% Mức: VD
Kết luận đúng khi vận dụng từ trường của dòng điện thẳng, dòng điện tròn và ống dây là
\choice
{Từ trường trong lòng ống dây dài luôn bằng không.}
{\True Đường sức quanh dây dẫn thẳng là các đường tròn đồng tâm; từ trường trong lòng ống dây dài gần đều.}
{Có thể bỏ qua phương, chiều và đơn vị của các đại lượng trong mọi trường hợp.}
{Đại lượng đang xét không phụ thuộc các điều kiện vật lí nêu trong bài.}
\loigiai{
Đường sức quanh dây dẫn thẳng là các đường tròn đồng tâm; từ trường trong lòng ống dây dài gần đều. Vì vậy phương án $B$ đúng; các phương án còn lại trái với định nghĩa, định luật hoặc điều kiện áp dụng.
}
\end{ex}

% ===== Câu 2 =====
% ID: L12C3B14-103-TN
% Mức: NB
\dangbt{So sánh, tổng hợp và biểu diễn các từ trường đơn giản}
\begin{ex}
% ID: L12C3B14-103-TN
% Mức: NB
Phát biểu nào sau đây đúng về so sánh, tổng hợp và biểu diễn các từ trường đơn giản?
\choice
{\True Vectơ cảm ứng từ tổng hợp tại một điểm bằng tổng vectơ các cảm ứng từ thành phần.}
{Độ lớn cảm ứng từ tổng hợp luôn bằng tổng độ lớn các cảm ứng từ thành phần.}
{Đại lượng đang xét không phụ thuộc các điều kiện vật lí nêu trong bài.}
{Có thể bỏ qua phương, chiều và đơn vị của các đại lượng trong mọi trường hợp.}
\loigiai{
Vectơ cảm ứng từ tổng hợp tại một điểm bằng tổng vectơ các cảm ứng từ thành phần. Vì vậy phương án $A$ đúng; các phương án còn lại trái với định nghĩa, định luật hoặc điều kiện áp dụng.
}
\end{ex}

% ===== Câu 3 =====
% ID: L12C3B14-104-TN
% Mức: NB
\begin{ex}
% ID: L12C3B14-104-TN
% Mức: NB
Trong so sánh, tổng hợp và biểu diễn các từ trường đơn giản, nhận định nào phù hợp với kiến thức Vật lí?
\choice
{Độ lớn cảm ứng từ tổng hợp luôn bằng tổng độ lớn các cảm ứng từ thành phần.}
{\True Vectơ cảm ứng từ tổng hợp tại một điểm bằng tổng vectơ các cảm ứng từ thành phần.}
{Có thể bỏ qua phương, chiều và đơn vị của các đại lượng trong mọi trường hợp.}
{Đại lượng đang xét không phụ thuộc các điều kiện vật lí nêu trong bài.}
\loigiai{
Vectơ cảm ứng từ tổng hợp tại một điểm bằng tổng vectơ các cảm ứng từ thành phần. Vì vậy phương án $B$ đúng; các phương án còn lại trái với định nghĩa, định luật hoặc điều kiện áp dụng.
}
\end{ex}

% ===== Câu 4 =====
% ID: L12C3B14-105-TN
% Mức: NB
\begin{ex}
% ID: L12C3B14-105-TN
% Mức: NB
Tình huống 3: chọn kết luận chính xác về so sánh, tổng hợp và biểu diễn các từ trường đơn giản.
\choice
{Đại lượng đang xét không phụ thuộc các điều kiện vật lí nêu trong bài.}
{Có thể bỏ qua phương, chiều và đơn vị của các đại lượng trong mọi trường hợp.}
{\True Vectơ cảm ứng từ tổng hợp tại một điểm bằng tổng vectơ các cảm ứng từ thành phần.}
{Độ lớn cảm ứng từ tổng hợp luôn bằng tổng độ lớn các cảm ứng từ thành phần.}
\loigiai{
Vectơ cảm ứng từ tổng hợp tại một điểm bằng tổng vectơ các cảm ứng từ thành phần. Vì vậy phương án $C$ đúng; các phương án còn lại trái với định nghĩa, định luật hoặc điều kiện áp dụng.
}
\end{ex}

% ===== Câu 5 =====
% ID: L12C3B14-106-TN
% Mức: TH
\begin{ex}
% ID: L12C3B14-106-TN
% Mức: TH
Nội dung nào mô tả đúng nhất so sánh, tổng hợp và biểu diễn các từ trường đơn giản?
\choice
{Có thể bỏ qua phương, chiều và đơn vị của các đại lượng trong mọi trường hợp.}
{Đại lượng đang xét không phụ thuộc các điều kiện vật lí nêu trong bài.}
{Độ lớn cảm ứng từ tổng hợp luôn bằng tổng độ lớn các cảm ứng từ thành phần.}
{\True Vectơ cảm ứng từ tổng hợp tại một điểm bằng tổng vectơ các cảm ứng từ thành phần.}
\loigiai{
Vectơ cảm ứng từ tổng hợp tại một điểm bằng tổng vectơ các cảm ứng từ thành phần. Vì vậy phương án $D$ đúng; các phương án còn lại trái với định nghĩa, định luật hoặc điều kiện áp dụng.
}
\end{ex}

% ===== Câu 6 =====
% ID: L12C3B14-107-TN
% Mức: TH
\begin{ex}
% ID: L12C3B14-107-TN
% Mức: TH
Khi phân tích so sánh, tổng hợp và biểu diễn các từ trường đơn giản, phát biểu nào đúng?
\choice
{\True Vectơ cảm ứng từ tổng hợp tại một điểm bằng tổng vectơ các cảm ứng từ thành phần.}
{Độ lớn cảm ứng từ tổng hợp luôn bằng tổng độ lớn các cảm ứng từ thành phần.}
{Đại lượng đang xét không phụ thuộc các điều kiện vật lí nêu trong bài.}
{Có thể bỏ qua phương, chiều và đơn vị của các đại lượng trong mọi trường hợp.}
\loigiai{
Vectơ cảm ứng từ tổng hợp tại một điểm bằng tổng vectơ các cảm ứng từ thành phần. Vì vậy phương án $A$ đúng; các phương án còn lại trái với định nghĩa, định luật hoặc điều kiện áp dụng.
}
\end{ex}

% ===== Câu 7 =====
% ID: L12C3B14-108-TN
% Mức: TH
\begin{ex}
% ID: L12C3B14-108-TN
% Mức: TH
Một học sinh ôn tập so sánh, tổng hợp và biểu diễn các từ trường đơn giản. Kết luận nào của học sinh là đúng?
\choice
{Độ lớn cảm ứng từ tổng hợp luôn bằng tổng độ lớn các cảm ứng từ thành phần.}
{\True Vectơ cảm ứng từ tổng hợp tại một điểm bằng tổng vectơ các cảm ứng từ thành phần.}
{Có thể bỏ qua phương, chiều và đơn vị của các đại lượng trong mọi trường hợp.}
{Đại lượng đang xét không phụ thuộc các điều kiện vật lí nêu trong bài.}
\loigiai{
Vectơ cảm ứng từ tổng hợp tại một điểm bằng tổng vectơ các cảm ứng từ thành phần. Vì vậy phương án $B$ đúng; các phương án còn lại trái với định nghĩa, định luật hoặc điều kiện áp dụng.
}
\end{ex}

% ===== Câu 8 =====
% ID: L12C3B14-109-TN
% Mức: TH
\begin{ex}
% ID: L12C3B14-109-TN
% Mức: TH
Chọn phát biểu không trái với kiến thức về so sánh, tổng hợp và biểu diễn các từ trường đơn giản.
\choice
{Đại lượng đang xét không phụ thuộc các điều kiện vật lí nêu trong bài.}
{Có thể bỏ qua phương, chiều và đơn vị của các đại lượng trong mọi trường hợp.}
{\True Vectơ cảm ứng từ tổng hợp tại một điểm bằng tổng vectơ các cảm ứng từ thành phần.}
{Độ lớn cảm ứng từ tổng hợp luôn bằng tổng độ lớn các cảm ứng từ thành phần.}
\loigiai{
Vectơ cảm ứng từ tổng hợp tại một điểm bằng tổng vectơ các cảm ứng từ thành phần. Vì vậy phương án $C$ đúng; các phương án còn lại trái với định nghĩa, định luật hoặc điều kiện áp dụng.
}
\end{ex}

% ===== Câu 9 =====
% ID: L12C3B14-11-DS
% Mức: TH
\dangbt{Nhận biết từ trường, nguồn gây ra từ trường và tương tác từ}
\begin{ex}
% ID: L12C3B14-11-DS
% Mức: TH
Xét nhận biết từ trường, nguồn gây ra từ trường và tương tác từ (Tình huống 1). Hãy xác định tính đúng, sai của bốn phát biểu sau.
\choiceTF
{\True Từ trường tồn tại xung quanh nam châm và dòng điện; nó gây lực từ lên nam châm hoặc dòng điện khác đặt trong đó.}
{Mọi điện tích đứng yên đều tạo ra từ trường giống dòng điện.}
{\True Khi giải bài thuộc dạng này phải xác định đúng phương, chiều, điều kiện áp dụng và đơn vị của các đại lượng.}
{Trong dạng bài này, kết quả không phụ thuộc dữ kiện và có thể bỏ qua điều kiện áp dụng.}
\loigiai{
\begin{itemchoice}
\item (Đúng) Từ trường tồn tại xung quanh nam châm và dòng điện; nó gây lực từ lên nam châm hoặc dòng điện khác đặt trong đó. (Vì): Từ trường tồn tại xung quanh nam châm và dòng điện; nó gây lực từ lên nam châm hoặc dòng điện khác đặt trong đó. b) Sai: Mọi điện tích đứng yên đều tạo ra từ trường giống dòng điện. c) Đúng: Khi giải bài thuộc dạng này phải xác định đúng phương, chiều, điều kiện áp dụng và đơn vị của các đại lượng. d) Sai: Trong dạng bài này, kết quả không phụ thuộc dữ kiện và có thể bỏ qua điều kiện áp dụng. Kiến thức cốt lõi: Từ trường tồn tại xung quanh nam châm và dòng điện; nó gây lực từ lên nam châm hoặc dòng điện khác đặt trong đó
\item (Sai) Mọi điện tích đứng yên đều tạo ra từ trường giống dòng điện.
\item (Đúng) Khi giải bài thuộc dạng này phải xác định đúng phương, chiều, điều kiện áp dụng và đơn vị của các đại lượng.
\item (Sai) Trong dạng bài này, kết quả không phụ thuộc dữ kiện và có thể bỏ qua điều kiện áp dụng.
\end{itemchoice}
}
\end{ex}

% ===== Câu 10 =====
% ID: L12C3B14-11-TL
% Mức: VD
\dangbt{Mô tả đường sức từ và xác định chiều đường sức từ}
\begin{bt}
% ID: L12C3B14-11-TL
% Mức: VD
Nêu một tình huống thực tiễn liên quan đến mô tả đường sức từ và xác định chiều đường sức từ và giải thích bằng kiến thức Vật lí.
\loigiai{
Cần nêu được: Tiếp tuyến với đường sức từ tại một điểm cho phương của vectơ cảm ứng từ; bên ngoài nam châm, đường sức đi từ cực Bắc đến cực Nam. Khi vận dụng phải xác định đúng dữ kiện, phương chiều nếu có, điều kiện áp dụng, hệ đơn vị và kiểm tra tính hợp lí của kết quả. Nhận định “Các đường sức từ có thể cắt nhau tại một điểm.” sai vì trái với kiến thức nêu trên.
}
\end{bt}

% ===== Câu 11 =====
% ID: L12C3B14-11-TLN
% Mức: VD
\begin{ex}
% ID: L12C3B14-11-TLN
% Mức: VD
Trong bốn phát biểu sau về mô tả đường sức từ và xác định chiều đường sức từ, có bao nhiêu phát biểu đúng? (Chỉ ghi một số nguyên.) (1) Các đường sức từ có thể cắt nhau tại một điểm. (2) Cần đọc đúng dữ kiện và giả thiết. (3) Tiếp tuyến với đường sức từ tại một điểm cho phương của vectơ cảm ứng từ; bên ngoài nam châm, đường sức đi từ cực Bắc đến cực Nam. (4) Có thể thay số trước khi xác định công thức.
\shortans{2}
\loigiai{
Đối chiếu từng phát biểu với kiến thức: Tiếp tuyến với đường sức từ tại một điểm cho phương của vectơ cảm ứng từ; bên ngoài nam châm, đường sức đi từ cực Bắc đến cực Nam. Các phát biểu đúng có số lượng là 2.
}
\end{ex}

% ===== Câu 12 =====
% ID: L12C3B14-11-TN
% Mức: NB
\begin{ex}
% ID: L12C3B14-11-TN
% Mức: NB
Phát biểu nào sau đây đúng về mô tả đường sức từ và xác định chiều đường sức từ?
\choice
{\True Tiếp tuyến với đường sức từ tại một điểm cho phương của vectơ cảm ứng từ; bên ngoài nam châm, đường sức đi từ cực Bắc đến cực Nam.}
{Các đường sức từ có thể cắt nhau tại một điểm.}
{Đại lượng đang xét không phụ thuộc các điều kiện vật lí nêu trong bài.}
{Có thể bỏ qua phương, chiều và đơn vị của các đại lượng trong mọi trường hợp.}
\loigiai{
Tiếp tuyến với đường sức từ tại một điểm cho phương của vectơ cảm ứng từ; bên ngoài nam châm, đường sức đi từ cực Bắc đến cực Nam. Vì vậy phương án $A$ đúng; các phương án còn lại trái với định nghĩa, định luật hoặc điều kiện áp dụng.
}
\end{ex}

% ===== Câu 13 =====
% ID: L12C3B14-110-TN
% Mức: VD
\dangbt{So sánh, tổng hợp và biểu diễn các từ trường đơn giản}
\begin{ex}
% ID: L12C3B14-110-TN
% Mức: VD
Cơ sở Vật lí đúng dùng để giải so sánh, tổng hợp và biểu diễn các từ trường đơn giản là gì?
\choice
{Có thể bỏ qua phương, chiều và đơn vị của các đại lượng trong mọi trường hợp.}
{Đại lượng đang xét không phụ thuộc các điều kiện vật lí nêu trong bài.}
{Độ lớn cảm ứng từ tổng hợp luôn bằng tổng độ lớn các cảm ứng từ thành phần.}
{\True Vectơ cảm ứng từ tổng hợp tại một điểm bằng tổng vectơ các cảm ứng từ thành phần.}
\loigiai{
Vectơ cảm ứng từ tổng hợp tại một điểm bằng tổng vectơ các cảm ứng từ thành phần. Vì vậy phương án $D$ đúng; các phương án còn lại trái với định nghĩa, định luật hoặc điều kiện áp dụng.
}
\end{ex}

% ===== Câu 14 =====
% ID: L12C3B14-111-TN
% Mức: VD
\begin{ex}
% ID: L12C3B14-111-TN
% Mức: VD
Trong một bài thực hành về so sánh, tổng hợp và biểu diễn các từ trường đơn giản, nhận xét nào đúng?
\choice
{\True Vectơ cảm ứng từ tổng hợp tại một điểm bằng tổng vectơ các cảm ứng từ thành phần.}
{Độ lớn cảm ứng từ tổng hợp luôn bằng tổng độ lớn các cảm ứng từ thành phần.}
{Đại lượng đang xét không phụ thuộc các điều kiện vật lí nêu trong bài.}
{Có thể bỏ qua phương, chiều và đơn vị của các đại lượng trong mọi trường hợp.}
\loigiai{
Vectơ cảm ứng từ tổng hợp tại một điểm bằng tổng vectơ các cảm ứng từ thành phần. Vì vậy phương án $A$ đúng; các phương án còn lại trái với định nghĩa, định luật hoặc điều kiện áp dụng.
}
\end{ex}

% ===== Câu 15 =====
% ID: L12C3B14-112-TN
% Mức: VD
\begin{ex}
% ID: L12C3B14-112-TN
% Mức: VD
Kết luận đúng khi vận dụng so sánh, tổng hợp và biểu diễn các từ trường đơn giản là
\choice
{Độ lớn cảm ứng từ tổng hợp luôn bằng tổng độ lớn các cảm ứng từ thành phần.}
{\True Vectơ cảm ứng từ tổng hợp tại một điểm bằng tổng vectơ các cảm ứng từ thành phần.}
{Có thể bỏ qua phương, chiều và đơn vị của các đại lượng trong mọi trường hợp.}
{Đại lượng đang xét không phụ thuộc các điều kiện vật lí nêu trong bài.}
\loigiai{
Vectơ cảm ứng từ tổng hợp tại một điểm bằng tổng vectơ các cảm ứng từ thành phần. Vì vậy phương án $B$ đúng; các phương án còn lại trái với định nghĩa, định luật hoặc điều kiện áp dụng.
}
\end{ex}

% ===== Câu 16 =====
% ID: L12C3B14-113-TN
% Mức: NB
\dangbt{Giải thích hiện tượng và ứng dụng của từ trường trong thực tiễn}
\begin{ex}
% ID: L12C3B14-113-TN
% Mức: NB
Phát biểu nào sau đây đúng về giải thích hiện tượng và ứng dụng của từ trường trong thực tiễn?
\choice
{\True La bàn định hướng nhờ kim nam châm tương tác với từ trường Trái Đất.}
{La bàn hoạt động chủ yếu nhờ điện trường của Trái Đất.}
{Đại lượng đang xét không phụ thuộc các điều kiện vật lí nêu trong bài.}
{Có thể bỏ qua phương, chiều và đơn vị của các đại lượng trong mọi trường hợp.}
\loigiai{
La bàn định hướng nhờ kim nam châm tương tác với từ trường Trái Đất. Vì vậy phương án $A$ đúng; các phương án còn lại trái với định nghĩa, định luật hoặc điều kiện áp dụng.
}
\end{ex}

% ===== Câu 17 =====
% ID: L12C3B14-114-TN
% Mức: NB
\begin{ex}
% ID: L12C3B14-114-TN
% Mức: NB
Trong giải thích hiện tượng và ứng dụng của từ trường trong thực tiễn, nhận định nào phù hợp với kiến thức Vật lí?
\choice
{La bàn hoạt động chủ yếu nhờ điện trường của Trái Đất.}
{\True La bàn định hướng nhờ kim nam châm tương tác với từ trường Trái Đất.}
{Có thể bỏ qua phương, chiều và đơn vị của các đại lượng trong mọi trường hợp.}
{Đại lượng đang xét không phụ thuộc các điều kiện vật lí nêu trong bài.}
\loigiai{
La bàn định hướng nhờ kim nam châm tương tác với từ trường Trái Đất. Vì vậy phương án $B$ đúng; các phương án còn lại trái với định nghĩa, định luật hoặc điều kiện áp dụng.
}
\end{ex}

% ===== Câu 18 =====
% ID: L12C3B14-115-TN
% Mức: NB
\begin{ex}
% ID: L12C3B14-115-TN
% Mức: NB
Tình huống 3: chọn kết luận chính xác về giải thích hiện tượng và ứng dụng của từ trường trong thực tiễn.
\choice
{Đại lượng đang xét không phụ thuộc các điều kiện vật lí nêu trong bài.}
{Có thể bỏ qua phương, chiều và đơn vị của các đại lượng trong mọi trường hợp.}
{\True La bàn định hướng nhờ kim nam châm tương tác với từ trường Trái Đất.}
{La bàn hoạt động chủ yếu nhờ điện trường của Trái Đất.}
\loigiai{
La bàn định hướng nhờ kim nam châm tương tác với từ trường Trái Đất. Vì vậy phương án $C$ đúng; các phương án còn lại trái với định nghĩa, định luật hoặc điều kiện áp dụng.
}
\end{ex}

% ===== Câu 19 =====
% ID: L12C3B14-116-TN
% Mức: TH
\begin{ex}
% ID: L12C3B14-116-TN
% Mức: TH
Nội dung nào mô tả đúng nhất giải thích hiện tượng và ứng dụng của từ trường trong thực tiễn?
\choice
{Có thể bỏ qua phương, chiều và đơn vị của các đại lượng trong mọi trường hợp.}
{Đại lượng đang xét không phụ thuộc các điều kiện vật lí nêu trong bài.}
{La bàn hoạt động chủ yếu nhờ điện trường của Trái Đất.}
{\True La bàn định hướng nhờ kim nam châm tương tác với từ trường Trái Đất.}
\loigiai{
La bàn định hướng nhờ kim nam châm tương tác với từ trường Trái Đất. Vì vậy phương án $D$ đúng; các phương án còn lại trái với định nghĩa, định luật hoặc điều kiện áp dụng.
}
\end{ex}

% ===== Câu 20 =====
% ID: L12C3B14-117-TN
% Mức: TH
\begin{ex}
% ID: L12C3B14-117-TN
% Mức: TH
Khi phân tích giải thích hiện tượng và ứng dụng của từ trường trong thực tiễn, phát biểu nào đúng?
\choice
{\True La bàn định hướng nhờ kim nam châm tương tác với từ trường Trái Đất.}
{La bàn hoạt động chủ yếu nhờ điện trường của Trái Đất.}
{Đại lượng đang xét không phụ thuộc các điều kiện vật lí nêu trong bài.}
{Có thể bỏ qua phương, chiều và đơn vị của các đại lượng trong mọi trường hợp.}
\loigiai{
La bàn định hướng nhờ kim nam châm tương tác với từ trường Trái Đất. Vì vậy phương án $A$ đúng; các phương án còn lại trái với định nghĩa, định luật hoặc điều kiện áp dụng.
}
\end{ex}

% ===== Câu 21 =====
% ID: L12C3B14-118-TN
% Mức: TH
\begin{ex}
% ID: L12C3B14-118-TN
% Mức: TH
Một học sinh ôn tập giải thích hiện tượng và ứng dụng của từ trường trong thực tiễn. Kết luận nào của học sinh là đúng?
\choice
{La bàn hoạt động chủ yếu nhờ điện trường của Trái Đất.}
{\True La bàn định hướng nhờ kim nam châm tương tác với từ trường Trái Đất.}
{Có thể bỏ qua phương, chiều và đơn vị của các đại lượng trong mọi trường hợp.}
{Đại lượng đang xét không phụ thuộc các điều kiện vật lí nêu trong bài.}
\loigiai{
La bàn định hướng nhờ kim nam châm tương tác với từ trường Trái Đất. Vì vậy phương án $B$ đúng; các phương án còn lại trái với định nghĩa, định luật hoặc điều kiện áp dụng.
}
\end{ex}

% ===== Câu 22 =====
% ID: L12C3B14-119-TN
% Mức: TH
\begin{ex}
% ID: L12C3B14-119-TN
% Mức: TH
Chọn phát biểu không trái với kiến thức về giải thích hiện tượng và ứng dụng của từ trường trong thực tiễn.
\choice
{Đại lượng đang xét không phụ thuộc các điều kiện vật lí nêu trong bài.}
{Có thể bỏ qua phương, chiều và đơn vị của các đại lượng trong mọi trường hợp.}
{\True La bàn định hướng nhờ kim nam châm tương tác với từ trường Trái Đất.}
{La bàn hoạt động chủ yếu nhờ điện trường của Trái Đất.}
\loigiai{
La bàn định hướng nhờ kim nam châm tương tác với từ trường Trái Đất. Vì vậy phương án $C$ đúng; các phương án còn lại trái với định nghĩa, định luật hoặc điều kiện áp dụng.
}
\end{ex}

% ===== Câu 23 =====
% ID: L12C3B14-12-DS
% Mức: TH
\dangbt{Nhận biết từ trường, nguồn gây ra từ trường và tương tác từ}
\begin{ex}
% ID: L12C3B14-12-DS
% Mức: TH
Xét nhận biết từ trường, nguồn gây ra từ trường và tương tác từ (Tình huống 2). Hãy xác định tính đúng, sai của bốn phát biểu sau.
\choiceTF
{Mọi điện tích đứng yên đều tạo ra từ trường giống dòng điện.}
{\True Từ trường tồn tại xung quanh nam châm và dòng điện; nó gây lực từ lên nam châm hoặc dòng điện khác đặt trong đó.}
{Trong dạng bài này, kết quả không phụ thuộc dữ kiện và có thể bỏ qua điều kiện áp dụng.}
{\True Khi giải bài thuộc dạng này phải xác định đúng phương, chiều, điều kiện áp dụng và đơn vị của các đại lượng.}
\loigiai{
\begin{itemchoice}
\item (Sai) Mọi điện tích đứng yên đều tạo ra từ trường giống dòng điện. (Vì): Mọi điện tích đứng yên đều tạo ra từ trường giống dòng điện. b) Đúng: Từ trường tồn tại xung quanh nam châm và dòng điện; nó gây lực từ lên nam châm hoặc dòng điện khác đặt trong đó. c) Sai: Trong dạng bài này, kết quả không phụ thuộc dữ kiện và có thể bỏ qua điều kiện áp dụng. d) Đúng: Khi giải bài thuộc dạng này phải xác định đúng phương, chiều, điều kiện áp dụng và đơn vị của các đại lượng. Kiến thức cốt lõi: Từ trường tồn tại xung quanh nam châm và dòng điện; nó gây lực từ lên nam châm hoặc dòng điện khác đặt trong đó
\item (Đúng) Từ trường tồn tại xung quanh nam châm và dòng điện; nó gây lực từ lên nam châm hoặc dòng điện khác đặt trong đó.
\item (Sai) Trong dạng bài này, kết quả không phụ thuộc dữ kiện và có thể bỏ qua điều kiện áp dụng.
\item (Đúng) Khi giải bài thuộc dạng này phải xác định đúng phương, chiều, điều kiện áp dụng và đơn vị của các đại lượng.
\end{itemchoice}
}
\end{ex}

% ===== Câu 24 =====
% ID: L12C3B14-12-TL
% Mức: VDC
\dangbt{Mô tả đường sức từ và xác định chiều đường sức từ}
\begin{bt}
% ID: L12C3B14-12-TL
% Mức: VDC
So sánh nhận định đúng “Tiếp tuyến với đường sức từ tại một điểm cho phương của vectơ cảm ứng từ; bên ngoài nam châm, đường sức đi từ cực Bắc đến cực Nam.” với nhận định sai “Các đường sức từ có thể cắt nhau tại một điểm.”; chỉ rõ căn cứ Vật lí.
\loigiai{
Cần nêu được: Tiếp tuyến với đường sức từ tại một điểm cho phương của vectơ cảm ứng từ; bên ngoài nam châm, đường sức đi từ cực Bắc đến cực Nam. Khi vận dụng phải xác định đúng dữ kiện, phương chiều nếu có, điều kiện áp dụng, hệ đơn vị và kiểm tra tính hợp lí của kết quả. Nhận định “Các đường sức từ có thể cắt nhau tại một điểm.” sai vì trái với kiến thức nêu trên.
}
\end{bt}

% ===== Câu 25 =====
% ID: L12C3B14-12-TLN
% Mức: VDC
\begin{ex}
% ID: L12C3B14-12-TLN
% Mức: VDC
Trong bốn phát biểu sau về mô tả đường sức từ và xác định chiều đường sức từ, có bao nhiêu phát biểu đúng? (Chỉ ghi một số nguyên.) (1) Kết quả cần có ý nghĩa vật lí. (2) Các đường sức từ có thể cắt nhau tại một điểm. (3) Phải dùng đúng định luật cho tình huống. (4) Tiếp tuyến với đường sức từ tại một điểm cho phương của vectơ cảm ứng từ; bên ngoài nam châm, đường sức đi từ cực Bắc đến cực Nam.
\shortans{3}
\loigiai{
Đối chiếu từng phát biểu với kiến thức: Tiếp tuyến với đường sức từ tại một điểm cho phương của vectơ cảm ứng từ; bên ngoài nam châm, đường sức đi từ cực Bắc đến cực Nam. Các phát biểu đúng có số lượng là 3.
}
\end{ex}

% ===== Câu 26 =====
% ID: L12C3B14-12-TN
% Mức: NB
\begin{ex}
% ID: L12C3B14-12-TN
% Mức: NB
Trong mô tả đường sức từ và xác định chiều đường sức từ, nhận định nào phù hợp với kiến thức Vật lí?
\choice
{Các đường sức từ có thể cắt nhau tại một điểm.}
{\True Tiếp tuyến với đường sức từ tại một điểm cho phương của vectơ cảm ứng từ; bên ngoài nam châm, đường sức đi từ cực Bắc đến cực Nam.}
{Có thể bỏ qua phương, chiều và đơn vị của các đại lượng trong mọi trường hợp.}
{Đại lượng đang xét không phụ thuộc các điều kiện vật lí nêu trong bài.}
\loigiai{
Tiếp tuyến với đường sức từ tại một điểm cho phương của vectơ cảm ứng từ; bên ngoài nam châm, đường sức đi từ cực Bắc đến cực Nam. Vì vậy phương án $B$ đúng; các phương án còn lại trái với định nghĩa, định luật hoặc điều kiện áp dụng.
}
\end{ex}

% ===== Câu 27 =====
% ID: L12C3B14-120-TN
% Mức: VD
\dangbt{Giải thích hiện tượng và ứng dụng của từ trường trong thực tiễn}
\begin{ex}
% ID: L12C3B14-120-TN
% Mức: VD
Cơ sở Vật lí đúng dùng để giải giải thích hiện tượng và ứng dụng của từ trường trong thực tiễn là gì?
\choice
{Có thể bỏ qua phương, chiều và đơn vị của các đại lượng trong mọi trường hợp.}
{Đại lượng đang xét không phụ thuộc các điều kiện vật lí nêu trong bài.}
{La bàn hoạt động chủ yếu nhờ điện trường của Trái Đất.}
{\True La bàn định hướng nhờ kim nam châm tương tác với từ trường Trái Đất.}
\loigiai{
La bàn định hướng nhờ kim nam châm tương tác với từ trường Trái Đất. Vì vậy phương án $D$ đúng; các phương án còn lại trái với định nghĩa, định luật hoặc điều kiện áp dụng.
}
\end{ex}

% ===== Câu 28 =====
% ID: L12C3B14-121-TN
% Mức: VD
\begin{ex}
% ID: L12C3B14-121-TN
% Mức: VD
Trong một bài thực hành về giải thích hiện tượng và ứng dụng của từ trường trong thực tiễn, nhận xét nào đúng?
\choice
{\True La bàn định hướng nhờ kim nam châm tương tác với từ trường Trái Đất.}
{La bàn hoạt động chủ yếu nhờ điện trường của Trái Đất.}
{Đại lượng đang xét không phụ thuộc các điều kiện vật lí nêu trong bài.}
{Có thể bỏ qua phương, chiều và đơn vị của các đại lượng trong mọi trường hợp.}
\loigiai{
La bàn định hướng nhờ kim nam châm tương tác với từ trường Trái Đất. Vì vậy phương án $A$ đúng; các phương án còn lại trái với định nghĩa, định luật hoặc điều kiện áp dụng.
}
\end{ex}

% ===== Câu 29 =====
% ID: L12C3B14-122-TN
% Mức: VD
\begin{ex}
% ID: L12C3B14-122-TN
% Mức: VD
Kết luận đúng khi vận dụng giải thích hiện tượng và ứng dụng của từ trường trong thực tiễn là
\choice
{La bàn hoạt động chủ yếu nhờ điện trường của Trái Đất.}
{\True La bàn định hướng nhờ kim nam châm tương tác với từ trường Trái Đất.}
{Có thể bỏ qua phương, chiều và đơn vị của các đại lượng trong mọi trường hợp.}
{Đại lượng đang xét không phụ thuộc các điều kiện vật lí nêu trong bài.}
\loigiai{
La bàn định hướng nhờ kim nam châm tương tác với từ trường Trái Đất. Vì vậy phương án $B$ đúng; các phương án còn lại trái với định nghĩa, định luật hoặc điều kiện áp dụng.
}
\end{ex}

% ===== Câu 30 =====
% ID: L12C3B14-13-DS
% Mức: VD
\dangbt{Nhận biết từ trường, nguồn gây ra từ trường và tương tác từ}
\begin{ex}
% ID: L12C3B14-13-DS
% Mức: VD
Xét nhận biết từ trường, nguồn gây ra từ trường và tương tác từ (Tình huống 3). Hãy xác định tính đúng, sai của bốn phát biểu sau.
\choiceTF
{\True Khi giải bài thuộc dạng này phải xác định đúng phương, chiều, điều kiện áp dụng và đơn vị của các đại lượng.}
{Trong dạng bài này, kết quả không phụ thuộc dữ kiện và có thể bỏ qua điều kiện áp dụng.}
{\True Từ trường tồn tại xung quanh nam châm và dòng điện; nó gây lực từ lên nam châm hoặc dòng điện khác đặt trong đó.}
{Mọi điện tích đứng yên đều tạo ra từ trường giống dòng điện.}
\loigiai{
\begin{itemchoice}
\item (Đúng) Khi giải bài thuộc dạng này phải xác định đúng phương, chiều, điều kiện áp dụng và đơn vị của các đại lượng. (Vì): Khi giải bài thuộc dạng này phải xác định đúng phương, chiều, điều kiện áp dụng và đơn vị của các đại lượng. b) Sai: Trong dạng bài này, kết quả không phụ thuộc dữ kiện và có thể bỏ qua điều kiện áp dụng. c) Đúng: Từ trường tồn tại xung quanh nam châm và dòng điện; nó gây lực từ lên nam châm hoặc dòng điện khác đặt trong đó. d) Sai: Mọi điện tích đứng yên đều tạo ra từ trường giống dòng điện. Kiến thức cốt lõi: Từ trường tồn tại xung quanh nam châm và dòng điện; nó gây lực từ lên nam châm hoặc dòng điện khác đặt trong đó
\item (Sai) Trong dạng bài này, kết quả không phụ thuộc dữ kiện và có thể bỏ qua điều kiện áp dụng.
\item (Đúng) Từ trường tồn tại xung quanh nam châm và dòng điện; nó gây lực từ lên nam châm hoặc dòng điện khác đặt trong đó.
\item (Sai) Mọi điện tích đứng yên đều tạo ra từ trường giống dòng điện.
\end{itemchoice}
}
\end{ex}

% ===== Câu 31 =====
% ID: L12C3B14-13-TL
% Mức: TH
\begin{bt}
% ID: L12C3B14-13-TL
% Mức: TH
Trình bày kiến thức cốt lõi của nhận biết từ trường, nguồn gây ra từ trường và tương tác từ và nêu điều kiện áp dụng.
\loigiai{
Cần nêu được: Từ trường tồn tại xung quanh nam châm và dòng điện; nó gây lực từ lên nam châm hoặc dòng điện khác đặt trong đó. Khi vận dụng phải xác định đúng dữ kiện, phương chiều nếu có, điều kiện áp dụng, hệ đơn vị và kiểm tra tính hợp lí của kết quả. Nhận định “Mọi điện tích đứng yên đều tạo ra từ trường giống dòng điện.” sai vì trái với kiến thức nêu trên.
}
\end{bt}

% ===== Câu 32 =====
% ID: L12C3B14-13-TLN
% Mức: TH
\begin{ex}
% ID: L12C3B14-13-TLN
% Mức: TH
Trong bốn phát biểu sau về nhận biết từ trường, nguồn gây ra từ trường và tương tác từ, có bao nhiêu phát biểu đúng? (Chỉ ghi một số nguyên.) (1) Từ trường tồn tại xung quanh nam châm và dòng điện; nó gây lực từ lên nam châm hoặc dòng điện khác đặt trong đó. (2) Mọi điện tích đứng yên đều tạo ra từ trường giống dòng điện. (3) Cần kiểm tra điều kiện áp dụng trước khi tính toán. (4) Có thể bỏ qua đơn vị của kết quả.
\shortans{2}
\loigiai{
Đối chiếu từng phát biểu với kiến thức: Từ trường tồn tại xung quanh nam châm và dòng điện; nó gây lực từ lên nam châm hoặc dòng điện khác đặt trong đó. Các phát biểu đúng có số lượng là 2.
}
\end{ex}

% ===== Câu 33 =====
% ID: L12C3B14-13-TN
% Mức: NB
\dangbt{Mô tả đường sức từ và xác định chiều đường sức từ}
\begin{ex}
% ID: L12C3B14-13-TN
% Mức: NB
Tình huống 3: chọn kết luận chính xác về mô tả đường sức từ và xác định chiều đường sức từ.
\choice
{Đại lượng đang xét không phụ thuộc các điều kiện vật lí nêu trong bài.}
{Có thể bỏ qua phương, chiều và đơn vị của các đại lượng trong mọi trường hợp.}
{\True Tiếp tuyến với đường sức từ tại một điểm cho phương của vectơ cảm ứng từ; bên ngoài nam châm, đường sức đi từ cực Bắc đến cực Nam.}
{Các đường sức từ có thể cắt nhau tại một điểm.}
\loigiai{
Tiếp tuyến với đường sức từ tại một điểm cho phương của vectơ cảm ứng từ; bên ngoài nam châm, đường sức đi từ cực Bắc đến cực Nam. Vì vậy phương án $C$ đúng; các phương án còn lại trái với định nghĩa, định luật hoặc điều kiện áp dụng.
}
\end{ex}

% ===== Câu 34 =====
% ID: L12C3B14-14-DS
% Mức: VD
\dangbt{Nhận biết từ trường, nguồn gây ra từ trường và tương tác từ}
\begin{ex}
% ID: L12C3B14-14-DS
% Mức: VD
Xét nhận biết từ trường, nguồn gây ra từ trường và tương tác từ (Tình huống 4). Hãy xác định tính đúng, sai của bốn phát biểu sau.
\choiceTF
{Trong dạng bài này, kết quả không phụ thuộc dữ kiện và có thể bỏ qua điều kiện áp dụng.}
{\True Khi giải bài thuộc dạng này phải xác định đúng phương, chiều, điều kiện áp dụng và đơn vị của các đại lượng.}
{Mọi điện tích đứng yên đều tạo ra từ trường giống dòng điện.}
{\True Từ trường tồn tại xung quanh nam châm và dòng điện; nó gây lực từ lên nam châm hoặc dòng điện khác đặt trong đó.}
\loigiai{
\begin{itemchoice}
\item (Sai) Trong dạng bài này, kết quả không phụ thuộc dữ kiện và có thể bỏ qua điều kiện áp dụng. (Vì): Trong dạng bài này, kết quả không phụ thuộc dữ kiện và có thể bỏ qua điều kiện áp dụng. b) Đúng: Khi giải bài thuộc dạng này phải xác định đúng phương, chiều, điều kiện áp dụng và đơn vị của các đại lượng. c) Sai: Mọi điện tích đứng yên đều tạo ra từ trường giống dòng điện. d) Đúng: Từ trường tồn tại xung quanh nam châm và dòng điện; nó gây lực từ lên nam châm hoặc dòng điện khác đặt trong đó. Kiến thức cốt lõi: Từ trường tồn tại xung quanh nam châm và dòng điện; nó gây lực từ lên nam châm hoặc dòng điện khác đặt trong đó
\item (Đúng) Khi giải bài thuộc dạng này phải xác định đúng phương, chiều, điều kiện áp dụng và đơn vị của các đại lượng.
\item (Sai) Mọi điện tích đứng yên đều tạo ra từ trường giống dòng điện.
\item (Đúng) Từ trường tồn tại xung quanh nam châm và dòng điện; nó gây lực từ lên nam châm hoặc dòng điện khác đặt trong đó.
\end{itemchoice}
}
\end{ex}

% ===== Câu 35 =====
% ID: L12C3B14-14-TL
% Mức: TH
\begin{bt}
% ID: L12C3B14-14-TL
% Mức: TH
Giải thích vì sao nhận định sau đúng: “Từ trường tồn tại xung quanh nam châm và dòng điện; nó gây lực từ lên nam châm hoặc dòng điện khác đặt trong đó.”
\loigiai{
Cần nêu được: Từ trường tồn tại xung quanh nam châm và dòng điện; nó gây lực từ lên nam châm hoặc dòng điện khác đặt trong đó. Khi vận dụng phải xác định đúng dữ kiện, phương chiều nếu có, điều kiện áp dụng, hệ đơn vị và kiểm tra tính hợp lí của kết quả. Nhận định “Mọi điện tích đứng yên đều tạo ra từ trường giống dòng điện.” sai vì trái với kiến thức nêu trên.
}
\end{bt}

% ===== Câu 36 =====
% ID: L12C3B14-14-TLN
% Mức: TH
\begin{ex}
% ID: L12C3B14-14-TLN
% Mức: TH
Trong bốn phát biểu sau về nhận biết từ trường, nguồn gây ra từ trường và tương tác từ, có bao nhiêu phát biểu đúng? (Chỉ ghi một số nguyên.) (1) Mọi điện tích đứng yên đều tạo ra từ trường giống dòng điện. (2) Từ trường tồn tại xung quanh nam châm và dòng điện; nó gây lực từ lên nam châm hoặc dòng điện khác đặt trong đó. (3) Kết quả phải phù hợp ý nghĩa vật lí. (4) Mọi đại lượng đều là vô hướng.
\shortans{1}
\loigiai{
Đối chiếu từng phát biểu với kiến thức: Từ trường tồn tại xung quanh nam châm và dòng điện; nó gây lực từ lên nam châm hoặc dòng điện khác đặt trong đó. Các phát biểu đúng có số lượng là 1.
}
\end{ex}

% ===== Câu 37 =====
% ID: L12C3B14-14-TN
% Mức: TH
\dangbt{Mô tả đường sức từ và xác định chiều đường sức từ}
\begin{ex}
% ID: L12C3B14-14-TN
% Mức: TH
Nội dung nào mô tả đúng nhất mô tả đường sức từ và xác định chiều đường sức từ?
\choice
{Có thể bỏ qua phương, chiều và đơn vị của các đại lượng trong mọi trường hợp.}
{Đại lượng đang xét không phụ thuộc các điều kiện vật lí nêu trong bài.}
{Các đường sức từ có thể cắt nhau tại một điểm.}
{\True Tiếp tuyến với đường sức từ tại một điểm cho phương của vectơ cảm ứng từ; bên ngoài nam châm, đường sức đi từ cực Bắc đến cực Nam.}
\loigiai{
Tiếp tuyến với đường sức từ tại một điểm cho phương của vectơ cảm ứng từ; bên ngoài nam châm, đường sức đi từ cực Bắc đến cực Nam. Vì vậy phương án $D$ đúng; các phương án còn lại trái với định nghĩa, định luật hoặc điều kiện áp dụng.
}
\end{ex}

% ===== Câu 38 =====
% ID: L12C3B14-15-DS
% Mức: VD
\dangbt{Nhận biết từ trường, nguồn gây ra từ trường và tương tác từ}
\begin{ex}
% ID: L12C3B14-15-DS
% Mức: VD
Xét nhận biết từ trường, nguồn gây ra từ trường và tương tác từ (Tình huống 5). Hãy xác định tính đúng, sai của bốn phát biểu sau.
\choiceTF
{\True Từ trường tồn tại xung quanh nam châm và dòng điện; nó gây lực từ lên nam châm hoặc dòng điện khác đặt trong đó.}
{Trong dạng bài này, kết quả không phụ thuộc dữ kiện và có thể bỏ qua điều kiện áp dụng.}
{Mọi điện tích đứng yên đều tạo ra từ trường giống dòng điện.}
{\True Khi giải bài thuộc dạng này phải xác định đúng phương, chiều, điều kiện áp dụng và đơn vị của các đại lượng.}
\loigiai{
\begin{itemchoice}
\item (Đúng) Từ trường tồn tại xung quanh nam châm và dòng điện; nó gây lực từ lên nam châm hoặc dòng điện khác đặt trong đó. (Vì): Từ trường tồn tại xung quanh nam châm và dòng điện; nó gây lực từ lên nam châm hoặc dòng điện khác đặt trong đó. b) Sai: Trong dạng bài này, kết quả không phụ thuộc dữ kiện và có thể bỏ qua điều kiện áp dụng. c) Sai: Mọi điện tích đứng yên đều tạo ra từ trường giống dòng điện. d) Đúng: Khi giải bài thuộc dạng này phải xác định đúng phương, chiều, điều kiện áp dụng và đơn vị của các đại lượng. Kiến thức cốt lõi: Từ trường tồn tại xung quanh nam châm và dòng điện; nó gây lực từ lên nam châm hoặc dòng điện khác đặt trong đó
\item (Sai) Trong dạng bài này, kết quả không phụ thuộc dữ kiện và có thể bỏ qua điều kiện áp dụng.
\item (Sai) Mọi điện tích đứng yên đều tạo ra từ trường giống dòng điện.
\item (Đúng) Khi giải bài thuộc dạng này phải xác định đúng phương, chiều, điều kiện áp dụng và đơn vị của các đại lượng.
\end{itemchoice}
}
\end{ex}

% ===== Câu 39 =====
% ID: L12C3B14-15-TL
% Mức: VD
\begin{bt}
% ID: L12C3B14-15-TL
% Mức: VD
Phân tích sai lầm trong nhận định: “Mọi điện tích đứng yên đều tạo ra từ trường giống dòng điện.”
\loigiai{
Cần nêu được: Từ trường tồn tại xung quanh nam châm và dòng điện; nó gây lực từ lên nam châm hoặc dòng điện khác đặt trong đó. Khi vận dụng phải xác định đúng dữ kiện, phương chiều nếu có, điều kiện áp dụng, hệ đơn vị và kiểm tra tính hợp lí của kết quả. Nhận định “Mọi điện tích đứng yên đều tạo ra từ trường giống dòng điện.” sai vì trái với kiến thức nêu trên.
}
\end{bt}

% ===== Câu 40 =====
% ID: L12C3B14-15-TLN
% Mức: TH
\begin{ex}
% ID: L12C3B14-15-TLN
% Mức: TH
Trong bốn phát biểu sau về nhận biết từ trường, nguồn gây ra từ trường và tương tác từ, có bao nhiêu phát biểu đúng? (Chỉ ghi một số nguyên.) (1) Cần đổi các đại lượng về hệ đơn vị thống nhất. (2) Từ trường tồn tại xung quanh nam châm và dòng điện; nó gây lực từ lên nam châm hoặc dòng điện khác đặt trong đó. (3) Mọi điện tích đứng yên đều tạo ra từ trường giống dòng điện. (4) Không cần xét chiều của đại lượng vectơ.
\shortans{2}
\loigiai{
Đối chiếu từng phát biểu với kiến thức: Từ trường tồn tại xung quanh nam châm và dòng điện; nó gây lực từ lên nam châm hoặc dòng điện khác đặt trong đó. Các phát biểu đúng có số lượng là 2.
}
\end{ex}

% ===== Câu 41 =====
% ID: L12C3B14-15-TN
% Mức: TH
\dangbt{Mô tả đường sức từ và xác định chiều đường sức từ}
\begin{ex}
% ID: L12C3B14-15-TN
% Mức: TH
Khi phân tích mô tả đường sức từ và xác định chiều đường sức từ, phát biểu nào đúng?
\choice
{\True Tiếp tuyến với đường sức từ tại một điểm cho phương của vectơ cảm ứng từ; bên ngoài nam châm, đường sức đi từ cực Bắc đến cực Nam.}
{Các đường sức từ có thể cắt nhau tại một điểm.}
{Đại lượng đang xét không phụ thuộc các điều kiện vật lí nêu trong bài.}
{Có thể bỏ qua phương, chiều và đơn vị của các đại lượng trong mọi trường hợp.}
\loigiai{
Tiếp tuyến với đường sức từ tại một điểm cho phương của vectơ cảm ứng từ; bên ngoài nam châm, đường sức đi từ cực Bắc đến cực Nam. Vì vậy phương án $A$ đúng; các phương án còn lại trái với định nghĩa, định luật hoặc điều kiện áp dụng.
}
\end{ex}

% ===== Câu 42 =====
% ID: L12C3B14-16-DS
% Mức: TH
\dangbt{So sánh, tổng hợp và biểu diễn các từ trường đơn giản}
\begin{ex}
% ID: L12C3B14-16-DS
% Mức: TH
Xét so sánh, tổng hợp và biểu diễn các từ trường đơn giản (Tình huống 1). Hãy xác định tính đúng, sai của bốn phát biểu sau.
\choiceTF
{\True Vectơ cảm ứng từ tổng hợp tại một điểm bằng tổng vectơ các cảm ứng từ thành phần.}
{Độ lớn cảm ứng từ tổng hợp luôn bằng tổng độ lớn các cảm ứng từ thành phần.}
{\True Khi giải bài thuộc dạng này phải xác định đúng phương, chiều, điều kiện áp dụng và đơn vị của các đại lượng.}
{Trong dạng bài này, kết quả không phụ thuộc dữ kiện và có thể bỏ qua điều kiện áp dụng.}
\loigiai{
\begin{itemchoice}
\item (Đúng) Vectơ cảm ứng từ tổng hợp tại một điểm bằng tổng vectơ các cảm ứng từ thành phần. (Vì): Vectơ cảm ứng từ tổng hợp tại một điểm bằng tổng vectơ các cảm ứng từ thành phần. b) Sai: Độ lớn cảm ứng từ tổng hợp luôn bằng tổng độ lớn các cảm ứng từ thành phần. c) Đúng: Khi giải bài thuộc dạng này phải xác định đúng phương, chiều, điều kiện áp dụng và đơn vị của các đại lượng. d) Sai: Trong dạng bài này, kết quả không phụ thuộc dữ kiện và có thể bỏ qua điều kiện áp dụng. Kiến thức cốt lõi: Vectơ cảm ứng từ tổng hợp tại một điểm bằng tổng vectơ các cảm ứng từ thành phần
\item (Sai) Độ lớn cảm ứng từ tổng hợp luôn bằng tổng độ lớn các cảm ứng từ thành phần.
\item (Đúng) Khi giải bài thuộc dạng này phải xác định đúng phương, chiều, điều kiện áp dụng và đơn vị của các đại lượng.
\item (Sai) Trong dạng bài này, kết quả không phụ thuộc dữ kiện và có thể bỏ qua điều kiện áp dụng.
\end{itemchoice}
}
\end{ex}

% ===== Câu 43 =====
% ID: L12C3B14-16-TL
% Mức: VD
\dangbt{Nhận biết từ trường, nguồn gây ra từ trường và tương tác từ}
\begin{bt}
% ID: L12C3B14-16-TL
% Mức: VD
Đề xuất các bước giải một bài toán thuộc dạng nhận biết từ trường, nguồn gây ra từ trường và tương tác từ.
\loigiai{
Cần nêu được: Từ trường tồn tại xung quanh nam châm và dòng điện; nó gây lực từ lên nam châm hoặc dòng điện khác đặt trong đó. Khi vận dụng phải xác định đúng dữ kiện, phương chiều nếu có, điều kiện áp dụng, hệ đơn vị và kiểm tra tính hợp lí của kết quả. Nhận định “Mọi điện tích đứng yên đều tạo ra từ trường giống dòng điện.” sai vì trái với kiến thức nêu trên.
}
\end{bt}

% ===== Câu 44 =====
% ID: L12C3B14-16-TLN
% Mức: VD
\begin{ex}
% ID: L12C3B14-16-TLN
% Mức: VD
Trong bốn phát biểu sau về nhận biết từ trường, nguồn gây ra từ trường và tương tác từ, có bao nhiêu phát biểu đúng? (Chỉ ghi một số nguyên.) (1) Từ trường tồn tại xung quanh nam châm và dòng điện; nó gây lực từ lên nam châm hoặc dòng điện khác đặt trong đó. (2) Cần đối chiếu kết quả với điều kiện bài toán. (3) Mọi điện tích đứng yên đều tạo ra từ trường giống dòng điện. (4) Mọi trường hợp đều cho cùng một kết quả.
\shortans{2}
\loigiai{
Đối chiếu từng phát biểu với kiến thức: Từ trường tồn tại xung quanh nam châm và dòng điện; nó gây lực từ lên nam châm hoặc dòng điện khác đặt trong đó. Các phát biểu đúng có số lượng là 2.
}
\end{ex}

% ===== Câu 45 =====
% ID: L12C3B14-16-TN
% Mức: TH
\dangbt{Mô tả đường sức từ và xác định chiều đường sức từ}
\begin{ex}
% ID: L12C3B14-16-TN
% Mức: TH
Một học sinh ôn tập mô tả đường sức từ và xác định chiều đường sức từ. Kết luận nào của học sinh là đúng?
\choice
{Các đường sức từ có thể cắt nhau tại một điểm.}
{\True Tiếp tuyến với đường sức từ tại một điểm cho phương của vectơ cảm ứng từ; bên ngoài nam châm, đường sức đi từ cực Bắc đến cực Nam.}
{Có thể bỏ qua phương, chiều và đơn vị của các đại lượng trong mọi trường hợp.}
{Đại lượng đang xét không phụ thuộc các điều kiện vật lí nêu trong bài.}
\loigiai{
Tiếp tuyến với đường sức từ tại một điểm cho phương của vectơ cảm ứng từ; bên ngoài nam châm, đường sức đi từ cực Bắc đến cực Nam. Vì vậy phương án $B$ đúng; các phương án còn lại trái với định nghĩa, định luật hoặc điều kiện áp dụng.
}
\end{ex}

% ===== Câu 46 =====
% ID: L12C3B14-17-DS
% Mức: TH
\dangbt{So sánh, tổng hợp và biểu diễn các từ trường đơn giản}
\begin{ex}
% ID: L12C3B14-17-DS
% Mức: TH
Xét so sánh, tổng hợp và biểu diễn các từ trường đơn giản (Tình huống 2). Hãy xác định tính đúng, sai của bốn phát biểu sau.
\choiceTF
{Độ lớn cảm ứng từ tổng hợp luôn bằng tổng độ lớn các cảm ứng từ thành phần.}
{\True Vectơ cảm ứng từ tổng hợp tại một điểm bằng tổng vectơ các cảm ứng từ thành phần.}
{Trong dạng bài này, kết quả không phụ thuộc dữ kiện và có thể bỏ qua điều kiện áp dụng.}
{\True Khi giải bài thuộc dạng này phải xác định đúng phương, chiều, điều kiện áp dụng và đơn vị của các đại lượng.}
\loigiai{
\begin{itemchoice}
\item (Sai) Độ lớn cảm ứng từ tổng hợp luôn bằng tổng độ lớn các cảm ứng từ thành phần. (Vì): Độ lớn cảm ứng từ tổng hợp luôn bằng tổng độ lớn các cảm ứng từ thành phần. b) Đúng: Vectơ cảm ứng từ tổng hợp tại một điểm bằng tổng vectơ các cảm ứng từ thành phần. c) Sai: Trong dạng bài này, kết quả không phụ thuộc dữ kiện và có thể bỏ qua điều kiện áp dụng. d) Đúng: Khi giải bài thuộc dạng này phải xác định đúng phương, chiều, điều kiện áp dụng và đơn vị của các đại lượng. Kiến thức cốt lõi: Vectơ cảm ứng từ tổng hợp tại một điểm bằng tổng vectơ các cảm ứng từ thành phần
\item (Đúng) Vectơ cảm ứng từ tổng hợp tại một điểm bằng tổng vectơ các cảm ứng từ thành phần.
\item (Sai) Trong dạng bài này, kết quả không phụ thuộc dữ kiện và có thể bỏ qua điều kiện áp dụng.
\item (Đúng) Khi giải bài thuộc dạng này phải xác định đúng phương, chiều, điều kiện áp dụng và đơn vị của các đại lượng.
\end{itemchoice}
}
\end{ex}

% ===== Câu 47 =====
% ID: L12C3B14-17-TL
% Mức: VD
\dangbt{Nhận biết từ trường, nguồn gây ra từ trường và tương tác từ}
\begin{bt}
% ID: L12C3B14-17-TL
% Mức: VD
Nêu một tình huống thực tiễn liên quan đến nhận biết từ trường, nguồn gây ra từ trường và tương tác từ và giải thích bằng kiến thức Vật lí.
\loigiai{
Cần nêu được: Từ trường tồn tại xung quanh nam châm và dòng điện; nó gây lực từ lên nam châm hoặc dòng điện khác đặt trong đó. Khi vận dụng phải xác định đúng dữ kiện, phương chiều nếu có, điều kiện áp dụng, hệ đơn vị và kiểm tra tính hợp lí của kết quả. Nhận định “Mọi điện tích đứng yên đều tạo ra từ trường giống dòng điện.” sai vì trái với kiến thức nêu trên.
}
\end{bt}

% ===== Câu 48 =====
% ID: L12C3B14-17-TLN
% Mức: VD
\begin{ex}
% ID: L12C3B14-17-TLN
% Mức: VD
Trong bốn phát biểu sau về nhận biết từ trường, nguồn gây ra từ trường và tương tác từ, có bao nhiêu phát biểu đúng? (Chỉ ghi một số nguyên.) (1) Mọi điện tích đứng yên đều tạo ra từ trường giống dòng điện. (2) Cần đọc đúng dữ kiện và giả thiết. (3) Từ trường tồn tại xung quanh nam châm và dòng điện; nó gây lực từ lên nam châm hoặc dòng điện khác đặt trong đó. (4) Có thể thay số trước khi xác định công thức.
\shortans{2}
\loigiai{
Đối chiếu từng phát biểu với kiến thức: Từ trường tồn tại xung quanh nam châm và dòng điện; nó gây lực từ lên nam châm hoặc dòng điện khác đặt trong đó. Các phát biểu đúng có số lượng là 2.
}
\end{ex}

% ===== Câu 49 =====
% ID: L12C3B14-17-TN
% Mức: TH
\dangbt{Mô tả đường sức từ và xác định chiều đường sức từ}
\begin{ex}
% ID: L12C3B14-17-TN
% Mức: TH
Chọn phát biểu không trái với kiến thức về mô tả đường sức từ và xác định chiều đường sức từ.
\choice
{Đại lượng đang xét không phụ thuộc các điều kiện vật lí nêu trong bài.}
{Có thể bỏ qua phương, chiều và đơn vị của các đại lượng trong mọi trường hợp.}
{\True Tiếp tuyến với đường sức từ tại một điểm cho phương của vectơ cảm ứng từ; bên ngoài nam châm, đường sức đi từ cực Bắc đến cực Nam.}
{Các đường sức từ có thể cắt nhau tại một điểm.}
\loigiai{
Tiếp tuyến với đường sức từ tại một điểm cho phương của vectơ cảm ứng từ; bên ngoài nam châm, đường sức đi từ cực Bắc đến cực Nam. Vì vậy phương án $C$ đúng; các phương án còn lại trái với định nghĩa, định luật hoặc điều kiện áp dụng.
}
\end{ex}

% ===== Câu 50 =====
% ID: L12C3B14-18-DS
% Mức: VD
\dangbt{So sánh, tổng hợp và biểu diễn các từ trường đơn giản}
\begin{ex}
% ID: L12C3B14-18-DS
% Mức: VD
Xét so sánh, tổng hợp và biểu diễn các từ trường đơn giản (Tình huống 3). Hãy xác định tính đúng, sai của bốn phát biểu sau.
\choiceTF
{\True Khi giải bài thuộc dạng này phải xác định đúng phương, chiều, điều kiện áp dụng và đơn vị của các đại lượng.}
{Trong dạng bài này, kết quả không phụ thuộc dữ kiện và có thể bỏ qua điều kiện áp dụng.}
{\True Vectơ cảm ứng từ tổng hợp tại một điểm bằng tổng vectơ các cảm ứng từ thành phần.}
{Độ lớn cảm ứng từ tổng hợp luôn bằng tổng độ lớn các cảm ứng từ thành phần.}
\loigiai{
\begin{itemchoice}
\item (Đúng) Khi giải bài thuộc dạng này phải xác định đúng phương, chiều, điều kiện áp dụng và đơn vị của các đại lượng. (Vì): Khi giải bài thuộc dạng này phải xác định đúng phương, chiều, điều kiện áp dụng và đơn vị của các đại lượng. b) Sai: Trong dạng bài này, kết quả không phụ thuộc dữ kiện và có thể bỏ qua điều kiện áp dụng. c) Đúng: Vectơ cảm ứng từ tổng hợp tại một điểm bằng tổng vectơ các cảm ứng từ thành phần. d) Sai: Độ lớn cảm ứng từ tổng hợp luôn bằng tổng độ lớn các cảm ứng từ thành phần. Kiến thức cốt lõi: Vectơ cảm ứng từ tổng hợp tại một điểm bằng tổng vectơ các cảm ứng từ thành phần
\item (Sai) Trong dạng bài này, kết quả không phụ thuộc dữ kiện và có thể bỏ qua điều kiện áp dụng.
\item (Đúng) Vectơ cảm ứng từ tổng hợp tại một điểm bằng tổng vectơ các cảm ứng từ thành phần.
\item (Sai) Độ lớn cảm ứng từ tổng hợp luôn bằng tổng độ lớn các cảm ứng từ thành phần.
\end{itemchoice}
}
\end{ex}

% ===== Câu 51 =====
% ID: L12C3B14-18-TL
% Mức: VDC
\dangbt{Nhận biết từ trường, nguồn gây ra từ trường và tương tác từ}
\begin{bt}
% ID: L12C3B14-18-TL
% Mức: VDC
So sánh nhận định đúng “Từ trường tồn tại xung quanh nam châm và dòng điện; nó gây lực từ lên nam châm hoặc dòng điện khác đặt trong đó.” với nhận định sai “Mọi điện tích đứng yên đều tạo ra từ trường giống dòng điện.”; chỉ rõ căn cứ Vật lí.
\loigiai{
Cần nêu được: Từ trường tồn tại xung quanh nam châm và dòng điện; nó gây lực từ lên nam châm hoặc dòng điện khác đặt trong đó. Khi vận dụng phải xác định đúng dữ kiện, phương chiều nếu có, điều kiện áp dụng, hệ đơn vị và kiểm tra tính hợp lí của kết quả. Nhận định “Mọi điện tích đứng yên đều tạo ra từ trường giống dòng điện.” sai vì trái với kiến thức nêu trên.
}
\end{bt}

% ===== Câu 52 =====
% ID: L12C3B14-18-TLN
% Mức: VDC
\begin{ex}
% ID: L12C3B14-18-TLN
% Mức: VDC
Trong bốn phát biểu sau về nhận biết từ trường, nguồn gây ra từ trường và tương tác từ, có bao nhiêu phát biểu đúng? (Chỉ ghi một số nguyên.) (1) Kết quả cần có ý nghĩa vật lí. (2) Mọi điện tích đứng yên đều tạo ra từ trường giống dòng điện. (3) Phải dùng đúng định luật cho tình huống. (4) Từ trường tồn tại xung quanh nam châm và dòng điện; nó gây lực từ lên nam châm hoặc dòng điện khác đặt trong đó.
\shortans{3}
\loigiai{
Đối chiếu từng phát biểu với kiến thức: Từ trường tồn tại xung quanh nam châm và dòng điện; nó gây lực từ lên nam châm hoặc dòng điện khác đặt trong đó. Các phát biểu đúng có số lượng là 3.
}
\end{ex}

% ===== Câu 53 =====
% ID: L12C3B14-18-TN
% Mức: TH
\dangbt{Mô tả đường sức từ và xác định chiều đường sức từ}
\begin{ex}
% ID: L12C3B14-18-TN
% Mức: TH
Tính chất nào sau đây mô tả đúng về các đường sức từ?
\choice
{Nơi nào từ trường mạnh hơn thì các đường sức từ ở đó vẽ thưa hơn}
{Các đường sức từ là những đường cong không khép kín}
{\True Chiều của đường sức từ là chiều của vectơ cảm ứng từ}
{Tại mỗi điểm trong từ trường có thể vẽ nhiều đường sức từ giao nhau}
\loigiai{
Các đặc điểm của đường sức từ:
 
  
• Tại mỗi điểm trong từ trường, chỉ có thể vẽ được một đường sức từ đi qua và chỉ một mà thôi.
  
• Các đường sức từ là những đường cong khép kín.
  
• Nơi nào từ trường mạnh hơn thì các đường sức từ ở đó vẽ dày hơn; nơi nào từ trường yếu hơn thì các đường sức từ vẽ thưa hơn.
  
• Tiếp tuyến tại mỗi điểm trên đường sức từ trùng với phương của từ trường tại điểm đó.
}
\end{ex}

% ===== Câu 54 =====
% ID: L12C3B14-19-DS
% Mức: VD
\dangbt{So sánh, tổng hợp và biểu diễn các từ trường đơn giản}
\begin{ex}
% ID: L12C3B14-19-DS
% Mức: VD
Xét so sánh, tổng hợp và biểu diễn các từ trường đơn giản (Tình huống 4). Hãy xác định tính đúng, sai của bốn phát biểu sau.
\choiceTF
{Trong dạng bài này, kết quả không phụ thuộc dữ kiện và có thể bỏ qua điều kiện áp dụng.}
{\True Khi giải bài thuộc dạng này phải xác định đúng phương, chiều, điều kiện áp dụng và đơn vị của các đại lượng.}
{Độ lớn cảm ứng từ tổng hợp luôn bằng tổng độ lớn các cảm ứng từ thành phần.}
{\True Vectơ cảm ứng từ tổng hợp tại một điểm bằng tổng vectơ các cảm ứng từ thành phần.}
\loigiai{
\begin{itemchoice}
\item (Sai) Trong dạng bài này, kết quả không phụ thuộc dữ kiện và có thể bỏ qua điều kiện áp dụng. (Vì): Trong dạng bài này, kết quả không phụ thuộc dữ kiện và có thể bỏ qua điều kiện áp dụng. b) Đúng: Khi giải bài thuộc dạng này phải xác định đúng phương, chiều, điều kiện áp dụng và đơn vị của các đại lượng. c) Sai: Độ lớn cảm ứng từ tổng hợp luôn bằng tổng độ lớn các cảm ứng từ thành phần. d) Đúng: Vectơ cảm ứng từ tổng hợp tại một điểm bằng tổng vectơ các cảm ứng từ thành phần. Kiến thức cốt lõi: Vectơ cảm ứng từ tổng hợp tại một điểm bằng tổng vectơ các cảm ứng từ thành phần
\item (Đúng) Khi giải bài thuộc dạng này phải xác định đúng phương, chiều, điều kiện áp dụng và đơn vị của các đại lượng.
\item (Sai) Độ lớn cảm ứng từ tổng hợp luôn bằng tổng độ lớn các cảm ứng từ thành phần.
\item (Đúng) Vectơ cảm ứng từ tổng hợp tại một điểm bằng tổng vectơ các cảm ứng từ thành phần.
\end{itemchoice}
}
\end{ex}

% ===== Câu 55 =====
% ID: L12C3B14-19-TL
% Mức: TH
\begin{bt}
% ID: L12C3B14-19-TL
% Mức: TH
Trình bày kiến thức cốt lõi của so sánh, tổng hợp và biểu diễn các từ trường đơn giản và nêu điều kiện áp dụng.
\loigiai{
Cần nêu được: Vectơ cảm ứng từ tổng hợp tại một điểm bằng tổng vectơ các cảm ứng từ thành phần. Khi vận dụng phải xác định đúng dữ kiện, phương chiều nếu có, điều kiện áp dụng, hệ đơn vị và kiểm tra tính hợp lí của kết quả. Nhận định “Độ lớn cảm ứng từ tổng hợp luôn bằng tổng độ lớn các cảm ứng từ thành phần.” sai vì trái với kiến thức nêu trên.
}
\end{bt}

% ===== Câu 56 =====
% ID: L12C3B14-19-TLN
% Mức: TH
\begin{ex}
% ID: L12C3B14-19-TLN
% Mức: TH
Trong bốn phát biểu sau về so sánh, tổng hợp và biểu diễn các từ trường đơn giản, có bao nhiêu phát biểu đúng? (Chỉ ghi một số nguyên.) (1) Vectơ cảm ứng từ tổng hợp tại một điểm bằng tổng vectơ các cảm ứng từ thành phần. (2) Độ lớn cảm ứng từ tổng hợp luôn bằng tổng độ lớn các cảm ứng từ thành phần. (3) Cần kiểm tra điều kiện áp dụng trước khi tính toán. (4) Có thể bỏ qua đơn vị của kết quả.
\shortans{2}
\loigiai{
Đối chiếu từng phát biểu với kiến thức: Vectơ cảm ứng từ tổng hợp tại một điểm bằng tổng vectơ các cảm ứng từ thành phần. Các phát biểu đúng có số lượng là 2.
}
\end{ex}

% ===== Câu 57 =====
% ID: L12C3B14-19-TN
% Mức: VD
\dangbt{Mô tả đường sức từ và xác định chiều đường sức từ}
\begin{ex}
% ID: L12C3B14-19-TN
% Mức: VD
Cơ sở Vật lí đúng dùng để giải mô tả đường sức từ và xác định chiều đường sức từ là gì?
\choice
{Có thể bỏ qua phương, chiều và đơn vị của các đại lượng trong mọi trường hợp.}
{Đại lượng đang xét không phụ thuộc các điều kiện vật lí nêu trong bài.}
{Các đường sức từ có thể cắt nhau tại một điểm.}
{\True Tiếp tuyến với đường sức từ tại một điểm cho phương của vectơ cảm ứng từ; bên ngoài nam châm, đường sức đi từ cực Bắc đến cực Nam.}
\loigiai{
Tiếp tuyến với đường sức từ tại một điểm cho phương của vectơ cảm ứng từ; bên ngoài nam châm, đường sức đi từ cực Bắc đến cực Nam. Vì vậy phương án $D$ đúng; các phương án còn lại trái với định nghĩa, định luật hoặc điều kiện áp dụng.
}
\end{ex}

% ===== Câu 58 =====
% ID: L12C3B14-20-DS
% Mức: VD
\dangbt{So sánh, tổng hợp và biểu diễn các từ trường đơn giản}
\begin{ex}
% ID: L12C3B14-20-DS
% Mức: VD
Xét so sánh, tổng hợp và biểu diễn các từ trường đơn giản (Tình huống 5). Hãy xác định tính đúng, sai của bốn phát biểu sau.
\choiceTF
{\True Vectơ cảm ứng từ tổng hợp tại một điểm bằng tổng vectơ các cảm ứng từ thành phần.}
{Trong dạng bài này, kết quả không phụ thuộc dữ kiện và có thể bỏ qua điều kiện áp dụng.}
{Độ lớn cảm ứng từ tổng hợp luôn bằng tổng độ lớn các cảm ứng từ thành phần.}
{\True Khi giải bài thuộc dạng này phải xác định đúng phương, chiều, điều kiện áp dụng và đơn vị của các đại lượng.}
\loigiai{
\begin{itemchoice}
\item (Đúng) Vectơ cảm ứng từ tổng hợp tại một điểm bằng tổng vectơ các cảm ứng từ thành phần. (Vì): Vectơ cảm ứng từ tổng hợp tại một điểm bằng tổng vectơ các cảm ứng từ thành phần. b) Sai: Trong dạng bài này, kết quả không phụ thuộc dữ kiện và có thể bỏ qua điều kiện áp dụng. c) Sai: Độ lớn cảm ứng từ tổng hợp luôn bằng tổng độ lớn các cảm ứng từ thành phần. d) Đúng: Khi giải bài thuộc dạng này phải xác định đúng phương, chiều, điều kiện áp dụng và đơn vị của các đại lượng. Kiến thức cốt lõi: Vectơ cảm ứng từ tổng hợp tại một điểm bằng tổng vectơ các cảm ứng từ thành phần
\item (Sai) Trong dạng bài này, kết quả không phụ thuộc dữ kiện và có thể bỏ qua điều kiện áp dụng.
\item (Sai) Độ lớn cảm ứng từ tổng hợp luôn bằng tổng độ lớn các cảm ứng từ thành phần.
\item (Đúng) Khi giải bài thuộc dạng này phải xác định đúng phương, chiều, điều kiện áp dụng và đơn vị của các đại lượng.
\end{itemchoice}
}
\end{ex}

% ===== Câu 59 =====
% ID: L12C3B14-20-TL
% Mức: TH
\begin{bt}
% ID: L12C3B14-20-TL
% Mức: TH
Giải thích vì sao nhận định sau đúng: “Vectơ cảm ứng từ tổng hợp tại một điểm bằng tổng vectơ các cảm ứng từ thành phần.”
\loigiai{
Cần nêu được: Vectơ cảm ứng từ tổng hợp tại một điểm bằng tổng vectơ các cảm ứng từ thành phần. Khi vận dụng phải xác định đúng dữ kiện, phương chiều nếu có, điều kiện áp dụng, hệ đơn vị và kiểm tra tính hợp lí của kết quả. Nhận định “Độ lớn cảm ứng từ tổng hợp luôn bằng tổng độ lớn các cảm ứng từ thành phần.” sai vì trái với kiến thức nêu trên.
}
\end{bt}

% ===== Câu 60 =====
% ID: L12C3B14-20-TLN
% Mức: TH
\begin{ex}
% ID: L12C3B14-20-TLN
% Mức: TH
Trong bốn phát biểu sau về so sánh, tổng hợp và biểu diễn các từ trường đơn giản, có bao nhiêu phát biểu đúng? (Chỉ ghi một số nguyên.) (1) Độ lớn cảm ứng từ tổng hợp luôn bằng tổng độ lớn các cảm ứng từ thành phần. (2) Vectơ cảm ứng từ tổng hợp tại một điểm bằng tổng vectơ các cảm ứng từ thành phần. (3) Kết quả phải phù hợp ý nghĩa vật lí. (4) Mọi đại lượng đều là vô hướng.
\shortans{1}
\loigiai{
Đối chiếu từng phát biểu với kiến thức: Vectơ cảm ứng từ tổng hợp tại một điểm bằng tổng vectơ các cảm ứng từ thành phần. Các phát biểu đúng có số lượng là 1.
}
\end{ex}

% ===== Câu 61 =====
% ID: L12C3B14-20-TN
% Mức: VD
\dangbt{Mô tả đường sức từ và xác định chiều đường sức từ}
\begin{ex}
% ID: L12C3B14-20-TN
% Mức: VD
Trong một bài thực hành về mô tả đường sức từ và xác định chiều đường sức từ, nhận xét nào đúng?
\choice
{\True Tiếp tuyến với đường sức từ tại một điểm cho phương của vectơ cảm ứng từ; bên ngoài nam châm, đường sức đi từ cực Bắc đến cực Nam.}
{Các đường sức từ có thể cắt nhau tại một điểm.}
{Đại lượng đang xét không phụ thuộc các điều kiện vật lí nêu trong bài.}
{Có thể bỏ qua phương, chiều và đơn vị của các đại lượng trong mọi trường hợp.}
\loigiai{
Tiếp tuyến với đường sức từ tại một điểm cho phương của vectơ cảm ứng từ; bên ngoài nam châm, đường sức đi từ cực Bắc đến cực Nam. Vì vậy phương án $A$ đúng; các phương án còn lại trái với định nghĩa, định luật hoặc điều kiện áp dụng.
}
\end{ex}

% ===== Câu 62 =====
% ID: L12C3B14-21-DS
% Mức: TH
\dangbt{Từ trường của dòng điện thẳng, dòng điện tròn và ống dây}
\begin{ex}
% ID: L12C3B14-21-DS
% Mức: TH
Xét từ trường của dòng điện thẳng, dòng điện tròn và ống dây (Tình huống 1). Hãy xác định tính đúng, sai của bốn phát biểu sau.
\choiceTF
{\True Đường sức quanh dây dẫn thẳng là các đường tròn đồng tâm; từ trường trong lòng ống dây dài gần đều.}
{Từ trường trong lòng ống dây dài luôn bằng không.}
{\True Khi giải bài thuộc dạng này phải xác định đúng phương, chiều, điều kiện áp dụng và đơn vị của các đại lượng.}
{Trong dạng bài này, kết quả không phụ thuộc dữ kiện và có thể bỏ qua điều kiện áp dụng.}
\loigiai{
\begin{itemchoice}
\item (Đúng) Đường sức quanh dây dẫn thẳng là các đường tròn đồng tâm; từ trường trong lòng ống dây dài gần đều. (Vì): ường sức quanh dây dẫn thẳng là các đường tròn đồng tâm; từ trường trong lòng ống dây dài gần đều. b) Sai: Từ trường trong lòng ống dây dài luôn bằng không. c) Đúng: Khi giải bài thuộc dạng này phải xác định đúng phương, chiều, điều kiện áp dụng và đơn vị của các đại lượng. d) Sai: Trong dạng bài này, kết quả không phụ thuộc dữ kiện và có thể bỏ qua điều kiện áp dụng. Kiến thức cốt lõi: Đường sức quanh dây dẫn thẳng là các đường tròn đồng tâm; từ trường trong lòng ống dây dài gần đều
\item (Sai) Từ trường trong lòng ống dây dài luôn bằng không.
\item (Đúng) Khi giải bài thuộc dạng này phải xác định đúng phương, chiều, điều kiện áp dụng và đơn vị của các đại lượng.
\item (Sai) Trong dạng bài này, kết quả không phụ thuộc dữ kiện và có thể bỏ qua điều kiện áp dụng.
\end{itemchoice}
}
\end{ex}

% ===== Câu 63 =====
% ID: L12C3B14-21-TL
% Mức: VD
\dangbt{So sánh, tổng hợp và biểu diễn các từ trường đơn giản}
\begin{bt}
% ID: L12C3B14-21-TL
% Mức: VD
Phân tích sai lầm trong nhận định: “Độ lớn cảm ứng từ tổng hợp luôn bằng tổng độ lớn các cảm ứng từ thành phần.”
\loigiai{
Cần nêu được: Vectơ cảm ứng từ tổng hợp tại một điểm bằng tổng vectơ các cảm ứng từ thành phần. Khi vận dụng phải xác định đúng dữ kiện, phương chiều nếu có, điều kiện áp dụng, hệ đơn vị và kiểm tra tính hợp lí của kết quả. Nhận định “Độ lớn cảm ứng từ tổng hợp luôn bằng tổng độ lớn các cảm ứng từ thành phần.” sai vì trái với kiến thức nêu trên.
}
\end{bt}

% ===== Câu 64 =====
% ID: L12C3B14-21-TLN
% Mức: TH
\begin{ex}
% ID: L12C3B14-21-TLN
% Mức: TH
Trong bốn phát biểu sau về so sánh, tổng hợp và biểu diễn các từ trường đơn giản, có bao nhiêu phát biểu đúng? (Chỉ ghi một số nguyên.) (1) Cần đổi các đại lượng về hệ đơn vị thống nhất. (2) Vectơ cảm ứng từ tổng hợp tại một điểm bằng tổng vectơ các cảm ứng từ thành phần. (3) Độ lớn cảm ứng từ tổng hợp luôn bằng tổng độ lớn các cảm ứng từ thành phần. (4) Không cần xét chiều của đại lượng vectơ.
\shortans{2}
\loigiai{
Đối chiếu từng phát biểu với kiến thức: Vectơ cảm ứng từ tổng hợp tại một điểm bằng tổng vectơ các cảm ứng từ thành phần. Các phát biểu đúng có số lượng là 2.
}
\end{ex}

% ===== Câu 65 =====
% ID: L12C3B14-21-TN
% Mức: VD
\dangbt{Mô tả đường sức từ và xác định chiều đường sức từ}
\begin{ex}
% ID: L12C3B14-21-TN
% Mức: VD
Kết luận đúng khi vận dụng mô tả đường sức từ và xác định chiều đường sức từ là
\choice
{Các đường sức từ có thể cắt nhau tại một điểm.}
{\True Tiếp tuyến với đường sức từ tại một điểm cho phương của vectơ cảm ứng từ; bên ngoài nam châm, đường sức đi từ cực Bắc đến cực Nam.}
{Có thể bỏ qua phương, chiều và đơn vị của các đại lượng trong mọi trường hợp.}
{Đại lượng đang xét không phụ thuộc các điều kiện vật lí nêu trong bài.}
\loigiai{
Tiếp tuyến với đường sức từ tại một điểm cho phương của vectơ cảm ứng từ; bên ngoài nam châm, đường sức đi từ cực Bắc đến cực Nam. Vì vậy phương án $B$ đúng; các phương án còn lại trái với định nghĩa, định luật hoặc điều kiện áp dụng.
}
\end{ex}

% ===== Câu 66 =====
% ID: L12C3B14-22-DS
% Mức: TH
\dangbt{Từ trường của dòng điện thẳng, dòng điện tròn và ống dây}
\begin{ex}
% ID: L12C3B14-22-DS
% Mức: TH
Xét từ trường của dòng điện thẳng, dòng điện tròn và ống dây (Tình huống 2). Hãy xác định tính đúng, sai của bốn phát biểu sau.
\choiceTF
{Từ trường trong lòng ống dây dài luôn bằng không.}
{\True Đường sức quanh dây dẫn thẳng là các đường tròn đồng tâm; từ trường trong lòng ống dây dài gần đều.}
{Trong dạng bài này, kết quả không phụ thuộc dữ kiện và có thể bỏ qua điều kiện áp dụng.}
{\True Khi giải bài thuộc dạng này phải xác định đúng phương, chiều, điều kiện áp dụng và đơn vị của các đại lượng.}
\loigiai{
\begin{itemchoice}
\item (Sai) Từ trường trong lòng ống dây dài luôn bằng không. (Vì): Từ trường trong lòng ống dây dài luôn bằng không. b) Đúng: Đường sức quanh dây dẫn thẳng là các đường tròn đồng tâm; từ trường trong lòng ống dây dài gần đều. c) Sai: Trong dạng bài này, kết quả không phụ thuộc dữ kiện và có thể bỏ qua điều kiện áp dụng. d) Đúng: Khi giải bài thuộc dạng này phải xác định đúng phương, chiều, điều kiện áp dụng và đơn vị của các đại lượng. Kiến thức cốt lõi: Đường sức quanh dây dẫn thẳng là các đường tròn đồng tâm; từ trường trong lòng ống dây dài gần đều
\item (Đúng) Đường sức quanh dây dẫn thẳng là các đường tròn đồng tâm; từ trường trong lòng ống dây dài gần đều.
\item (Sai) Trong dạng bài này, kết quả không phụ thuộc dữ kiện và có thể bỏ qua điều kiện áp dụng.
\item (Đúng) Khi giải bài thuộc dạng này phải xác định đúng phương, chiều, điều kiện áp dụng và đơn vị của các đại lượng.
\end{itemchoice}
}
\end{ex}

% ===== Câu 67 =====
% ID: L12C3B14-22-TL
% Mức: VD
\dangbt{So sánh, tổng hợp và biểu diễn các từ trường đơn giản}
\begin{bt}
% ID: L12C3B14-22-TL
% Mức: VD
Đề xuất các bước giải một bài toán thuộc dạng so sánh, tổng hợp và biểu diễn các từ trường đơn giản.
\loigiai{
Cần nêu được: Vectơ cảm ứng từ tổng hợp tại một điểm bằng tổng vectơ các cảm ứng từ thành phần. Khi vận dụng phải xác định đúng dữ kiện, phương chiều nếu có, điều kiện áp dụng, hệ đơn vị và kiểm tra tính hợp lí của kết quả. Nhận định “Độ lớn cảm ứng từ tổng hợp luôn bằng tổng độ lớn các cảm ứng từ thành phần.” sai vì trái với kiến thức nêu trên.
}
\end{bt}

% ===== Câu 68 =====
% ID: L12C3B14-22-TLN
% Mức: VD
\begin{ex}
% ID: L12C3B14-22-TLN
% Mức: VD
Trong bốn phát biểu sau về so sánh, tổng hợp và biểu diễn các từ trường đơn giản, có bao nhiêu phát biểu đúng? (Chỉ ghi một số nguyên.) (1) Vectơ cảm ứng từ tổng hợp tại một điểm bằng tổng vectơ các cảm ứng từ thành phần. (2) Cần đối chiếu kết quả với điều kiện bài toán. (3) Độ lớn cảm ứng từ tổng hợp luôn bằng tổng độ lớn các cảm ứng từ thành phần. (4) Mọi trường hợp đều cho cùng một kết quả.
\shortans{2}
\loigiai{
Đối chiếu từng phát biểu với kiến thức: Vectơ cảm ứng từ tổng hợp tại một điểm bằng tổng vectơ các cảm ứng từ thành phần. Các phát biểu đúng có số lượng là 2.
}
\end{ex}

% ===== Câu 69 =====
% ID: L12C3B14-22-TN
% Mức: NB
\dangbt{Nhận biết từ trường, nguồn gây ra từ trường và tương tác từ}
\begin{ex}
% ID: L12C3B14-22-TN
% Mức: NB
Phát biểu nào sau đây đúng về nhận biết từ trường, nguồn gây ra từ trường và tương tác từ?
\choice
{\True Từ trường tồn tại xung quanh nam châm và dòng điện; nó gây lực từ lên nam châm hoặc dòng điện khác đặt trong đó.}
{Mọi điện tích đứng yên đều tạo ra từ trường giống dòng điện.}
{Đại lượng đang xét không phụ thuộc các điều kiện vật lí nêu trong bài.}
{Có thể bỏ qua phương, chiều và đơn vị của các đại lượng trong mọi trường hợp.}
\loigiai{
Từ trường tồn tại xung quanh nam châm và dòng điện; nó gây lực từ lên nam châm hoặc dòng điện khác đặt trong đó. Vì vậy phương án $A$ đúng; các phương án còn lại trái với định nghĩa, định luật hoặc điều kiện áp dụng.
}
\end{ex}

% ===== Câu 70 =====
% ID: L12C3B14-23-DS
% Mức: VD
\dangbt{Từ trường của dòng điện thẳng, dòng điện tròn và ống dây}
\begin{ex}
% ID: L12C3B14-23-DS
% Mức: VD
Xét từ trường của dòng điện thẳng, dòng điện tròn và ống dây (Tình huống 3). Hãy xác định tính đúng, sai của bốn phát biểu sau.
\choiceTF
{\True Khi giải bài thuộc dạng này phải xác định đúng phương, chiều, điều kiện áp dụng và đơn vị của các đại lượng.}
{Trong dạng bài này, kết quả không phụ thuộc dữ kiện và có thể bỏ qua điều kiện áp dụng.}
{\True Đường sức quanh dây dẫn thẳng là các đường tròn đồng tâm; từ trường trong lòng ống dây dài gần đều.}
{Từ trường trong lòng ống dây dài luôn bằng không.}
\loigiai{
\begin{itemchoice}
\item (Đúng) Khi giải bài thuộc dạng này phải xác định đúng phương, chiều, điều kiện áp dụng và đơn vị của các đại lượng. (Vì): Khi giải bài thuộc dạng này phải xác định đúng phương, chiều, điều kiện áp dụng và đơn vị của các đại lượng. b) Sai: Trong dạng bài này, kết quả không phụ thuộc dữ kiện và có thể bỏ qua điều kiện áp dụng. c) Đúng: Đường sức quanh dây dẫn thẳng là các đường tròn đồng tâm; từ trường trong lòng ống dây dài gần đều. d) Sai: Từ trường trong lòng ống dây dài luôn bằng không. Kiến thức cốt lõi: Đường sức quanh dây dẫn thẳng là các đường tròn đồng tâm; từ trường trong lòng ống dây dài gần đều
\item (Sai) Trong dạng bài này, kết quả không phụ thuộc dữ kiện và có thể bỏ qua điều kiện áp dụng.
\item (Đúng) Đường sức quanh dây dẫn thẳng là các đường tròn đồng tâm; từ trường trong lòng ống dây dài gần đều.
\item (Sai) Từ trường trong lòng ống dây dài luôn bằng không.
\end{itemchoice}
}
\end{ex}

% ===== Câu 71 =====
% ID: L12C3B14-23-TL
% Mức: VD
\dangbt{So sánh, tổng hợp và biểu diễn các từ trường đơn giản}
\begin{bt}
% ID: L12C3B14-23-TL
% Mức: VD
Nêu một tình huống thực tiễn liên quan đến so sánh, tổng hợp và biểu diễn các từ trường đơn giản và giải thích bằng kiến thức Vật lí.
\loigiai{
Cần nêu được: Vectơ cảm ứng từ tổng hợp tại một điểm bằng tổng vectơ các cảm ứng từ thành phần. Khi vận dụng phải xác định đúng dữ kiện, phương chiều nếu có, điều kiện áp dụng, hệ đơn vị và kiểm tra tính hợp lí của kết quả. Nhận định “Độ lớn cảm ứng từ tổng hợp luôn bằng tổng độ lớn các cảm ứng từ thành phần.” sai vì trái với kiến thức nêu trên.
}
\end{bt}

% ===== Câu 72 =====
% ID: L12C3B14-23-TLN
% Mức: VD
\begin{ex}
% ID: L12C3B14-23-TLN
% Mức: VD
Trong bốn phát biểu sau về so sánh, tổng hợp và biểu diễn các từ trường đơn giản, có bao nhiêu phát biểu đúng? (Chỉ ghi một số nguyên.) (1) Độ lớn cảm ứng từ tổng hợp luôn bằng tổng độ lớn các cảm ứng từ thành phần. (2) Cần đọc đúng dữ kiện và giả thiết. (3) Vectơ cảm ứng từ tổng hợp tại một điểm bằng tổng vectơ các cảm ứng từ thành phần. (4) Có thể thay số trước khi xác định công thức.
\shortans{2}
\loigiai{
Đối chiếu từng phát biểu với kiến thức: Vectơ cảm ứng từ tổng hợp tại một điểm bằng tổng vectơ các cảm ứng từ thành phần. Các phát biểu đúng có số lượng là 2.
}
\end{ex}

% ===== Câu 73 =====
% ID: L12C3B14-23-TN
% Mức: NB
\dangbt{Nhận biết từ trường, nguồn gây ra từ trường và tương tác từ}
\begin{ex}
% ID: L12C3B14-23-TN
% Mức: NB
Trong nhận biết từ trường, nguồn gây ra từ trường và tương tác từ, nhận định nào phù hợp với kiến thức Vật lí?
\choice
{Mọi điện tích đứng yên đều tạo ra từ trường giống dòng điện.}
{\True Từ trường tồn tại xung quanh nam châm và dòng điện; nó gây lực từ lên nam châm hoặc dòng điện khác đặt trong đó.}
{Có thể bỏ qua phương, chiều và đơn vị của các đại lượng trong mọi trường hợp.}
{Đại lượng đang xét không phụ thuộc các điều kiện vật lí nêu trong bài.}
\loigiai{
Từ trường tồn tại xung quanh nam châm và dòng điện; nó gây lực từ lên nam châm hoặc dòng điện khác đặt trong đó. Vì vậy phương án $B$ đúng; các phương án còn lại trái với định nghĩa, định luật hoặc điều kiện áp dụng.
}
\end{ex}

% ===== Câu 74 =====
% ID: L12C3B14-24-DS
% Mức: VD
\dangbt{Từ trường của dòng điện thẳng, dòng điện tròn và ống dây}
\begin{ex}
% ID: L12C3B14-24-DS
% Mức: VD
Xét từ trường của dòng điện thẳng, dòng điện tròn và ống dây (Tình huống 4). Hãy xác định tính đúng, sai của bốn phát biểu sau.
\choiceTF
{Trong dạng bài này, kết quả không phụ thuộc dữ kiện và có thể bỏ qua điều kiện áp dụng.}
{\True Khi giải bài thuộc dạng này phải xác định đúng phương, chiều, điều kiện áp dụng và đơn vị của các đại lượng.}
{Từ trường trong lòng ống dây dài luôn bằng không.}
{\True Đường sức quanh dây dẫn thẳng là các đường tròn đồng tâm; từ trường trong lòng ống dây dài gần đều.}
\loigiai{
\begin{itemchoice}
\item (Sai) Trong dạng bài này, kết quả không phụ thuộc dữ kiện và có thể bỏ qua điều kiện áp dụng. (Vì): Trong dạng bài này, kết quả không phụ thuộc dữ kiện và có thể bỏ qua điều kiện áp dụng. b) Đúng: Khi giải bài thuộc dạng này phải xác định đúng phương, chiều, điều kiện áp dụng và đơn vị của các đại lượng. c) Sai: Từ trường trong lòng ống dây dài luôn bằng không. d) Đúng: Đường sức quanh dây dẫn thẳng là các đường tròn đồng tâm; từ trường trong lòng ống dây dài gần đều. Kiến thức cốt lõi: Đường sức quanh dây dẫn thẳng là các đường tròn đồng tâm; từ trường trong lòng ống dây dài gần đều
\item (Đúng) Khi giải bài thuộc dạng này phải xác định đúng phương, chiều, điều kiện áp dụng và đơn vị của các đại lượng.
\item (Sai) Từ trường trong lòng ống dây dài luôn bằng không.
\item (Đúng) Đường sức quanh dây dẫn thẳng là các đường tròn đồng tâm; từ trường trong lòng ống dây dài gần đều.
\end{itemchoice}
}
\end{ex}

% ===== Câu 75 =====
% ID: L12C3B14-24-TL
% Mức: VDC
\dangbt{So sánh, tổng hợp và biểu diễn các từ trường đơn giản}
\begin{bt}
% ID: L12C3B14-24-TL
% Mức: VDC
So sánh nhận định đúng “Vectơ cảm ứng từ tổng hợp tại một điểm bằng tổng vectơ các cảm ứng từ thành phần.” với nhận định sai “Độ lớn cảm ứng từ tổng hợp luôn bằng tổng độ lớn các cảm ứng từ thành phần.”; chỉ rõ căn cứ Vật lí.
\loigiai{
Cần nêu được: Vectơ cảm ứng từ tổng hợp tại một điểm bằng tổng vectơ các cảm ứng từ thành phần. Khi vận dụng phải xác định đúng dữ kiện, phương chiều nếu có, điều kiện áp dụng, hệ đơn vị và kiểm tra tính hợp lí của kết quả. Nhận định “Độ lớn cảm ứng từ tổng hợp luôn bằng tổng độ lớn các cảm ứng từ thành phần.” sai vì trái với kiến thức nêu trên.
}
\end{bt}

% ===== Câu 76 =====
% ID: L12C3B14-24-TLN
% Mức: VDC
\begin{ex}
% ID: L12C3B14-24-TLN
% Mức: VDC
Trong bốn phát biểu sau về so sánh, tổng hợp và biểu diễn các từ trường đơn giản, có bao nhiêu phát biểu đúng? (Chỉ ghi một số nguyên.) (1) Kết quả cần có ý nghĩa vật lí. (2) Độ lớn cảm ứng từ tổng hợp luôn bằng tổng độ lớn các cảm ứng từ thành phần. (3) Phải dùng đúng định luật cho tình huống. (4) Vectơ cảm ứng từ tổng hợp tại một điểm bằng tổng vectơ các cảm ứng từ thành phần.
\shortans{3}
\loigiai{
Đối chiếu từng phát biểu với kiến thức: Vectơ cảm ứng từ tổng hợp tại một điểm bằng tổng vectơ các cảm ứng từ thành phần. Các phát biểu đúng có số lượng là 3.
}
\end{ex}

% ===== Câu 77 =====
% ID: L12C3B14-24-TN
% Mức: NB
\dangbt{Nhận biết từ trường, nguồn gây ra từ trường và tương tác từ}
\begin{ex}
% ID: L12C3B14-24-TN
% Mức: NB
Tình huống 3: chọn kết luận chính xác về nhận biết từ trường, nguồn gây ra từ trường và tương tác từ.
\choice
{Đại lượng đang xét không phụ thuộc các điều kiện vật lí nêu trong bài.}
{Có thể bỏ qua phương, chiều và đơn vị của các đại lượng trong mọi trường hợp.}
{\True Từ trường tồn tại xung quanh nam châm và dòng điện; nó gây lực từ lên nam châm hoặc dòng điện khác đặt trong đó.}
{Mọi điện tích đứng yên đều tạo ra từ trường giống dòng điện.}
\loigiai{
Từ trường tồn tại xung quanh nam châm và dòng điện; nó gây lực từ lên nam châm hoặc dòng điện khác đặt trong đó. Vì vậy phương án $C$ đúng; các phương án còn lại trái với định nghĩa, định luật hoặc điều kiện áp dụng.
}
\end{ex}

% ===== Câu 78 =====
% ID: L12C3B14-25-DS
% Mức: VD
\dangbt{Từ trường của dòng điện thẳng, dòng điện tròn và ống dây}
\begin{ex}
% ID: L12C3B14-25-DS
% Mức: VD
Xét từ trường của dòng điện thẳng, dòng điện tròn và ống dây (Tình huống 5). Hãy xác định tính đúng, sai của bốn phát biểu sau.
\choiceTF
{\True Đường sức quanh dây dẫn thẳng là các đường tròn đồng tâm; từ trường trong lòng ống dây dài gần đều.}
{Trong dạng bài này, kết quả không phụ thuộc dữ kiện và có thể bỏ qua điều kiện áp dụng.}
{Từ trường trong lòng ống dây dài luôn bằng không.}
{\True Khi giải bài thuộc dạng này phải xác định đúng phương, chiều, điều kiện áp dụng và đơn vị của các đại lượng.}
\loigiai{
\begin{itemchoice}
\item (Đúng) Đường sức quanh dây dẫn thẳng là các đường tròn đồng tâm; từ trường trong lòng ống dây dài gần đều. (Vì): ường sức quanh dây dẫn thẳng là các đường tròn đồng tâm; từ trường trong lòng ống dây dài gần đều. b) Sai: Trong dạng bài này, kết quả không phụ thuộc dữ kiện và có thể bỏ qua điều kiện áp dụng. c) Sai: Từ trường trong lòng ống dây dài luôn bằng không. d) Đúng: Khi giải bài thuộc dạng này phải xác định đúng phương, chiều, điều kiện áp dụng và đơn vị của các đại lượng. Kiến thức cốt lõi: Đường sức quanh dây dẫn thẳng là các đường tròn đồng tâm; từ trường trong lòng ống dây dài gần đều
\item (Sai) Trong dạng bài này, kết quả không phụ thuộc dữ kiện và có thể bỏ qua điều kiện áp dụng.
\item (Sai) Từ trường trong lòng ống dây dài luôn bằng không.
\item (Đúng) Khi giải bài thuộc dạng này phải xác định đúng phương, chiều, điều kiện áp dụng và đơn vị của các đại lượng.
\end{itemchoice}
}
\end{ex}

% ===== Câu 79 =====
% ID: L12C3B14-25-TL
% Mức: TH
\begin{bt}
% ID: L12C3B14-25-TL
% Mức: TH
Trình bày kiến thức cốt lõi của từ trường của dòng điện thẳng, dòng điện tròn và ống dây và nêu điều kiện áp dụng.
\loigiai{
Cần nêu được: Đường sức quanh dây dẫn thẳng là các đường tròn đồng tâm; từ trường trong lòng ống dây dài gần đều. Khi vận dụng phải xác định đúng dữ kiện, phương chiều nếu có, điều kiện áp dụng, hệ đơn vị và kiểm tra tính hợp lí của kết quả. Nhận định “Từ trường trong lòng ống dây dài luôn bằng không.” sai vì trái với kiến thức nêu trên.
}
\end{bt}

% ===== Câu 80 =====
% ID: L12C3B14-25-TLN
% Mức: TH
\begin{ex}
% ID: L12C3B14-25-TLN
% Mức: TH
Trong bốn phát biểu sau về từ trường của dòng điện thẳng, dòng điện tròn và ống dây, có bao nhiêu phát biểu đúng? (Chỉ ghi một số nguyên.) (1) Đường sức quanh dây dẫn thẳng là các đường tròn đồng tâm; từ trường trong lòng ống dây dài gần đều. (2) Từ trường trong lòng ống dây dài luôn bằng không. (3) Cần kiểm tra điều kiện áp dụng trước khi tính toán. (4) Có thể bỏ qua đơn vị của kết quả.
\shortans{2}
\loigiai{
Đối chiếu từng phát biểu với kiến thức: Đường sức quanh dây dẫn thẳng là các đường tròn đồng tâm; từ trường trong lòng ống dây dài gần đều. Các phát biểu đúng có số lượng là 2.
}
\end{ex}

% ===== Câu 81 =====
% ID: L12C3B14-25-TN
% Mức: NB
\dangbt{Nhận biết từ trường, nguồn gây ra từ trường và tương tác từ}
\begin{ex}
% ID: L12C3B14-25-TN
% Mức: NB
Nếu một proton đứng yên, nó sẽ
\choice
{tạo ra từ trường khi nó nằm gần một dây dẫn.}
{không tạo ra điện trường.}
{\True không tạo ra từ trường.}
{tạo ra từ trường giống như một nam châm nhỏ.}
\loigiai{
Điện tích điểm đứng yên chỉ tạo ra điện trường xung quanh nó. Từ trường chỉ được tạo ra bởi điện tích chuyển động (dòng điện).
}
\end{ex}

% ===== Câu 82 =====
% ID: L12C3B14-26-DS
% Mức: TH
\dangbt{Xác định chiều từ trường bằng kim nam châm và quy tắc nắm tay phải}
\begin{ex}
% ID: L12C3B14-26-DS
% Mức: TH
Xét xác định chiều từ trường bằng kim nam châm và quy tắc nắm tay phải (Tình huống 1). Hãy xác định tính đúng, sai của bốn phát biểu sau.
\choiceTF
{\True Đầu Bắc của kim nam châm nhỏ chỉ theo chiều của vectơ cảm ứng từ; quy tắc nắm tay phải xác định chiều đường sức quanh dòng điện.}
{Quy tắc bàn tay trái dùng để xác định chiều đường sức từ quanh dây dẫn thẳng.}
{\True Khi giải bài thuộc dạng này phải xác định đúng phương, chiều, điều kiện áp dụng và đơn vị của các đại lượng.}
{Trong dạng bài này, kết quả không phụ thuộc dữ kiện và có thể bỏ qua điều kiện áp dụng.}
\loigiai{
\begin{itemchoice}
\item (Đúng) Đầu Bắc của kim nam châm nhỏ chỉ theo chiều của vectơ cảm ứng từ; quy tắc nắm tay phải xác định chiều đường sức quanh dòng điện. (Vì): ầu Bắc của kim nam châm nhỏ chỉ theo chiều của vectơ cảm ứng từ; quy tắc nắm tay phải xác định chiều đường sức quanh dòng điện. b) Sai: Quy tắc bàn tay trái dùng để xác định chiều đường sức từ quanh dây dẫn thẳng. c) Đúng: Khi giải bài thuộc dạng này phải xác định đúng phương, chiều, điều kiện áp dụng và đơn vị của các đại lượng. d) Sai: Trong dạng bài này, kết quả không phụ thuộc dữ kiện và có thể bỏ qua điều kiện áp dụng. Kiến thức cốt lõi: Đầu Bắc của kim nam châm nhỏ chỉ theo chiều của vectơ cảm ứng từ; quy tắc nắm tay phải xác định chiều đường sức quanh dòng điện
\item (Sai) Quy tắc bàn tay trái dùng để xác định chiều đường sức từ quanh dây dẫn thẳng.
\item (Đúng) Khi giải bài thuộc dạng này phải xác định đúng phương, chiều, điều kiện áp dụng và đơn vị của các đại lượng.
\item (Sai) Trong dạng bài này, kết quả không phụ thuộc dữ kiện và có thể bỏ qua điều kiện áp dụng.
\end{itemchoice}
}
\end{ex}

% ===== Câu 83 =====
% ID: L12C3B14-26-TL
% Mức: TH
\dangbt{Từ trường của dòng điện thẳng, dòng điện tròn và ống dây}
\begin{bt}
% ID: L12C3B14-26-TL
% Mức: TH
Giải thích vì sao nhận định sau đúng: “Đường sức quanh dây dẫn thẳng là các đường tròn đồng tâm; từ trường trong lòng ống dây dài gần đều.”
\loigiai{
Cần nêu được: Đường sức quanh dây dẫn thẳng là các đường tròn đồng tâm; từ trường trong lòng ống dây dài gần đều. Khi vận dụng phải xác định đúng dữ kiện, phương chiều nếu có, điều kiện áp dụng, hệ đơn vị và kiểm tra tính hợp lí của kết quả. Nhận định “Từ trường trong lòng ống dây dài luôn bằng không.” sai vì trái với kiến thức nêu trên.
}
\end{bt}

% ===== Câu 84 =====
% ID: L12C3B14-26-TLN
% Mức: TH
\begin{ex}
% ID: L12C3B14-26-TLN
% Mức: TH
Trong bốn phát biểu sau về từ trường của dòng điện thẳng, dòng điện tròn và ống dây, có bao nhiêu phát biểu đúng? (Chỉ ghi một số nguyên.) (1) Từ trường trong lòng ống dây dài luôn bằng không. (2) Đường sức quanh dây dẫn thẳng là các đường tròn đồng tâm; từ trường trong lòng ống dây dài gần đều. (3) Kết quả phải phù hợp ý nghĩa vật lí. (4) Mọi đại lượng đều là vô hướng.
\shortans{1}
\loigiai{
Đối chiếu từng phát biểu với kiến thức: Đường sức quanh dây dẫn thẳng là các đường tròn đồng tâm; từ trường trong lòng ống dây dài gần đều. Các phát biểu đúng có số lượng là 1.
}
\end{ex}

% ===== Câu 85 =====
% ID: L12C3B14-26-TN
% Mức: TH
\dangbt{Nhận biết từ trường, nguồn gây ra từ trường và tương tác từ}
\begin{ex}
% ID: L12C3B14-26-TN
% Mức: TH
Nội dung nào mô tả đúng nhất nhận biết từ trường, nguồn gây ra từ trường và tương tác từ?
\choice
{Có thể bỏ qua phương, chiều và đơn vị của các đại lượng trong mọi trường hợp.}
{Đại lượng đang xét không phụ thuộc các điều kiện vật lí nêu trong bài.}
{Mọi điện tích đứng yên đều tạo ra từ trường giống dòng điện.}
{\True Từ trường tồn tại xung quanh nam châm và dòng điện; nó gây lực từ lên nam châm hoặc dòng điện khác đặt trong đó.}
\loigiai{
Từ trường tồn tại xung quanh nam châm và dòng điện; nó gây lực từ lên nam châm hoặc dòng điện khác đặt trong đó. Vì vậy phương án $D$ đúng; các phương án còn lại trái với định nghĩa, định luật hoặc điều kiện áp dụng.
}
\end{ex}

% ===== Câu 86 =====
% ID: L12C3B14-27-DS
% Mức: TH
\dangbt{Xác định chiều từ trường bằng kim nam châm và quy tắc nắm tay phải}
\begin{ex}
% ID: L12C3B14-27-DS
% Mức: TH
Xét xác định chiều từ trường bằng kim nam châm và quy tắc nắm tay phải (Tình huống 2). Hãy xác định tính đúng, sai của bốn phát biểu sau.
\choiceTF
{Quy tắc bàn tay trái dùng để xác định chiều đường sức từ quanh dây dẫn thẳng.}
{\True Đầu Bắc của kim nam châm nhỏ chỉ theo chiều của vectơ cảm ứng từ; quy tắc nắm tay phải xác định chiều đường sức quanh dòng điện.}
{Trong dạng bài này, kết quả không phụ thuộc dữ kiện và có thể bỏ qua điều kiện áp dụng.}
{\True Khi giải bài thuộc dạng này phải xác định đúng phương, chiều, điều kiện áp dụng và đơn vị của các đại lượng.}
\loigiai{
\begin{itemchoice}
\item (Sai) Quy tắc bàn tay trái dùng để xác định chiều đường sức từ quanh dây dẫn thẳng. (Vì): Quy tắc bàn tay trái dùng để xác định chiều đường sức từ quanh dây dẫn thẳng. b) Đúng: Đầu Bắc của kim nam châm nhỏ chỉ theo chiều của vectơ cảm ứng từ; quy tắc nắm tay phải xác định chiều đường sức quanh dòng điện. c) Sai: Trong dạng bài này, kết quả không phụ thuộc dữ kiện và có thể bỏ qua điều kiện áp dụng. d) Đúng: Khi giải bài thuộc dạng này phải xác định đúng phương, chiều, điều kiện áp dụng và đơn vị của các đại lượng. Kiến thức cốt lõi: Đầu Bắc của kim nam châm nhỏ chỉ theo chiều của vectơ cảm ứng từ; quy tắc nắm tay phải xác định chiều đường sức quanh dòng điện
\item (Đúng) Đầu Bắc của kim nam châm nhỏ chỉ theo chiều của vectơ cảm ứng từ; quy tắc nắm tay phải xác định chiều đường sức quanh dòng điện.
\item (Sai) Trong dạng bài này, kết quả không phụ thuộc dữ kiện và có thể bỏ qua điều kiện áp dụng.
\item (Đúng) Khi giải bài thuộc dạng này phải xác định đúng phương, chiều, điều kiện áp dụng và đơn vị của các đại lượng.
\end{itemchoice}
}
\end{ex}

% ===== Câu 87 =====
% ID: L12C3B14-27-TL
% Mức: VD
\dangbt{Từ trường của dòng điện thẳng, dòng điện tròn và ống dây}
\begin{bt}
% ID: L12C3B14-27-TL
% Mức: VD
Phân tích sai lầm trong nhận định: “Từ trường trong lòng ống dây dài luôn bằng không.”
\loigiai{
Cần nêu được: Đường sức quanh dây dẫn thẳng là các đường tròn đồng tâm; từ trường trong lòng ống dây dài gần đều. Khi vận dụng phải xác định đúng dữ kiện, phương chiều nếu có, điều kiện áp dụng, hệ đơn vị và kiểm tra tính hợp lí của kết quả. Nhận định “Từ trường trong lòng ống dây dài luôn bằng không.” sai vì trái với kiến thức nêu trên.
}
\end{bt}

% ===== Câu 88 =====
% ID: L12C3B14-27-TLN
% Mức: TH
\begin{ex}
% ID: L12C3B14-27-TLN
% Mức: TH
Trong bốn phát biểu sau về từ trường của dòng điện thẳng, dòng điện tròn và ống dây, có bao nhiêu phát biểu đúng? (Chỉ ghi một số nguyên.) (1) Cần đổi các đại lượng về hệ đơn vị thống nhất. (2) Đường sức quanh dây dẫn thẳng là các đường tròn đồng tâm; từ trường trong lòng ống dây dài gần đều. (3) Từ trường trong lòng ống dây dài luôn bằng không. (4) Không cần xét chiều của đại lượng vectơ.
\shortans{2}
\loigiai{
Đối chiếu từng phát biểu với kiến thức: Đường sức quanh dây dẫn thẳng là các đường tròn đồng tâm; từ trường trong lòng ống dây dài gần đều. Các phát biểu đúng có số lượng là 2.
}
\end{ex}

% ===== Câu 89 =====
% ID: L12C3B14-27-TN
% Mức: TH
\dangbt{Nhận biết từ trường, nguồn gây ra từ trường và tương tác từ}
\begin{ex}
% ID: L12C3B14-27-TN
% Mức: TH
Khi phân tích nhận biết từ trường, nguồn gây ra từ trường và tương tác từ, phát biểu nào đúng?
\choice
{\True Từ trường tồn tại xung quanh nam châm và dòng điện; nó gây lực từ lên nam châm hoặc dòng điện khác đặt trong đó.}
{Mọi điện tích đứng yên đều tạo ra từ trường giống dòng điện.}
{Đại lượng đang xét không phụ thuộc các điều kiện vật lí nêu trong bài.}
{Có thể bỏ qua phương, chiều và đơn vị của các đại lượng trong mọi trường hợp.}
\loigiai{
Từ trường tồn tại xung quanh nam châm và dòng điện; nó gây lực từ lên nam châm hoặc dòng điện khác đặt trong đó. Vì vậy phương án $A$ đúng; các phương án còn lại trái với định nghĩa, định luật hoặc điều kiện áp dụng.
}
\end{ex}

% ===== Câu 90 =====
% ID: L12C3B14-28-DS
% Mức: VD
\dangbt{Xác định chiều từ trường bằng kim nam châm và quy tắc nắm tay phải}
\begin{ex}
% ID: L12C3B14-28-DS
% Mức: VD
Xét xác định chiều từ trường bằng kim nam châm và quy tắc nắm tay phải (Tình huống 3). Hãy xác định tính đúng, sai của bốn phát biểu sau.
\choiceTF
{\True Khi giải bài thuộc dạng này phải xác định đúng phương, chiều, điều kiện áp dụng và đơn vị của các đại lượng.}
{Trong dạng bài này, kết quả không phụ thuộc dữ kiện và có thể bỏ qua điều kiện áp dụng.}
{\True Đầu Bắc của kim nam châm nhỏ chỉ theo chiều của vectơ cảm ứng từ; quy tắc nắm tay phải xác định chiều đường sức quanh dòng điện.}
{Quy tắc bàn tay trái dùng để xác định chiều đường sức từ quanh dây dẫn thẳng.}
\loigiai{
\begin{itemchoice}
\item (Đúng) Khi giải bài thuộc dạng này phải xác định đúng phương, chiều, điều kiện áp dụng và đơn vị của các đại lượng. (Vì): Khi giải bài thuộc dạng này phải xác định đúng phương, chiều, điều kiện áp dụng và đơn vị của các đại lượng. b) Sai: Trong dạng bài này, kết quả không phụ thuộc dữ kiện và có thể bỏ qua điều kiện áp dụng. c) Đúng: Đầu Bắc của kim nam châm nhỏ chỉ theo chiều của vectơ cảm ứng từ; quy tắc nắm tay phải xác định chiều đường sức quanh dòng điện. d) Sai: Quy tắc bàn tay trái dùng để xác định chiều đường sức từ quanh dây dẫn thẳng. Kiến thức cốt lõi: Đầu Bắc của kim nam châm nhỏ chỉ theo chiều của vectơ cảm ứng từ; quy tắc nắm tay phải xác định chiều đường sức quanh dòng điện
\item (Sai) Trong dạng bài này, kết quả không phụ thuộc dữ kiện và có thể bỏ qua điều kiện áp dụng.
\item (Đúng) Đầu Bắc của kim nam châm nhỏ chỉ theo chiều của vectơ cảm ứng từ; quy tắc nắm tay phải xác định chiều đường sức quanh dòng điện.
\item (Sai) Quy tắc bàn tay trái dùng để xác định chiều đường sức từ quanh dây dẫn thẳng.
\end{itemchoice}
}
\end{ex}

% ===== Câu 91 =====
% ID: L12C3B14-28-TL
% Mức: VD
\dangbt{Từ trường của dòng điện thẳng, dòng điện tròn và ống dây}
\begin{bt}
% ID: L12C3B14-28-TL
% Mức: VD
Đề xuất các bước giải một bài toán thuộc dạng từ trường của dòng điện thẳng, dòng điện tròn và ống dây.
\loigiai{
Cần nêu được: Đường sức quanh dây dẫn thẳng là các đường tròn đồng tâm; từ trường trong lòng ống dây dài gần đều. Khi vận dụng phải xác định đúng dữ kiện, phương chiều nếu có, điều kiện áp dụng, hệ đơn vị và kiểm tra tính hợp lí của kết quả. Nhận định “Từ trường trong lòng ống dây dài luôn bằng không.” sai vì trái với kiến thức nêu trên.
}
\end{bt}

% ===== Câu 92 =====
% ID: L12C3B14-28-TLN
% Mức: VD
\begin{ex}
% ID: L12C3B14-28-TLN
% Mức: VD
Trong bốn phát biểu sau về từ trường của dòng điện thẳng, dòng điện tròn và ống dây, có bao nhiêu phát biểu đúng? (Chỉ ghi một số nguyên.) (1) Đường sức quanh dây dẫn thẳng là các đường tròn đồng tâm; từ trường trong lòng ống dây dài gần đều. (2) Cần đối chiếu kết quả với điều kiện bài toán. (3) Từ trường trong lòng ống dây dài luôn bằng không. (4) Mọi trường hợp đều cho cùng một kết quả.
\shortans{2}
\loigiai{
Đối chiếu từng phát biểu với kiến thức: Đường sức quanh dây dẫn thẳng là các đường tròn đồng tâm; từ trường trong lòng ống dây dài gần đều. Các phát biểu đúng có số lượng là 2.
}
\end{ex}

% ===== Câu 93 =====
% ID: L12C3B14-28-TN
% Mức: TH
\dangbt{Nhận biết từ trường, nguồn gây ra từ trường và tương tác từ}
\begin{ex}
% ID: L12C3B14-28-TN
% Mức: TH
Một học sinh ôn tập nhận biết từ trường, nguồn gây ra từ trường và tương tác từ. Kết luận nào của học sinh là đúng?
\choice
{Mọi điện tích đứng yên đều tạo ra từ trường giống dòng điện.}
{\True Từ trường tồn tại xung quanh nam châm và dòng điện; nó gây lực từ lên nam châm hoặc dòng điện khác đặt trong đó.}
{Có thể bỏ qua phương, chiều và đơn vị của các đại lượng trong mọi trường hợp.}
{Đại lượng đang xét không phụ thuộc các điều kiện vật lí nêu trong bài.}
\loigiai{
Từ trường tồn tại xung quanh nam châm và dòng điện; nó gây lực từ lên nam châm hoặc dòng điện khác đặt trong đó. Vì vậy phương án $B$ đúng; các phương án còn lại trái với định nghĩa, định luật hoặc điều kiện áp dụng.
}
\end{ex}

% ===== Câu 94 =====
% ID: L12C3B14-29-DS
% Mức: VD
\dangbt{Xác định chiều từ trường bằng kim nam châm và quy tắc nắm tay phải}
\begin{ex}
% ID: L12C3B14-29-DS
% Mức: VD
Xét xác định chiều từ trường bằng kim nam châm và quy tắc nắm tay phải (Tình huống 4). Hãy xác định tính đúng, sai của bốn phát biểu sau.
\choiceTF
{Trong dạng bài này, kết quả không phụ thuộc dữ kiện và có thể bỏ qua điều kiện áp dụng.}
{\True Khi giải bài thuộc dạng này phải xác định đúng phương, chiều, điều kiện áp dụng và đơn vị của các đại lượng.}
{Quy tắc bàn tay trái dùng để xác định chiều đường sức từ quanh dây dẫn thẳng.}
{\True Đầu Bắc của kim nam châm nhỏ chỉ theo chiều của vectơ cảm ứng từ; quy tắc nắm tay phải xác định chiều đường sức quanh dòng điện.}
\loigiai{
\begin{itemchoice}
\item (Sai) Trong dạng bài này, kết quả không phụ thuộc dữ kiện và có thể bỏ qua điều kiện áp dụng. (Vì): Trong dạng bài này, kết quả không phụ thuộc dữ kiện và có thể bỏ qua điều kiện áp dụng. b) Đúng: Khi giải bài thuộc dạng này phải xác định đúng phương, chiều, điều kiện áp dụng và đơn vị của các đại lượng. c) Sai: Quy tắc bàn tay trái dùng để xác định chiều đường sức từ quanh dây dẫn thẳng. d) Đúng: Đầu Bắc của kim nam châm nhỏ chỉ theo chiều của vectơ cảm ứng từ; quy tắc nắm tay phải xác định chiều đường sức quanh dòng điện. Kiến thức cốt lõi: Đầu Bắc của kim nam châm nhỏ chỉ theo chiều của vectơ cảm ứng từ; quy tắc nắm tay phải xác định chiều đường sức quanh dòng điện
\item (Đúng) Khi giải bài thuộc dạng này phải xác định đúng phương, chiều, điều kiện áp dụng và đơn vị của các đại lượng.
\item (Sai) Quy tắc bàn tay trái dùng để xác định chiều đường sức từ quanh dây dẫn thẳng.
\item (Đúng) Đầu Bắc của kim nam châm nhỏ chỉ theo chiều của vectơ cảm ứng từ; quy tắc nắm tay phải xác định chiều đường sức quanh dòng điện.
\end{itemchoice}
}
\end{ex}

% ===== Câu 95 =====
% ID: L12C3B14-29-TL
% Mức: VD
\dangbt{Từ trường của dòng điện thẳng, dòng điện tròn và ống dây}
\begin{bt}
% ID: L12C3B14-29-TL
% Mức: VD
Nêu một tình huống thực tiễn liên quan đến từ trường của dòng điện thẳng, dòng điện tròn và ống dây và giải thích bằng kiến thức Vật lí.
\loigiai{
Cần nêu được: Đường sức quanh dây dẫn thẳng là các đường tròn đồng tâm; từ trường trong lòng ống dây dài gần đều. Khi vận dụng phải xác định đúng dữ kiện, phương chiều nếu có, điều kiện áp dụng, hệ đơn vị và kiểm tra tính hợp lí của kết quả. Nhận định “Từ trường trong lòng ống dây dài luôn bằng không.” sai vì trái với kiến thức nêu trên.
}
\end{bt}

% ===== Câu 96 =====
% ID: L12C3B14-29-TLN
% Mức: VD
\begin{ex}
% ID: L12C3B14-29-TLN
% Mức: VD
Trong bốn phát biểu sau về từ trường của dòng điện thẳng, dòng điện tròn và ống dây, có bao nhiêu phát biểu đúng? (Chỉ ghi một số nguyên.) (1) Từ trường trong lòng ống dây dài luôn bằng không. (2) Cần đọc đúng dữ kiện và giả thiết. (3) Đường sức quanh dây dẫn thẳng là các đường tròn đồng tâm; từ trường trong lòng ống dây dài gần đều. (4) Có thể thay số trước khi xác định công thức.
\shortans{2}
\loigiai{
Đối chiếu từng phát biểu với kiến thức: Đường sức quanh dây dẫn thẳng là các đường tròn đồng tâm; từ trường trong lòng ống dây dài gần đều. Các phát biểu đúng có số lượng là 2.
}
\end{ex}

% ===== Câu 97 =====
% ID: L12C3B14-29-TN
% Mức: TH
\dangbt{Nhận biết từ trường, nguồn gây ra từ trường và tương tác từ}
\begin{ex}
% ID: L12C3B14-29-TN
% Mức: TH
Chọn phát biểu không trái với kiến thức về nhận biết từ trường, nguồn gây ra từ trường và tương tác từ.
\choice
{Đại lượng đang xét không phụ thuộc các điều kiện vật lí nêu trong bài.}
{Có thể bỏ qua phương, chiều và đơn vị của các đại lượng trong mọi trường hợp.}
{\True Từ trường tồn tại xung quanh nam châm và dòng điện; nó gây lực từ lên nam châm hoặc dòng điện khác đặt trong đó.}
{Mọi điện tích đứng yên đều tạo ra từ trường giống dòng điện.}
\loigiai{
Từ trường tồn tại xung quanh nam châm và dòng điện; nó gây lực từ lên nam châm hoặc dòng điện khác đặt trong đó. Vì vậy phương án $C$ đúng; các phương án còn lại trái với định nghĩa, định luật hoặc điều kiện áp dụng.
}
\end{ex}

% ===== Câu 98 =====
% ID: L12C3B14-30-DS
% Mức: VD
\dangbt{Xác định chiều từ trường bằng kim nam châm và quy tắc nắm tay phải}
\begin{ex}
% ID: L12C3B14-30-DS
% Mức: VD
Xét xác định chiều từ trường bằng kim nam châm và quy tắc nắm tay phải (Tình huống 5). Hãy xác định tính đúng, sai của bốn phát biểu sau.
\choiceTF
{\True Đầu Bắc của kim nam châm nhỏ chỉ theo chiều của vectơ cảm ứng từ; quy tắc nắm tay phải xác định chiều đường sức quanh dòng điện.}
{Trong dạng bài này, kết quả không phụ thuộc dữ kiện và có thể bỏ qua điều kiện áp dụng.}
{Quy tắc bàn tay trái dùng để xác định chiều đường sức từ quanh dây dẫn thẳng.}
{\True Khi giải bài thuộc dạng này phải xác định đúng phương, chiều, điều kiện áp dụng và đơn vị của các đại lượng.}
\loigiai{
\begin{itemchoice}
\item (Đúng) Đầu Bắc của kim nam châm nhỏ chỉ theo chiều của vectơ cảm ứng từ; quy tắc nắm tay phải xác định chiều đường sức quanh dòng điện. (Vì): ầu Bắc của kim nam châm nhỏ chỉ theo chiều của vectơ cảm ứng từ; quy tắc nắm tay phải xác định chiều đường sức quanh dòng điện. b) Sai: Trong dạng bài này, kết quả không phụ thuộc dữ kiện và có thể bỏ qua điều kiện áp dụng. c) Sai: Quy tắc bàn tay trái dùng để xác định chiều đường sức từ quanh dây dẫn thẳng. d) Đúng: Khi giải bài thuộc dạng này phải xác định đúng phương, chiều, điều kiện áp dụng và đơn vị của các đại lượng. Kiến thức cốt lõi: Đầu Bắc của kim nam châm nhỏ chỉ theo chiều của vectơ cảm ứng từ; quy tắc nắm tay phải xác định chiều đường sức quanh dòng điện
\item (Sai) Trong dạng bài này, kết quả không phụ thuộc dữ kiện và có thể bỏ qua điều kiện áp dụng.
\item (Sai) Quy tắc bàn tay trái dùng để xác định chiều đường sức từ quanh dây dẫn thẳng.
\item (Đúng) Khi giải bài thuộc dạng này phải xác định đúng phương, chiều, điều kiện áp dụng và đơn vị của các đại lượng.
\end{itemchoice}
}
\end{ex}

% ===== Câu 99 =====
% ID: L12C3B14-30-TL
% Mức: VDC
\dangbt{Từ trường của dòng điện thẳng, dòng điện tròn và ống dây}
\begin{bt}
% ID: L12C3B14-30-TL
% Mức: VDC
So sánh nhận định đúng “Đường sức quanh dây dẫn thẳng là các đường tròn đồng tâm; từ trường trong lòng ống dây dài gần đều.” với nhận định sai “Từ trường trong lòng ống dây dài luôn bằng không.”; chỉ rõ căn cứ Vật lí.
\loigiai{
Cần nêu được: Đường sức quanh dây dẫn thẳng là các đường tròn đồng tâm; từ trường trong lòng ống dây dài gần đều. Khi vận dụng phải xác định đúng dữ kiện, phương chiều nếu có, điều kiện áp dụng, hệ đơn vị và kiểm tra tính hợp lí của kết quả. Nhận định “Từ trường trong lòng ống dây dài luôn bằng không.” sai vì trái với kiến thức nêu trên.
}
\end{bt}

% ===== Câu 100 =====
% ID: L12C3B14-30-TLN
% Mức: VDC
\begin{ex}
% ID: L12C3B14-30-TLN
% Mức: VDC
Trong bốn phát biểu sau về từ trường của dòng điện thẳng, dòng điện tròn và ống dây, có bao nhiêu phát biểu đúng? (Chỉ ghi một số nguyên.) (1) Kết quả cần có ý nghĩa vật lí. (2) Từ trường trong lòng ống dây dài luôn bằng không. (3) Phải dùng đúng định luật cho tình huống. (4) Đường sức quanh dây dẫn thẳng là các đường tròn đồng tâm; từ trường trong lòng ống dây dài gần đều.
\shortans{3}
\loigiai{
Đối chiếu từng phát biểu với kiến thức: Đường sức quanh dây dẫn thẳng là các đường tròn đồng tâm; từ trường trong lòng ống dây dài gần đều. Các phát biểu đúng có số lượng là 3.
}
\end{ex}

% ===== Câu 101 =====
% ID: L12C3B14-30-TN
% Mức: VD
\dangbt{Nhận biết từ trường, nguồn gây ra từ trường và tương tác từ}
\begin{ex}
% ID: L12C3B14-30-TN
% Mức: VD
Cơ sở Vật lí đúng dùng để giải nhận biết từ trường, nguồn gây ra từ trường và tương tác từ là gì?
\choice
{Có thể bỏ qua phương, chiều và đơn vị của các đại lượng trong mọi trường hợp.}
{Đại lượng đang xét không phụ thuộc các điều kiện vật lí nêu trong bài.}
{Mọi điện tích đứng yên đều tạo ra từ trường giống dòng điện.}
{\True Từ trường tồn tại xung quanh nam châm và dòng điện; nó gây lực từ lên nam châm hoặc dòng điện khác đặt trong đó.}
\loigiai{
Từ trường tồn tại xung quanh nam châm và dòng điện; nó gây lực từ lên nam châm hoặc dòng điện khác đặt trong đó. Vì vậy phương án $D$ đúng; các phương án còn lại trái với định nghĩa, định luật hoặc điều kiện áp dụng.
}
\end{ex}

% ===== Câu 102 =====
% ID: L12C3B14-31-DS
% Mức: TH
\begin{ex}
% ID: L12C3B14-31-DS
% Mức: TH
Xét nhận biết từ trường, nguồn gây ra từ trường và tương tác từ (Tình huống 1). Hãy xác định tính đúng, sai của bốn phát biểu sau.
\choiceTF
{\True Từ trường tồn tại xung quanh nam châm và dòng điện; nó gây lực từ lên nam châm hoặc dòng điện khác đặt trong đó.}
{Mọi điện tích đứng yên đều tạo ra từ trường giống dòng điện.}
{\True Khi giải bài thuộc dạng này phải xác định đúng phương, chiều, điều kiện áp dụng và đơn vị của các đại lượng.}
{Trong dạng bài này, kết quả không phụ thuộc dữ kiện và có thể bỏ qua điều kiện áp dụng.}
\loigiai{
\begin{itemchoice}
\item (Đúng) Từ trường tồn tại xung quanh nam châm và dòng điện; nó gây lực từ lên nam châm hoặc dòng điện khác đặt trong đó. (Vì): Từ trường tồn tại xung quanh nam châm và dòng điện; nó gây lực từ lên nam châm hoặc dòng điện khác đặt trong đó. b) Sai: Mọi điện tích đứng yên đều tạo ra từ trường giống dòng điện. c) Đúng: Khi giải bài thuộc dạng này phải xác định đúng phương, chiều, điều kiện áp dụng và đơn vị của các đại lượng. d) Sai: Trong dạng bài này, kết quả không phụ thuộc dữ kiện và có thể bỏ qua điều kiện áp dụng. Kiến thức cốt lõi: Từ trường tồn tại xung quanh nam châm và dòng điện; nó gây lực từ lên nam châm hoặc dòng điện khác đặt trong đó
\item (Sai) Mọi điện tích đứng yên đều tạo ra từ trường giống dòng điện.
\item (Đúng) Khi giải bài thuộc dạng này phải xác định đúng phương, chiều, điều kiện áp dụng và đơn vị của các đại lượng.
\item (Sai) Trong dạng bài này, kết quả không phụ thuộc dữ kiện và có thể bỏ qua điều kiện áp dụng.
\end{itemchoice}
}
\end{ex}

% ===== Câu 103 =====
% ID: L12C3B14-31-TL
% Mức: TH
\dangbt{Xác định chiều từ trường bằng kim nam châm và quy tắc nắm tay phải}
\begin{bt}
% ID: L12C3B14-31-TL
% Mức: TH
Trình bày kiến thức cốt lõi của xác định chiều từ trường bằng kim nam châm và quy tắc nắm tay phải và nêu điều kiện áp dụng.
\loigiai{
Cần nêu được: Đầu Bắc của kim nam châm nhỏ chỉ theo chiều của vectơ cảm ứng từ; quy tắc nắm tay phải xác định chiều đường sức quanh dòng điện. Khi vận dụng phải xác định đúng dữ kiện, phương chiều nếu có, điều kiện áp dụng, hệ đơn vị và kiểm tra tính hợp lí của kết quả. Nhận định “Quy tắc bàn tay trái dùng để xác định chiều đường sức từ quanh dây dẫn thẳng.” sai vì trái với kiến thức nêu trên.
}
\end{bt}

% ===== Câu 104 =====
% ID: L12C3B14-31-TLN
% Mức: TH
\begin{ex}
% ID: L12C3B14-31-TLN
% Mức: TH
Trong bốn phát biểu sau về xác định chiều từ trường bằng kim nam châm và quy tắc nắm tay phải, có bao nhiêu phát biểu đúng? (Chỉ ghi một số nguyên.) (1) Đầu Bắc của kim nam châm nhỏ chỉ theo chiều của vectơ cảm ứng từ; quy tắc nắm tay phải xác định chiều đường sức quanh dòng điện. (2) Quy tắc bàn tay trái dùng để xác định chiều đường sức từ quanh dây dẫn thẳng. (3) Cần kiểm tra điều kiện áp dụng trước khi tính toán. (4) Có thể bỏ qua đơn vị của kết quả.
\shortans{2}
\loigiai{
Đối chiếu từng phát biểu với kiến thức: Đầu Bắc của kim nam châm nhỏ chỉ theo chiều của vectơ cảm ứng từ; quy tắc nắm tay phải xác định chiều đường sức quanh dòng điện. Các phát biểu đúng có số lượng là 2.
}
\end{ex}

% ===== Câu 105 =====
% ID: L12C3B14-31-TN
% Mức: VD
\dangbt{Nhận biết từ trường, nguồn gây ra từ trường và tương tác từ}
\begin{ex}
% ID: L12C3B14-31-TN
% Mức: VD
Trong một bài thực hành về nhận biết từ trường, nguồn gây ra từ trường và tương tác từ, nhận xét nào đúng?
\choice
{\True Từ trường tồn tại xung quanh nam châm và dòng điện; nó gây lực từ lên nam châm hoặc dòng điện khác đặt trong đó.}
{Mọi điện tích đứng yên đều tạo ra từ trường giống dòng điện.}
{Đại lượng đang xét không phụ thuộc các điều kiện vật lí nêu trong bài.}
{Có thể bỏ qua phương, chiều và đơn vị của các đại lượng trong mọi trường hợp.}
\loigiai{
Từ trường tồn tại xung quanh nam châm và dòng điện; nó gây lực từ lên nam châm hoặc dòng điện khác đặt trong đó. Vì vậy phương án $A$ đúng; các phương án còn lại trái với định nghĩa, định luật hoặc điều kiện áp dụng.
}
\end{ex}

% ===== Câu 106 =====
% ID: L12C3B14-32-DS
% Mức: TH
\begin{ex}
% ID: L12C3B14-32-DS
% Mức: TH
Xét nhận biết từ trường, nguồn gây ra từ trường và tương tác từ (Tình huống 2). Hãy xác định tính đúng, sai của bốn phát biểu sau.
\choiceTF
{Mọi điện tích đứng yên đều tạo ra từ trường giống dòng điện.}
{\True Từ trường tồn tại xung quanh nam châm và dòng điện; nó gây lực từ lên nam châm hoặc dòng điện khác đặt trong đó.}
{Trong dạng bài này, kết quả không phụ thuộc dữ kiện và có thể bỏ qua điều kiện áp dụng.}
{\True Khi giải bài thuộc dạng này phải xác định đúng phương, chiều, điều kiện áp dụng và đơn vị của các đại lượng.}
\loigiai{
\begin{itemchoice}
\item (Sai) Mọi điện tích đứng yên đều tạo ra từ trường giống dòng điện. (Vì): Mọi điện tích đứng yên đều tạo ra từ trường giống dòng điện. b) Đúng: Từ trường tồn tại xung quanh nam châm và dòng điện; nó gây lực từ lên nam châm hoặc dòng điện khác đặt trong đó. c) Sai: Trong dạng bài này, kết quả không phụ thuộc dữ kiện và có thể bỏ qua điều kiện áp dụng. d) Đúng: Khi giải bài thuộc dạng này phải xác định đúng phương, chiều, điều kiện áp dụng và đơn vị của các đại lượng. Kiến thức cốt lõi: Từ trường tồn tại xung quanh nam châm và dòng điện; nó gây lực từ lên nam châm hoặc dòng điện khác đặt trong đó
\item (Đúng) Từ trường tồn tại xung quanh nam châm và dòng điện; nó gây lực từ lên nam châm hoặc dòng điện khác đặt trong đó.
\item (Sai) Trong dạng bài này, kết quả không phụ thuộc dữ kiện và có thể bỏ qua điều kiện áp dụng.
\item (Đúng) Khi giải bài thuộc dạng này phải xác định đúng phương, chiều, điều kiện áp dụng và đơn vị của các đại lượng.
\end{itemchoice}
}
\end{ex}

% ===== Câu 107 =====
% ID: L12C3B14-32-TL
% Mức: TH
\dangbt{Xác định chiều từ trường bằng kim nam châm và quy tắc nắm tay phải}
\begin{bt}
% ID: L12C3B14-32-TL
% Mức: TH
Giải thích vì sao nhận định sau đúng: “Đầu Bắc của kim nam châm nhỏ chỉ theo chiều của vectơ cảm ứng từ; quy tắc nắm tay phải xác định chiều đường sức quanh dòng điện.”
\loigiai{
Cần nêu được: Đầu Bắc của kim nam châm nhỏ chỉ theo chiều của vectơ cảm ứng từ; quy tắc nắm tay phải xác định chiều đường sức quanh dòng điện. Khi vận dụng phải xác định đúng dữ kiện, phương chiều nếu có, điều kiện áp dụng, hệ đơn vị và kiểm tra tính hợp lí của kết quả. Nhận định “Quy tắc bàn tay trái dùng để xác định chiều đường sức từ quanh dây dẫn thẳng.” sai vì trái với kiến thức nêu trên.
}
\end{bt}

% ===== Câu 108 =====
% ID: L12C3B14-32-TLN
% Mức: TH
\begin{ex}
% ID: L12C3B14-32-TLN
% Mức: TH
Trong bốn phát biểu sau về xác định chiều từ trường bằng kim nam châm và quy tắc nắm tay phải, có bao nhiêu phát biểu đúng? (Chỉ ghi một số nguyên.) (1) Quy tắc bàn tay trái dùng để xác định chiều đường sức từ quanh dây dẫn thẳng. (2) Đầu Bắc của kim nam châm nhỏ chỉ theo chiều của vectơ cảm ứng từ; quy tắc nắm tay phải xác định chiều đường sức quanh dòng điện. (3) Kết quả phải phù hợp ý nghĩa vật lí. (4) Mọi đại lượng đều là vô hướng.
\shortans{1}
\loigiai{
Đối chiếu từng phát biểu với kiến thức: Đầu Bắc của kim nam châm nhỏ chỉ theo chiều của vectơ cảm ứng từ; quy tắc nắm tay phải xác định chiều đường sức quanh dòng điện. Các phát biểu đúng có số lượng là 1.
}
\end{ex}

% ===== Câu 109 =====
% ID: L12C3B14-32-TN
% Mức: VD
\dangbt{Nhận biết từ trường, nguồn gây ra từ trường và tương tác từ}
\begin{ex}
% ID: L12C3B14-32-TN
% Mức: VD
Kết luận đúng khi vận dụng nhận biết từ trường, nguồn gây ra từ trường và tương tác từ là
\choice
{Mọi điện tích đứng yên đều tạo ra từ trường giống dòng điện.}
{\True Từ trường tồn tại xung quanh nam châm và dòng điện; nó gây lực từ lên nam châm hoặc dòng điện khác đặt trong đó.}
{Có thể bỏ qua phương, chiều và đơn vị của các đại lượng trong mọi trường hợp.}
{Đại lượng đang xét không phụ thuộc các điều kiện vật lí nêu trong bài.}
\loigiai{
Từ trường tồn tại xung quanh nam châm và dòng điện; nó gây lực từ lên nam châm hoặc dòng điện khác đặt trong đó. Vì vậy phương án $B$ đúng; các phương án còn lại trái với định nghĩa, định luật hoặc điều kiện áp dụng.
}
\end{ex}

% ===== Câu 110 =====
% ID: L12C3B14-33-DS
% Mức: VD
\begin{ex}
% ID: L12C3B14-33-DS
% Mức: VD
Xét nhận biết từ trường, nguồn gây ra từ trường và tương tác từ (Tình huống 3). Hãy xác định tính đúng, sai của bốn phát biểu sau.
\choiceTF
{\True Khi giải bài thuộc dạng này phải xác định đúng phương, chiều, điều kiện áp dụng và đơn vị của các đại lượng.}
{Trong dạng bài này, kết quả không phụ thuộc dữ kiện và có thể bỏ qua điều kiện áp dụng.}
{\True Từ trường tồn tại xung quanh nam châm và dòng điện; nó gây lực từ lên nam châm hoặc dòng điện khác đặt trong đó.}
{Mọi điện tích đứng yên đều tạo ra từ trường giống dòng điện.}
\loigiai{
\begin{itemchoice}
\item (Đúng) Khi giải bài thuộc dạng này phải xác định đúng phương, chiều, điều kiện áp dụng và đơn vị của các đại lượng. (Vì): Khi giải bài thuộc dạng này phải xác định đúng phương, chiều, điều kiện áp dụng và đơn vị của các đại lượng. b) Sai: Trong dạng bài này, kết quả không phụ thuộc dữ kiện và có thể bỏ qua điều kiện áp dụng. c) Đúng: Từ trường tồn tại xung quanh nam châm và dòng điện; nó gây lực từ lên nam châm hoặc dòng điện khác đặt trong đó. d) Sai: Mọi điện tích đứng yên đều tạo ra từ trường giống dòng điện. Kiến thức cốt lõi: Từ trường tồn tại xung quanh nam châm và dòng điện; nó gây lực từ lên nam châm hoặc dòng điện khác đặt trong đó
\item (Sai) Trong dạng bài này, kết quả không phụ thuộc dữ kiện và có thể bỏ qua điều kiện áp dụng.
\item (Đúng) Từ trường tồn tại xung quanh nam châm và dòng điện; nó gây lực từ lên nam châm hoặc dòng điện khác đặt trong đó.
\item (Sai) Mọi điện tích đứng yên đều tạo ra từ trường giống dòng điện.
\end{itemchoice}
}
\end{ex}

% ===== Câu 111 =====
% ID: L12C3B14-33-TL
% Mức: VD
\dangbt{Xác định chiều từ trường bằng kim nam châm và quy tắc nắm tay phải}
\begin{bt}
% ID: L12C3B14-33-TL
% Mức: VD
Phân tích sai lầm trong nhận định: “Quy tắc bàn tay trái dùng để xác định chiều đường sức từ quanh dây dẫn thẳng.”
\loigiai{
Cần nêu được: Đầu Bắc của kim nam châm nhỏ chỉ theo chiều của vectơ cảm ứng từ; quy tắc nắm tay phải xác định chiều đường sức quanh dòng điện. Khi vận dụng phải xác định đúng dữ kiện, phương chiều nếu có, điều kiện áp dụng, hệ đơn vị và kiểm tra tính hợp lí của kết quả. Nhận định “Quy tắc bàn tay trái dùng để xác định chiều đường sức từ quanh dây dẫn thẳng.” sai vì trái với kiến thức nêu trên.
}
\end{bt}

% ===== Câu 112 =====
% ID: L12C3B14-33-TLN
% Mức: TH
\begin{ex}
% ID: L12C3B14-33-TLN
% Mức: TH
Trong bốn phát biểu sau về xác định chiều từ trường bằng kim nam châm và quy tắc nắm tay phải, có bao nhiêu phát biểu đúng? (Chỉ ghi một số nguyên.) (1) Cần đổi các đại lượng về hệ đơn vị thống nhất. (2) Đầu Bắc của kim nam châm nhỏ chỉ theo chiều của vectơ cảm ứng từ; quy tắc nắm tay phải xác định chiều đường sức quanh dòng điện. (3) Quy tắc bàn tay trái dùng để xác định chiều đường sức từ quanh dây dẫn thẳng. (4) Không cần xét chiều của đại lượng vectơ.
\shortans{2}
\loigiai{
Đối chiếu từng phát biểu với kiến thức: Đầu Bắc của kim nam châm nhỏ chỉ theo chiều của vectơ cảm ứng từ; quy tắc nắm tay phải xác định chiều đường sức quanh dòng điện. Các phát biểu đúng có số lượng là 2.
}
\end{ex}

% ===== Câu 113 =====
% ID: L12C3B14-33-TN
% Mức: NB
\dangbt{So sánh, tổng hợp và biểu diễn các từ trường đơn giản}
\begin{ex}
% ID: L12C3B14-33-TN
% Mức: NB
Phát biểu nào sau đây đúng về so sánh, tổng hợp và biểu diễn các từ trường đơn giản?
\choice
{\True Vectơ cảm ứng từ tổng hợp tại một điểm bằng tổng vectơ các cảm ứng từ thành phần.}
{Độ lớn cảm ứng từ tổng hợp luôn bằng tổng độ lớn các cảm ứng từ thành phần.}
{Đại lượng đang xét không phụ thuộc các điều kiện vật lí nêu trong bài.}
{Có thể bỏ qua phương, chiều và đơn vị của các đại lượng trong mọi trường hợp.}
\loigiai{
Vectơ cảm ứng từ tổng hợp tại một điểm bằng tổng vectơ các cảm ứng từ thành phần. Vì vậy phương án $A$ đúng; các phương án còn lại trái với định nghĩa, định luật hoặc điều kiện áp dụng.
}
\end{ex}

% ===== Câu 114 =====
% ID: L12C3B14-34-DS
% Mức: VD
\dangbt{Nhận biết từ trường, nguồn gây ra từ trường và tương tác từ}
\begin{ex}
% ID: L12C3B14-34-DS
% Mức: VD
Xét nhận biết từ trường, nguồn gây ra từ trường và tương tác từ (Tình huống 4). Hãy xác định tính đúng, sai của bốn phát biểu sau.
\choiceTF
{Trong dạng bài này, kết quả không phụ thuộc dữ kiện và có thể bỏ qua điều kiện áp dụng.}
{\True Khi giải bài thuộc dạng này phải xác định đúng phương, chiều, điều kiện áp dụng và đơn vị của các đại lượng.}
{Mọi điện tích đứng yên đều tạo ra từ trường giống dòng điện.}
{\True Từ trường tồn tại xung quanh nam châm và dòng điện; nó gây lực từ lên nam châm hoặc dòng điện khác đặt trong đó.}
\loigiai{
\begin{itemchoice}
\item (Sai) Trong dạng bài này, kết quả không phụ thuộc dữ kiện và có thể bỏ qua điều kiện áp dụng. (Vì): Trong dạng bài này, kết quả không phụ thuộc dữ kiện và có thể bỏ qua điều kiện áp dụng. b) Đúng: Khi giải bài thuộc dạng này phải xác định đúng phương, chiều, điều kiện áp dụng và đơn vị của các đại lượng. c) Sai: Mọi điện tích đứng yên đều tạo ra từ trường giống dòng điện. d) Đúng: Từ trường tồn tại xung quanh nam châm và dòng điện; nó gây lực từ lên nam châm hoặc dòng điện khác đặt trong đó. Kiến thức cốt lõi: Từ trường tồn tại xung quanh nam châm và dòng điện; nó gây lực từ lên nam châm hoặc dòng điện khác đặt trong đó
\item (Đúng) Khi giải bài thuộc dạng này phải xác định đúng phương, chiều, điều kiện áp dụng và đơn vị của các đại lượng.
\item (Sai) Mọi điện tích đứng yên đều tạo ra từ trường giống dòng điện.
\item (Đúng) Từ trường tồn tại xung quanh nam châm và dòng điện; nó gây lực từ lên nam châm hoặc dòng điện khác đặt trong đó.
\end{itemchoice}
}
\end{ex}

% ===== Câu 115 =====
% ID: L12C3B14-34-TL
% Mức: VD
\dangbt{Xác định chiều từ trường bằng kim nam châm và quy tắc nắm tay phải}
\begin{bt}
% ID: L12C3B14-34-TL
% Mức: VD
Đề xuất các bước giải một bài toán thuộc dạng xác định chiều từ trường bằng kim nam châm và quy tắc nắm tay phải.
\loigiai{
Cần nêu được: Đầu Bắc của kim nam châm nhỏ chỉ theo chiều của vectơ cảm ứng từ; quy tắc nắm tay phải xác định chiều đường sức quanh dòng điện. Khi vận dụng phải xác định đúng dữ kiện, phương chiều nếu có, điều kiện áp dụng, hệ đơn vị và kiểm tra tính hợp lí của kết quả. Nhận định “Quy tắc bàn tay trái dùng để xác định chiều đường sức từ quanh dây dẫn thẳng.” sai vì trái với kiến thức nêu trên.
}
\end{bt}

% ===== Câu 116 =====
% ID: L12C3B14-34-TLN
% Mức: VD
\begin{ex}
% ID: L12C3B14-34-TLN
% Mức: VD
Trong bốn phát biểu sau về xác định chiều từ trường bằng kim nam châm và quy tắc nắm tay phải, có bao nhiêu phát biểu đúng? (Chỉ ghi một số nguyên.) (1) Đầu Bắc của kim nam châm nhỏ chỉ theo chiều của vectơ cảm ứng từ; quy tắc nắm tay phải xác định chiều đường sức quanh dòng điện. (2) Cần đối chiếu kết quả với điều kiện bài toán. (3) Quy tắc bàn tay trái dùng để xác định chiều đường sức từ quanh dây dẫn thẳng. (4) Mọi trường hợp đều cho cùng một kết quả.
\shortans{2}
\loigiai{
Đối chiếu từng phát biểu với kiến thức: Đầu Bắc của kim nam châm nhỏ chỉ theo chiều của vectơ cảm ứng từ; quy tắc nắm tay phải xác định chiều đường sức quanh dòng điện. Các phát biểu đúng có số lượng là 2.
}
\end{ex}

% ===== Câu 117 =====
% ID: L12C3B14-34-TN
% Mức: NB
\dangbt{So sánh, tổng hợp và biểu diễn các từ trường đơn giản}
\begin{ex}
% ID: L12C3B14-34-TN
% Mức: NB
Trong so sánh, tổng hợp và biểu diễn các từ trường đơn giản, nhận định nào phù hợp với kiến thức Vật lí?
\choice
{Độ lớn cảm ứng từ tổng hợp luôn bằng tổng độ lớn các cảm ứng từ thành phần.}
{\True Vectơ cảm ứng từ tổng hợp tại một điểm bằng tổng vectơ các cảm ứng từ thành phần.}
{Có thể bỏ qua phương, chiều và đơn vị của các đại lượng trong mọi trường hợp.}
{Đại lượng đang xét không phụ thuộc các điều kiện vật lí nêu trong bài.}
\loigiai{
Vectơ cảm ứng từ tổng hợp tại một điểm bằng tổng vectơ các cảm ứng từ thành phần. Vì vậy phương án $B$ đúng; các phương án còn lại trái với định nghĩa, định luật hoặc điều kiện áp dụng.
}
\end{ex}

% ===== Câu 118 =====
% ID: L12C3B14-35-DS
% Mức: VD
\dangbt{Nhận biết từ trường, nguồn gây ra từ trường và tương tác từ}
\begin{ex}
% ID: L12C3B14-35-DS
% Mức: VD
Xét nhận biết từ trường, nguồn gây ra từ trường và tương tác từ (Tình huống 5). Hãy xác định tính đúng, sai của bốn phát biểu sau.
\choiceTF
{\True Từ trường tồn tại xung quanh nam châm và dòng điện; nó gây lực từ lên nam châm hoặc dòng điện khác đặt trong đó.}
{Trong dạng bài này, kết quả không phụ thuộc dữ kiện và có thể bỏ qua điều kiện áp dụng.}
{Mọi điện tích đứng yên đều tạo ra từ trường giống dòng điện.}
{\True Khi giải bài thuộc dạng này phải xác định đúng phương, chiều, điều kiện áp dụng và đơn vị của các đại lượng.}
\loigiai{
\begin{itemchoice}
\item (Đúng) Từ trường tồn tại xung quanh nam châm và dòng điện; nó gây lực từ lên nam châm hoặc dòng điện khác đặt trong đó. (Vì): Từ trường tồn tại xung quanh nam châm và dòng điện; nó gây lực từ lên nam châm hoặc dòng điện khác đặt trong đó. b) Sai: Trong dạng bài này, kết quả không phụ thuộc dữ kiện và có thể bỏ qua điều kiện áp dụng. c) Sai: Mọi điện tích đứng yên đều tạo ra từ trường giống dòng điện. d) Đúng: Khi giải bài thuộc dạng này phải xác định đúng phương, chiều, điều kiện áp dụng và đơn vị của các đại lượng. Kiến thức cốt lõi: Từ trường tồn tại xung quanh nam châm và dòng điện; nó gây lực từ lên nam châm hoặc dòng điện khác đặt trong đó
\item (Sai) Trong dạng bài này, kết quả không phụ thuộc dữ kiện và có thể bỏ qua điều kiện áp dụng.
\item (Sai) Mọi điện tích đứng yên đều tạo ra từ trường giống dòng điện.
\item (Đúng) Khi giải bài thuộc dạng này phải xác định đúng phương, chiều, điều kiện áp dụng và đơn vị của các đại lượng.
\end{itemchoice}
}
\end{ex}

% ===== Câu 119 =====
% ID: L12C3B14-35-TL
% Mức: VD
\dangbt{Xác định chiều từ trường bằng kim nam châm và quy tắc nắm tay phải}
\begin{bt}
% ID: L12C3B14-35-TL
% Mức: VD
Nêu một tình huống thực tiễn liên quan đến xác định chiều từ trường bằng kim nam châm và quy tắc nắm tay phải và giải thích bằng kiến thức Vật lí.
\loigiai{
Cần nêu được: Đầu Bắc của kim nam châm nhỏ chỉ theo chiều của vectơ cảm ứng từ; quy tắc nắm tay phải xác định chiều đường sức quanh dòng điện. Khi vận dụng phải xác định đúng dữ kiện, phương chiều nếu có, điều kiện áp dụng, hệ đơn vị và kiểm tra tính hợp lí của kết quả. Nhận định “Quy tắc bàn tay trái dùng để xác định chiều đường sức từ quanh dây dẫn thẳng.” sai vì trái với kiến thức nêu trên.
}
\end{bt}

% ===== Câu 120 =====
% ID: L12C3B14-35-TLN
% Mức: VD
\begin{ex}
% ID: L12C3B14-35-TLN
% Mức: VD
Trong bốn phát biểu sau về xác định chiều từ trường bằng kim nam châm và quy tắc nắm tay phải, có bao nhiêu phát biểu đúng? (Chỉ ghi một số nguyên.) (1) Quy tắc bàn tay trái dùng để xác định chiều đường sức từ quanh dây dẫn thẳng. (2) Cần đọc đúng dữ kiện và giả thiết. (3) Đầu Bắc của kim nam châm nhỏ chỉ theo chiều của vectơ cảm ứng từ; quy tắc nắm tay phải xác định chiều đường sức quanh dòng điện. (4) Có thể thay số trước khi xác định công thức.
\shortans{2}
\loigiai{
Đối chiếu từng phát biểu với kiến thức: Đầu Bắc của kim nam châm nhỏ chỉ theo chiều của vectơ cảm ứng từ; quy tắc nắm tay phải xác định chiều đường sức quanh dòng điện. Các phát biểu đúng có số lượng là 2.
}
\end{ex}

% ===== Câu 121 =====
% ID: L12C3B14-35-TN
% Mức: NB
\dangbt{So sánh, tổng hợp và biểu diễn các từ trường đơn giản}
\begin{ex}
% ID: L12C3B14-35-TN
% Mức: NB
Tình huống 3: chọn kết luận chính xác về so sánh, tổng hợp và biểu diễn các từ trường đơn giản.
\choice
{Đại lượng đang xét không phụ thuộc các điều kiện vật lí nêu trong bài.}
{Có thể bỏ qua phương, chiều và đơn vị của các đại lượng trong mọi trường hợp.}
{\True Vectơ cảm ứng từ tổng hợp tại một điểm bằng tổng vectơ các cảm ứng từ thành phần.}
{Độ lớn cảm ứng từ tổng hợp luôn bằng tổng độ lớn các cảm ứng từ thành phần.}
\loigiai{
Vectơ cảm ứng từ tổng hợp tại một điểm bằng tổng vectơ các cảm ứng từ thành phần. Vì vậy phương án $C$ đúng; các phương án còn lại trái với định nghĩa, định luật hoặc điều kiện áp dụng.
}
\end{ex}

% ===== Câu 122 =====
% ID: L12C3B14-36-DS
% Mức: TH
\dangbt{Mô tả đường sức từ và xác định chiều đường sức từ}
\begin{ex}
% ID: L12C3B14-36-DS
% Mức: TH
Xét mô tả đường sức từ và xác định chiều đường sức từ (Tình huống 1). Hãy xác định tính đúng, sai của bốn phát biểu sau.
\choiceTF
{\True Tiếp tuyến với đường sức từ tại một điểm cho phương của vectơ cảm ứng từ; bên ngoài nam châm, đường sức đi từ cực Bắc đến cực Nam.}
{Các đường sức từ có thể cắt nhau tại một điểm.}
{\True Khi giải bài thuộc dạng này phải xác định đúng phương, chiều, điều kiện áp dụng và đơn vị của các đại lượng.}
{Trong dạng bài này, kết quả không phụ thuộc dữ kiện và có thể bỏ qua điều kiện áp dụng.}
\loigiai{
\begin{itemchoice}
\item (Đúng) Tiếp tuyến với đường sức từ tại một điểm cho phương của vectơ cảm ứng từ; bên ngoài nam châm, đường sức đi từ cực Bắc đến cực Nam. (Vì): Tiếp tuyến với đường sức từ tại một điểm cho phương của vectơ cảm ứng từ; bên ngoài nam châm, đường sức đi từ cực Bắc đến cực Nam. b) Sai: Các đường sức từ có thể cắt nhau tại một điểm. c) Đúng: Khi giải bài thuộc dạng này phải xác định đúng phương, chiều, điều kiện áp dụng và đơn vị của các đại lượng. d) Sai: Trong dạng bài này, kết quả không phụ thuộc dữ kiện và có thể bỏ qua điều kiện áp dụng. Kiến thức cốt lõi: Tiếp tuyến với đường sức từ tại một điểm cho phương của vectơ cảm ứng từ; bên ngoài nam châm, đường sức đi từ cực Bắc đến cực Nam
\item (Sai) Các đường sức từ có thể cắt nhau tại một điểm.
\item (Đúng) Khi giải bài thuộc dạng này phải xác định đúng phương, chiều, điều kiện áp dụng và đơn vị của các đại lượng.
\item (Sai) Trong dạng bài này, kết quả không phụ thuộc dữ kiện và có thể bỏ qua điều kiện áp dụng.
\end{itemchoice}
}
\end{ex}

% ===== Câu 123 =====
% ID: L12C3B14-36-TL
% Mức: VDC
\dangbt{Xác định chiều từ trường bằng kim nam châm và quy tắc nắm tay phải}
\begin{bt}
% ID: L12C3B14-36-TL
% Mức: VDC
So sánh nhận định đúng “Đầu Bắc của kim nam châm nhỏ chỉ theo chiều của vectơ cảm ứng từ; quy tắc nắm tay phải xác định chiều đường sức quanh dòng điện.” với nhận định sai “Quy tắc bàn tay trái dùng để xác định chiều đường sức từ quanh dây dẫn thẳng.”; chỉ rõ căn cứ Vật lí.
\loigiai{
Cần nêu được: Đầu Bắc của kim nam châm nhỏ chỉ theo chiều của vectơ cảm ứng từ; quy tắc nắm tay phải xác định chiều đường sức quanh dòng điện. Khi vận dụng phải xác định đúng dữ kiện, phương chiều nếu có, điều kiện áp dụng, hệ đơn vị và kiểm tra tính hợp lí của kết quả. Nhận định “Quy tắc bàn tay trái dùng để xác định chiều đường sức từ quanh dây dẫn thẳng.” sai vì trái với kiến thức nêu trên.
}
\end{bt}

% ===== Câu 124 =====
% ID: L12C3B14-36-TLN
% Mức: VDC
\begin{ex}
% ID: L12C3B14-36-TLN
% Mức: VDC
Trong bốn phát biểu sau về xác định chiều từ trường bằng kim nam châm và quy tắc nắm tay phải, có bao nhiêu phát biểu đúng? (Chỉ ghi một số nguyên.) (1) Kết quả cần có ý nghĩa vật lí. (2) Quy tắc bàn tay trái dùng để xác định chiều đường sức từ quanh dây dẫn thẳng. (3) Phải dùng đúng định luật cho tình huống. (4) Đầu Bắc của kim nam châm nhỏ chỉ theo chiều của vectơ cảm ứng từ; quy tắc nắm tay phải xác định chiều đường sức quanh dòng điện.
\shortans{3}
\loigiai{
Đối chiếu từng phát biểu với kiến thức: Đầu Bắc của kim nam châm nhỏ chỉ theo chiều của vectơ cảm ứng từ; quy tắc nắm tay phải xác định chiều đường sức quanh dòng điện. Các phát biểu đúng có số lượng là 3.
}
\end{ex}

% ===== Câu 125 =====
% ID: L12C3B14-36-TN
% Mức: TH
\dangbt{So sánh, tổng hợp và biểu diễn các từ trường đơn giản}
\begin{ex}
% ID: L12C3B14-36-TN
% Mức: TH
Nội dung nào mô tả đúng nhất so sánh, tổng hợp và biểu diễn các từ trường đơn giản?
\choice
{Có thể bỏ qua phương, chiều và đơn vị của các đại lượng trong mọi trường hợp.}
{Đại lượng đang xét không phụ thuộc các điều kiện vật lí nêu trong bài.}
{Độ lớn cảm ứng từ tổng hợp luôn bằng tổng độ lớn các cảm ứng từ thành phần.}
{\True Vectơ cảm ứng từ tổng hợp tại một điểm bằng tổng vectơ các cảm ứng từ thành phần.}
\loigiai{
Vectơ cảm ứng từ tổng hợp tại một điểm bằng tổng vectơ các cảm ứng từ thành phần. Vì vậy phương án $D$ đúng; các phương án còn lại trái với định nghĩa, định luật hoặc điều kiện áp dụng.
}
\end{ex}

% ===== Câu 126 =====
% ID: L12C3B14-37-DS
% Mức: TH
\dangbt{Mô tả đường sức từ và xác định chiều đường sức từ}
\begin{ex}
% ID: L12C3B14-37-DS
% Mức: TH
Xét mô tả đường sức từ và xác định chiều đường sức từ (Tình huống 2). Hãy xác định tính đúng, sai của bốn phát biểu sau.
\choiceTF
{Các đường sức từ có thể cắt nhau tại một điểm.}
{\True Tiếp tuyến với đường sức từ tại một điểm cho phương của vectơ cảm ứng từ; bên ngoài nam châm, đường sức đi từ cực Bắc đến cực Nam.}
{Trong dạng bài này, kết quả không phụ thuộc dữ kiện và có thể bỏ qua điều kiện áp dụng.}
{\True Khi giải bài thuộc dạng này phải xác định đúng phương, chiều, điều kiện áp dụng và đơn vị của các đại lượng.}
\loigiai{
\begin{itemchoice}
\item (Sai) Các đường sức từ có thể cắt nhau tại một điểm. (Vì): Các đường sức từ có thể cắt nhau tại một điểm. b) Đúng: Tiếp tuyến với đường sức từ tại một điểm cho phương của vectơ cảm ứng từ; bên ngoài nam châm, đường sức đi từ cực Bắc đến cực Nam. c) Sai: Trong dạng bài này, kết quả không phụ thuộc dữ kiện và có thể bỏ qua điều kiện áp dụng. d) Đúng: Khi giải bài thuộc dạng này phải xác định đúng phương, chiều, điều kiện áp dụng và đơn vị của các đại lượng. Kiến thức cốt lõi: Tiếp tuyến với đường sức từ tại một điểm cho phương của vectơ cảm ứng từ; bên ngoài nam châm, đường sức đi từ cực Bắc đến cực Nam
\item (Đúng) Tiếp tuyến với đường sức từ tại một điểm cho phương của vectơ cảm ứng từ; bên ngoài nam châm, đường sức đi từ cực Bắc đến cực Nam.
\item (Sai) Trong dạng bài này, kết quả không phụ thuộc dữ kiện và có thể bỏ qua điều kiện áp dụng.
\item (Đúng) Khi giải bài thuộc dạng này phải xác định đúng phương, chiều, điều kiện áp dụng và đơn vị của các đại lượng.
\end{itemchoice}
}
\end{ex}

% ===== Câu 127 =====
% ID: L12C3B14-37-TL
% Mức: TH
\dangbt{Nhận biết từ trường, nguồn gây ra từ trường và tương tác từ}
\begin{bt}
% ID: L12C3B14-37-TL
% Mức: TH
Trình bày kiến thức cốt lõi của nhận biết từ trường, nguồn gây ra từ trường và tương tác từ và nêu điều kiện áp dụng.
\loigiai{
Cần nêu được: Từ trường tồn tại xung quanh nam châm và dòng điện; nó gây lực từ lên nam châm hoặc dòng điện khác đặt trong đó. Khi vận dụng phải xác định đúng dữ kiện, phương chiều nếu có, điều kiện áp dụng, hệ đơn vị và kiểm tra tính hợp lí của kết quả. Nhận định “Mọi điện tích đứng yên đều tạo ra từ trường giống dòng điện.” sai vì trái với kiến thức nêu trên.
}
\end{bt}

% ===== Câu 128 =====
% ID: L12C3B14-37-TLN
% Mức: TH
\begin{ex}
% ID: L12C3B14-37-TLN
% Mức: TH
Trong bốn phát biểu sau về nhận biết từ trường, nguồn gây ra từ trường và tương tác từ, có bao nhiêu phát biểu đúng? (Chỉ ghi một số nguyên.) (1) Từ trường tồn tại xung quanh nam châm và dòng điện; nó gây lực từ lên nam châm hoặc dòng điện khác đặt trong đó. (2) Mọi điện tích đứng yên đều tạo ra từ trường giống dòng điện. (3) Cần kiểm tra điều kiện áp dụng trước khi tính toán. (4) Có thể bỏ qua đơn vị của kết quả.
\shortans{2}
\loigiai{
Đối chiếu từng phát biểu với kiến thức: Từ trường tồn tại xung quanh nam châm và dòng điện; nó gây lực từ lên nam châm hoặc dòng điện khác đặt trong đó. Các phát biểu đúng có số lượng là 2.
}
\end{ex}

% ===== Câu 129 =====
% ID: L12C3B14-37-TN
% Mức: TH
\dangbt{So sánh, tổng hợp và biểu diễn các từ trường đơn giản}
\begin{ex}
% ID: L12C3B14-37-TN
% Mức: TH
Khi phân tích so sánh, tổng hợp và biểu diễn các từ trường đơn giản, phát biểu nào đúng?
\choice
{\True Vectơ cảm ứng từ tổng hợp tại một điểm bằng tổng vectơ các cảm ứng từ thành phần.}
{Độ lớn cảm ứng từ tổng hợp luôn bằng tổng độ lớn các cảm ứng từ thành phần.}
{Đại lượng đang xét không phụ thuộc các điều kiện vật lí nêu trong bài.}
{Có thể bỏ qua phương, chiều và đơn vị của các đại lượng trong mọi trường hợp.}
\loigiai{
Vectơ cảm ứng từ tổng hợp tại một điểm bằng tổng vectơ các cảm ứng từ thành phần. Vì vậy phương án $A$ đúng; các phương án còn lại trái với định nghĩa, định luật hoặc điều kiện áp dụng.
}
\end{ex}

% ===== Câu 130 =====
% ID: L12C3B14-38-DS
% Mức: VD
\dangbt{Mô tả đường sức từ và xác định chiều đường sức từ}
\begin{ex}
% ID: L12C3B14-38-DS
% Mức: VD
Xét mô tả đường sức từ và xác định chiều đường sức từ (Tình huống 3). Hãy xác định tính đúng, sai của bốn phát biểu sau.
\choiceTF
{\True Khi giải bài thuộc dạng này phải xác định đúng phương, chiều, điều kiện áp dụng và đơn vị của các đại lượng.}
{Trong dạng bài này, kết quả không phụ thuộc dữ kiện và có thể bỏ qua điều kiện áp dụng.}
{\True Tiếp tuyến với đường sức từ tại một điểm cho phương của vectơ cảm ứng từ; bên ngoài nam châm, đường sức đi từ cực Bắc đến cực Nam.}
{Các đường sức từ có thể cắt nhau tại một điểm.}
\loigiai{
\begin{itemchoice}
\item (Đúng) Khi giải bài thuộc dạng này phải xác định đúng phương, chiều, điều kiện áp dụng và đơn vị của các đại lượng. (Vì): Khi giải bài thuộc dạng này phải xác định đúng phương, chiều, điều kiện áp dụng và đơn vị của các đại lượng. b) Sai: Trong dạng bài này, kết quả không phụ thuộc dữ kiện và có thể bỏ qua điều kiện áp dụng. c) Đúng: Tiếp tuyến với đường sức từ tại một điểm cho phương của vectơ cảm ứng từ; bên ngoài nam châm, đường sức đi từ cực Bắc đến cực Nam. d) Sai: Các đường sức từ có thể cắt nhau tại một điểm. Kiến thức cốt lõi: Tiếp tuyến với đường sức từ tại một điểm cho phương của vectơ cảm ứng từ; bên ngoài nam châm, đường sức đi từ cực Bắc đến cực Nam
\item (Sai) Trong dạng bài này, kết quả không phụ thuộc dữ kiện và có thể bỏ qua điều kiện áp dụng.
\item (Đúng) Tiếp tuyến với đường sức từ tại một điểm cho phương của vectơ cảm ứng từ; bên ngoài nam châm, đường sức đi từ cực Bắc đến cực Nam.
\item (Sai) Các đường sức từ có thể cắt nhau tại một điểm.
\end{itemchoice}
}
\end{ex}

% ===== Câu 131 =====
% ID: L12C3B14-38-TL
% Mức: TH
\dangbt{Nhận biết từ trường, nguồn gây ra từ trường và tương tác từ}
\begin{bt}
% ID: L12C3B14-38-TL
% Mức: TH
Giải thích vì sao nhận định sau đúng: “Từ trường tồn tại xung quanh nam châm và dòng điện; nó gây lực từ lên nam châm hoặc dòng điện khác đặt trong đó.”
\loigiai{
Cần nêu được: Từ trường tồn tại xung quanh nam châm và dòng điện; nó gây lực từ lên nam châm hoặc dòng điện khác đặt trong đó. Khi vận dụng phải xác định đúng dữ kiện, phương chiều nếu có, điều kiện áp dụng, hệ đơn vị và kiểm tra tính hợp lí của kết quả. Nhận định “Mọi điện tích đứng yên đều tạo ra từ trường giống dòng điện.” sai vì trái với kiến thức nêu trên.
}
\end{bt}

% ===== Câu 132 =====
% ID: L12C3B14-38-TLN
% Mức: TH
\begin{ex}
% ID: L12C3B14-38-TLN
% Mức: TH
Trong bốn phát biểu sau về nhận biết từ trường, nguồn gây ra từ trường và tương tác từ, có bao nhiêu phát biểu đúng? (Chỉ ghi một số nguyên.) (1) Mọi điện tích đứng yên đều tạo ra từ trường giống dòng điện. (2) Từ trường tồn tại xung quanh nam châm và dòng điện; nó gây lực từ lên nam châm hoặc dòng điện khác đặt trong đó. (3) Kết quả phải phù hợp ý nghĩa vật lí. (4) Mọi đại lượng đều là vô hướng.
\shortans{1}
\loigiai{
Đối chiếu từng phát biểu với kiến thức: Từ trường tồn tại xung quanh nam châm và dòng điện; nó gây lực từ lên nam châm hoặc dòng điện khác đặt trong đó. Các phát biểu đúng có số lượng là 1.
}
\end{ex}

% ===== Câu 133 =====
% ID: L12C3B14-38-TN
% Mức: TH
\dangbt{So sánh, tổng hợp và biểu diễn các từ trường đơn giản}
\begin{ex}
% ID: L12C3B14-38-TN
% Mức: TH
Một học sinh ôn tập so sánh, tổng hợp và biểu diễn các từ trường đơn giản. Kết luận nào của học sinh là đúng?
\choice
{Độ lớn cảm ứng từ tổng hợp luôn bằng tổng độ lớn các cảm ứng từ thành phần.}
{\True Vectơ cảm ứng từ tổng hợp tại một điểm bằng tổng vectơ các cảm ứng từ thành phần.}
{Có thể bỏ qua phương, chiều và đơn vị của các đại lượng trong mọi trường hợp.}
{Đại lượng đang xét không phụ thuộc các điều kiện vật lí nêu trong bài.}
\loigiai{
Vectơ cảm ứng từ tổng hợp tại một điểm bằng tổng vectơ các cảm ứng từ thành phần. Vì vậy phương án $B$ đúng; các phương án còn lại trái với định nghĩa, định luật hoặc điều kiện áp dụng.
}
\end{ex}

% ===== Câu 134 =====
% ID: L12C3B14-39-DS
% Mức: VD
\dangbt{Mô tả đường sức từ và xác định chiều đường sức từ}
\begin{ex}
% ID: L12C3B14-39-DS
% Mức: VD
Xét mô tả đường sức từ và xác định chiều đường sức từ (Tình huống 4). Hãy xác định tính đúng, sai của bốn phát biểu sau.
\choiceTF
{Trong dạng bài này, kết quả không phụ thuộc dữ kiện và có thể bỏ qua điều kiện áp dụng.}
{\True Khi giải bài thuộc dạng này phải xác định đúng phương, chiều, điều kiện áp dụng và đơn vị của các đại lượng.}
{Các đường sức từ có thể cắt nhau tại một điểm.}
{\True Tiếp tuyến với đường sức từ tại một điểm cho phương của vectơ cảm ứng từ; bên ngoài nam châm, đường sức đi từ cực Bắc đến cực Nam.}
\loigiai{
\begin{itemchoice}
\item (Sai) Trong dạng bài này, kết quả không phụ thuộc dữ kiện và có thể bỏ qua điều kiện áp dụng. (Vì): Trong dạng bài này, kết quả không phụ thuộc dữ kiện và có thể bỏ qua điều kiện áp dụng. b) Đúng: Khi giải bài thuộc dạng này phải xác định đúng phương, chiều, điều kiện áp dụng và đơn vị của các đại lượng. c) Sai: Các đường sức từ có thể cắt nhau tại một điểm. d) Đúng: Tiếp tuyến với đường sức từ tại một điểm cho phương của vectơ cảm ứng từ; bên ngoài nam châm, đường sức đi từ cực Bắc đến cực Nam. Kiến thức cốt lõi: Tiếp tuyến với đường sức từ tại một điểm cho phương của vectơ cảm ứng từ; bên ngoài nam châm, đường sức đi từ cực Bắc đến cực Nam
\item (Đúng) Khi giải bài thuộc dạng này phải xác định đúng phương, chiều, điều kiện áp dụng và đơn vị của các đại lượng.
\item (Sai) Các đường sức từ có thể cắt nhau tại một điểm.
\item (Đúng) Tiếp tuyến với đường sức từ tại một điểm cho phương của vectơ cảm ứng từ; bên ngoài nam châm, đường sức đi từ cực Bắc đến cực Nam.
\end{itemchoice}
}
\end{ex}

% ===== Câu 135 =====
% ID: L12C3B14-39-TL
% Mức: VD
\dangbt{Nhận biết từ trường, nguồn gây ra từ trường và tương tác từ}
\begin{bt}
% ID: L12C3B14-39-TL
% Mức: VD
Phân tích sai lầm trong nhận định: “Mọi điện tích đứng yên đều tạo ra từ trường giống dòng điện.”
\loigiai{
Cần nêu được: Từ trường tồn tại xung quanh nam châm và dòng điện; nó gây lực từ lên nam châm hoặc dòng điện khác đặt trong đó. Khi vận dụng phải xác định đúng dữ kiện, phương chiều nếu có, điều kiện áp dụng, hệ đơn vị và kiểm tra tính hợp lí của kết quả. Nhận định “Mọi điện tích đứng yên đều tạo ra từ trường giống dòng điện.” sai vì trái với kiến thức nêu trên.
}
\end{bt}

% ===== Câu 136 =====
% ID: L12C3B14-39-TLN
% Mức: TH
\begin{ex}
% ID: L12C3B14-39-TLN
% Mức: TH
Trong bốn phát biểu sau về nhận biết từ trường, nguồn gây ra từ trường và tương tác từ, có bao nhiêu phát biểu đúng? (Chỉ ghi một số nguyên.) (1) Cần đổi các đại lượng về hệ đơn vị thống nhất. (2) Từ trường tồn tại xung quanh nam châm và dòng điện; nó gây lực từ lên nam châm hoặc dòng điện khác đặt trong đó. (3) Mọi điện tích đứng yên đều tạo ra từ trường giống dòng điện. (4) Không cần xét chiều của đại lượng vectơ.
\shortans{2}
\loigiai{
Đối chiếu từng phát biểu với kiến thức: Từ trường tồn tại xung quanh nam châm và dòng điện; nó gây lực từ lên nam châm hoặc dòng điện khác đặt trong đó. Các phát biểu đúng có số lượng là 2.
}
\end{ex}

% ===== Câu 137 =====
% ID: L12C3B14-39-TN
% Mức: TH
\dangbt{So sánh, tổng hợp và biểu diễn các từ trường đơn giản}
\begin{ex}
% ID: L12C3B14-39-TN
% Mức: TH
Chọn phát biểu không trái với kiến thức về so sánh, tổng hợp và biểu diễn các từ trường đơn giản.
\choice
{Đại lượng đang xét không phụ thuộc các điều kiện vật lí nêu trong bài.}
{Có thể bỏ qua phương, chiều và đơn vị của các đại lượng trong mọi trường hợp.}
{\True Vectơ cảm ứng từ tổng hợp tại một điểm bằng tổng vectơ các cảm ứng từ thành phần.}
{Độ lớn cảm ứng từ tổng hợp luôn bằng tổng độ lớn các cảm ứng từ thành phần.}
\loigiai{
Vectơ cảm ứng từ tổng hợp tại một điểm bằng tổng vectơ các cảm ứng từ thành phần. Vì vậy phương án $C$ đúng; các phương án còn lại trái với định nghĩa, định luật hoặc điều kiện áp dụng.
}
\end{ex}

% ===== Câu 138 =====
% ID: L12C3B14-40-DS
% Mức: VD
\dangbt{Mô tả đường sức từ và xác định chiều đường sức từ}
\begin{ex}
% ID: L12C3B14-40-DS
% Mức: VD
Xét mô tả đường sức từ và xác định chiều đường sức từ (Tình huống 5). Hãy xác định tính đúng, sai của bốn phát biểu sau.
\choiceTF
{\True Tiếp tuyến với đường sức từ tại một điểm cho phương của vectơ cảm ứng từ; bên ngoài nam châm, đường sức đi từ cực Bắc đến cực Nam.}
{Trong dạng bài này, kết quả không phụ thuộc dữ kiện và có thể bỏ qua điều kiện áp dụng.}
{Các đường sức từ có thể cắt nhau tại một điểm.}
{\True Khi giải bài thuộc dạng này phải xác định đúng phương, chiều, điều kiện áp dụng và đơn vị của các đại lượng.}
\loigiai{
\begin{itemchoice}
\item (Đúng) Tiếp tuyến với đường sức từ tại một điểm cho phương của vectơ cảm ứng từ; bên ngoài nam châm, đường sức đi từ cực Bắc đến cực Nam. (Vì): Tiếp tuyến với đường sức từ tại một điểm cho phương của vectơ cảm ứng từ; bên ngoài nam châm, đường sức đi từ cực Bắc đến cực Nam. b) Sai: Trong dạng bài này, kết quả không phụ thuộc dữ kiện và có thể bỏ qua điều kiện áp dụng. c) Sai: Các đường sức từ có thể cắt nhau tại một điểm. d) Đúng: Khi giải bài thuộc dạng này phải xác định đúng phương, chiều, điều kiện áp dụng và đơn vị của các đại lượng. Kiến thức cốt lõi: Tiếp tuyến với đường sức từ tại một điểm cho phương của vectơ cảm ứng từ; bên ngoài nam châm, đường sức đi từ cực Bắc đến cực Nam
\item (Sai) Trong dạng bài này, kết quả không phụ thuộc dữ kiện và có thể bỏ qua điều kiện áp dụng.
\item (Sai) Các đường sức từ có thể cắt nhau tại một điểm.
\item (Đúng) Khi giải bài thuộc dạng này phải xác định đúng phương, chiều, điều kiện áp dụng và đơn vị của các đại lượng.
\end{itemchoice}
}
\end{ex}

% ===== Câu 139 =====
% ID: L12C3B14-40-TL
% Mức: VD
\dangbt{Nhận biết từ trường, nguồn gây ra từ trường và tương tác từ}
\begin{bt}
% ID: L12C3B14-40-TL
% Mức: VD
Đề xuất các bước giải một bài toán thuộc dạng nhận biết từ trường, nguồn gây ra từ trường và tương tác từ.
\loigiai{
Cần nêu được: Từ trường tồn tại xung quanh nam châm và dòng điện; nó gây lực từ lên nam châm hoặc dòng điện khác đặt trong đó. Khi vận dụng phải xác định đúng dữ kiện, phương chiều nếu có, điều kiện áp dụng, hệ đơn vị và kiểm tra tính hợp lí của kết quả. Nhận định “Mọi điện tích đứng yên đều tạo ra từ trường giống dòng điện.” sai vì trái với kiến thức nêu trên.
}
\end{bt}

% ===== Câu 140 =====
% ID: L12C3B14-40-TLN
% Mức: VD
\begin{ex}
% ID: L12C3B14-40-TLN
% Mức: VD
Trong bốn phát biểu sau về nhận biết từ trường, nguồn gây ra từ trường và tương tác từ, có bao nhiêu phát biểu đúng? (Chỉ ghi một số nguyên.) (1) Từ trường tồn tại xung quanh nam châm và dòng điện; nó gây lực từ lên nam châm hoặc dòng điện khác đặt trong đó. (2) Cần đối chiếu kết quả với điều kiện bài toán. (3) Mọi điện tích đứng yên đều tạo ra từ trường giống dòng điện. (4) Mọi trường hợp đều cho cùng một kết quả.
\shortans{2}
\loigiai{
Đối chiếu từng phát biểu với kiến thức: Từ trường tồn tại xung quanh nam châm và dòng điện; nó gây lực từ lên nam châm hoặc dòng điện khác đặt trong đó. Các phát biểu đúng có số lượng là 2.
}
\end{ex}

% ===== Câu 141 =====
% ID: L12C3B14-40-TN
% Mức: VD
\dangbt{So sánh, tổng hợp và biểu diễn các từ trường đơn giản}
\begin{ex}
% ID: L12C3B14-40-TN
% Mức: VD
Cơ sở Vật lí đúng dùng để giải so sánh, tổng hợp và biểu diễn các từ trường đơn giản là gì?
\choice
{Có thể bỏ qua phương, chiều và đơn vị của các đại lượng trong mọi trường hợp.}
{Đại lượng đang xét không phụ thuộc các điều kiện vật lí nêu trong bài.}
{Độ lớn cảm ứng từ tổng hợp luôn bằng tổng độ lớn các cảm ứng từ thành phần.}
{\True Vectơ cảm ứng từ tổng hợp tại một điểm bằng tổng vectơ các cảm ứng từ thành phần.}
\loigiai{
Vectơ cảm ứng từ tổng hợp tại một điểm bằng tổng vectơ các cảm ứng từ thành phần. Vì vậy phương án $D$ đúng; các phương án còn lại trái với định nghĩa, định luật hoặc điều kiện áp dụng.
}
\end{ex}

% ===== Câu 142 =====
% ID: L12C3B14-41-DS
% Mức: TH
\dangbt{Xác định chiều từ trường bằng kim nam châm và quy tắc nắm tay phải}
\begin{ex}
% ID: L12C3B14-41-DS
% Mức: TH
Xét xác định chiều từ trường bằng kim nam châm và quy tắc nắm tay phải (Tình huống 1). Hãy xác định tính đúng, sai của bốn phát biểu sau.
\choiceTF
{\True Đầu Bắc của kim nam châm nhỏ chỉ theo chiều của vectơ cảm ứng từ; quy tắc nắm tay phải xác định chiều đường sức quanh dòng điện.}
{Quy tắc bàn tay trái dùng để xác định chiều đường sức từ quanh dây dẫn thẳng.}
{\True Khi giải bài thuộc dạng này phải xác định đúng phương, chiều, điều kiện áp dụng và đơn vị của các đại lượng.}
{Trong dạng bài này, kết quả không phụ thuộc dữ kiện và có thể bỏ qua điều kiện áp dụng.}
\loigiai{
\begin{itemchoice}
\item (Đúng) Đầu Bắc của kim nam châm nhỏ chỉ theo chiều của vectơ cảm ứng từ; quy tắc nắm tay phải xác định chiều đường sức quanh dòng điện. (Vì): ầu Bắc của kim nam châm nhỏ chỉ theo chiều của vectơ cảm ứng từ; quy tắc nắm tay phải xác định chiều đường sức quanh dòng điện. b) Sai: Quy tắc bàn tay trái dùng để xác định chiều đường sức từ quanh dây dẫn thẳng. c) Đúng: Khi giải bài thuộc dạng này phải xác định đúng phương, chiều, điều kiện áp dụng và đơn vị của các đại lượng. d) Sai: Trong dạng bài này, kết quả không phụ thuộc dữ kiện và có thể bỏ qua điều kiện áp dụng. Kiến thức cốt lõi: Đầu Bắc của kim nam châm nhỏ chỉ theo chiều của vectơ cảm ứng từ; quy tắc nắm tay phải xác định chiều đường sức quanh dòng điện
\item (Sai) Quy tắc bàn tay trái dùng để xác định chiều đường sức từ quanh dây dẫn thẳng.
\item (Đúng) Khi giải bài thuộc dạng này phải xác định đúng phương, chiều, điều kiện áp dụng và đơn vị của các đại lượng.
\item (Sai) Trong dạng bài này, kết quả không phụ thuộc dữ kiện và có thể bỏ qua điều kiện áp dụng.
\end{itemchoice}
}
\end{ex}

% ===== Câu 143 =====
% ID: L12C3B14-41-TL
% Mức: VD
\dangbt{Nhận biết từ trường, nguồn gây ra từ trường và tương tác từ}
\begin{bt}
% ID: L12C3B14-41-TL
% Mức: VD
Nêu một tình huống thực tiễn liên quan đến nhận biết từ trường, nguồn gây ra từ trường và tương tác từ và giải thích bằng kiến thức Vật lí.
\loigiai{
Cần nêu được: Từ trường tồn tại xung quanh nam châm và dòng điện; nó gây lực từ lên nam châm hoặc dòng điện khác đặt trong đó. Khi vận dụng phải xác định đúng dữ kiện, phương chiều nếu có, điều kiện áp dụng, hệ đơn vị và kiểm tra tính hợp lí của kết quả. Nhận định “Mọi điện tích đứng yên đều tạo ra từ trường giống dòng điện.” sai vì trái với kiến thức nêu trên.
}
\end{bt}

% ===== Câu 144 =====
% ID: L12C3B14-41-TLN
% Mức: VD
\begin{ex}
% ID: L12C3B14-41-TLN
% Mức: VD
Trong bốn phát biểu sau về nhận biết từ trường, nguồn gây ra từ trường và tương tác từ, có bao nhiêu phát biểu đúng? (Chỉ ghi một số nguyên.) (1) Mọi điện tích đứng yên đều tạo ra từ trường giống dòng điện. (2) Cần đọc đúng dữ kiện và giả thiết. (3) Từ trường tồn tại xung quanh nam châm và dòng điện; nó gây lực từ lên nam châm hoặc dòng điện khác đặt trong đó. (4) Có thể thay số trước khi xác định công thức.
\shortans{2}
\loigiai{
Đối chiếu từng phát biểu với kiến thức: Từ trường tồn tại xung quanh nam châm và dòng điện; nó gây lực từ lên nam châm hoặc dòng điện khác đặt trong đó. Các phát biểu đúng có số lượng là 2.
}
\end{ex}

% ===== Câu 145 =====
% ID: L12C3B14-41-TN
% Mức: VD
\dangbt{So sánh, tổng hợp và biểu diễn các từ trường đơn giản}
\begin{ex}
% ID: L12C3B14-41-TN
% Mức: VD
Trong một bài thực hành về so sánh, tổng hợp và biểu diễn các từ trường đơn giản, nhận xét nào đúng?
\choice
{\True Vectơ cảm ứng từ tổng hợp tại một điểm bằng tổng vectơ các cảm ứng từ thành phần.}
{Độ lớn cảm ứng từ tổng hợp luôn bằng tổng độ lớn các cảm ứng từ thành phần.}
{Đại lượng đang xét không phụ thuộc các điều kiện vật lí nêu trong bài.}
{Có thể bỏ qua phương, chiều và đơn vị của các đại lượng trong mọi trường hợp.}
\loigiai{
Vectơ cảm ứng từ tổng hợp tại một điểm bằng tổng vectơ các cảm ứng từ thành phần. Vì vậy phương án $A$ đúng; các phương án còn lại trái với định nghĩa, định luật hoặc điều kiện áp dụng.
}
\end{ex}

% ===== Câu 146 =====
% ID: L12C3B14-42-DS
% Mức: TH
\dangbt{Xác định chiều từ trường bằng kim nam châm và quy tắc nắm tay phải}
\begin{ex}
% ID: L12C3B14-42-DS
% Mức: TH
Xét xác định chiều từ trường bằng kim nam châm và quy tắc nắm tay phải (Tình huống 2). Hãy xác định tính đúng, sai của bốn phát biểu sau.
\choiceTF
{Quy tắc bàn tay trái dùng để xác định chiều đường sức từ quanh dây dẫn thẳng.}
{\True Đầu Bắc của kim nam châm nhỏ chỉ theo chiều của vectơ cảm ứng từ; quy tắc nắm tay phải xác định chiều đường sức quanh dòng điện.}
{Trong dạng bài này, kết quả không phụ thuộc dữ kiện và có thể bỏ qua điều kiện áp dụng.}
{\True Khi giải bài thuộc dạng này phải xác định đúng phương, chiều, điều kiện áp dụng và đơn vị của các đại lượng.}
\loigiai{
\begin{itemchoice}
\item (Sai) Quy tắc bàn tay trái dùng để xác định chiều đường sức từ quanh dây dẫn thẳng. (Vì): Quy tắc bàn tay trái dùng để xác định chiều đường sức từ quanh dây dẫn thẳng. b) Đúng: Đầu Bắc của kim nam châm nhỏ chỉ theo chiều của vectơ cảm ứng từ; quy tắc nắm tay phải xác định chiều đường sức quanh dòng điện. c) Sai: Trong dạng bài này, kết quả không phụ thuộc dữ kiện và có thể bỏ qua điều kiện áp dụng. d) Đúng: Khi giải bài thuộc dạng này phải xác định đúng phương, chiều, điều kiện áp dụng và đơn vị của các đại lượng. Kiến thức cốt lõi: Đầu Bắc của kim nam châm nhỏ chỉ theo chiều của vectơ cảm ứng từ; quy tắc nắm tay phải xác định chiều đường sức quanh dòng điện
\item (Đúng) Đầu Bắc của kim nam châm nhỏ chỉ theo chiều của vectơ cảm ứng từ; quy tắc nắm tay phải xác định chiều đường sức quanh dòng điện.
\item (Sai) Trong dạng bài này, kết quả không phụ thuộc dữ kiện và có thể bỏ qua điều kiện áp dụng.
\item (Đúng) Khi giải bài thuộc dạng này phải xác định đúng phương, chiều, điều kiện áp dụng và đơn vị của các đại lượng.
\end{itemchoice}
}
\end{ex}

% ===== Câu 147 =====
% ID: L12C3B14-42-TL
% Mức: VDC
\dangbt{Nhận biết từ trường, nguồn gây ra từ trường và tương tác từ}
\begin{bt}
% ID: L12C3B14-42-TL
% Mức: VDC
So sánh nhận định đúng “Từ trường tồn tại xung quanh nam châm và dòng điện; nó gây lực từ lên nam châm hoặc dòng điện khác đặt trong đó.” với nhận định sai “Mọi điện tích đứng yên đều tạo ra từ trường giống dòng điện.”; chỉ rõ căn cứ Vật lí.
\loigiai{
Cần nêu được: Từ trường tồn tại xung quanh nam châm và dòng điện; nó gây lực từ lên nam châm hoặc dòng điện khác đặt trong đó. Khi vận dụng phải xác định đúng dữ kiện, phương chiều nếu có, điều kiện áp dụng, hệ đơn vị và kiểm tra tính hợp lí của kết quả. Nhận định “Mọi điện tích đứng yên đều tạo ra từ trường giống dòng điện.” sai vì trái với kiến thức nêu trên.
}
\end{bt}

% ===== Câu 148 =====
% ID: L12C3B14-42-TLN
% Mức: VDC
\begin{ex}
% ID: L12C3B14-42-TLN
% Mức: VDC
Trong bốn phát biểu sau về nhận biết từ trường, nguồn gây ra từ trường và tương tác từ, có bao nhiêu phát biểu đúng? (Chỉ ghi một số nguyên.) (1) Kết quả cần có ý nghĩa vật lí. (2) Mọi điện tích đứng yên đều tạo ra từ trường giống dòng điện. (3) Phải dùng đúng định luật cho tình huống. (4) Từ trường tồn tại xung quanh nam châm và dòng điện; nó gây lực từ lên nam châm hoặc dòng điện khác đặt trong đó.
\shortans{3}
\loigiai{
Đối chiếu từng phát biểu với kiến thức: Từ trường tồn tại xung quanh nam châm và dòng điện; nó gây lực từ lên nam châm hoặc dòng điện khác đặt trong đó. Các phát biểu đúng có số lượng là 3.
}
\end{ex}

% ===== Câu 149 =====
% ID: L12C3B14-42-TN
% Mức: VD
\dangbt{So sánh, tổng hợp và biểu diễn các từ trường đơn giản}
\begin{ex}
% ID: L12C3B14-42-TN
% Mức: VD
Kết luận đúng khi vận dụng so sánh, tổng hợp và biểu diễn các từ trường đơn giản là
\choice
{Độ lớn cảm ứng từ tổng hợp luôn bằng tổng độ lớn các cảm ứng từ thành phần.}
{\True Vectơ cảm ứng từ tổng hợp tại một điểm bằng tổng vectơ các cảm ứng từ thành phần.}
{Có thể bỏ qua phương, chiều và đơn vị của các đại lượng trong mọi trường hợp.}
{Đại lượng đang xét không phụ thuộc các điều kiện vật lí nêu trong bài.}
\loigiai{
Vectơ cảm ứng từ tổng hợp tại một điểm bằng tổng vectơ các cảm ứng từ thành phần. Vì vậy phương án $B$ đúng; các phương án còn lại trái với định nghĩa, định luật hoặc điều kiện áp dụng.
}
\end{ex}

% ===== Câu 150 =====
% ID: L12C3B14-43-DS
% Mức: VD
\dangbt{Xác định chiều từ trường bằng kim nam châm và quy tắc nắm tay phải}
\begin{ex}
% ID: L12C3B14-43-DS
% Mức: VD
Xét xác định chiều từ trường bằng kim nam châm và quy tắc nắm tay phải (Tình huống 3). Hãy xác định tính đúng, sai của bốn phát biểu sau.
\choiceTF
{\True Khi giải bài thuộc dạng này phải xác định đúng phương, chiều, điều kiện áp dụng và đơn vị của các đại lượng.}
{Trong dạng bài này, kết quả không phụ thuộc dữ kiện và có thể bỏ qua điều kiện áp dụng.}
{\True Đầu Bắc của kim nam châm nhỏ chỉ theo chiều của vectơ cảm ứng từ; quy tắc nắm tay phải xác định chiều đường sức quanh dòng điện.}
{Quy tắc bàn tay trái dùng để xác định chiều đường sức từ quanh dây dẫn thẳng.}
\loigiai{
\begin{itemchoice}
\item (Đúng) Khi giải bài thuộc dạng này phải xác định đúng phương, chiều, điều kiện áp dụng và đơn vị của các đại lượng. (Vì): Khi giải bài thuộc dạng này phải xác định đúng phương, chiều, điều kiện áp dụng và đơn vị của các đại lượng. b) Sai: Trong dạng bài này, kết quả không phụ thuộc dữ kiện và có thể bỏ qua điều kiện áp dụng. c) Đúng: Đầu Bắc của kim nam châm nhỏ chỉ theo chiều của vectơ cảm ứng từ; quy tắc nắm tay phải xác định chiều đường sức quanh dòng điện. d) Sai: Quy tắc bàn tay trái dùng để xác định chiều đường sức từ quanh dây dẫn thẳng. Kiến thức cốt lõi: Đầu Bắc của kim nam châm nhỏ chỉ theo chiều của vectơ cảm ứng từ; quy tắc nắm tay phải xác định chiều đường sức quanh dòng điện
\item (Sai) Trong dạng bài này, kết quả không phụ thuộc dữ kiện và có thể bỏ qua điều kiện áp dụng.
\item (Đúng) Đầu Bắc của kim nam châm nhỏ chỉ theo chiều của vectơ cảm ứng từ; quy tắc nắm tay phải xác định chiều đường sức quanh dòng điện.
\item (Sai) Quy tắc bàn tay trái dùng để xác định chiều đường sức từ quanh dây dẫn thẳng.
\end{itemchoice}
}
\end{ex}

% ===== Câu 151 =====
% ID: L12C3B14-43-TL
% Mức: TH
\dangbt{Mô tả đường sức từ và xác định chiều đường sức từ}
\begin{bt}
% ID: L12C3B14-43-TL
% Mức: TH
Trình bày kiến thức cốt lõi của mô tả đường sức từ và xác định chiều đường sức từ và nêu điều kiện áp dụng.
\loigiai{
Cần nêu được: Tiếp tuyến với đường sức từ tại một điểm cho phương của vectơ cảm ứng từ; bên ngoài nam châm, đường sức đi từ cực Bắc đến cực Nam. Khi vận dụng phải xác định đúng dữ kiện, phương chiều nếu có, điều kiện áp dụng, hệ đơn vị và kiểm tra tính hợp lí của kết quả. Nhận định “Các đường sức từ có thể cắt nhau tại một điểm.” sai vì trái với kiến thức nêu trên.
}
\end{bt}

% ===== Câu 152 =====
% ID: L12C3B14-43-TLN
% Mức: TH
\begin{ex}
% ID: L12C3B14-43-TLN
% Mức: TH
Trong bốn phát biểu sau về mô tả đường sức từ và xác định chiều đường sức từ, có bao nhiêu phát biểu đúng? (Chỉ ghi một số nguyên.) (1) Tiếp tuyến với đường sức từ tại một điểm cho phương của vectơ cảm ứng từ; bên ngoài nam châm, đường sức đi từ cực Bắc đến cực Nam. (2) Các đường sức từ có thể cắt nhau tại một điểm. (3) Cần kiểm tra điều kiện áp dụng trước khi tính toán. (4) Có thể bỏ qua đơn vị của kết quả.
\shortans{2}
\loigiai{
Đối chiếu từng phát biểu với kiến thức: Tiếp tuyến với đường sức từ tại một điểm cho phương của vectơ cảm ứng từ; bên ngoài nam châm, đường sức đi từ cực Bắc đến cực Nam. Các phát biểu đúng có số lượng là 2.
}
\end{ex}

% ===== Câu 153 =====
% ID: L12C3B14-43-TN
% Mức: NB
\dangbt{Từ trường của dòng điện thẳng, dòng điện tròn và ống dây}
\begin{ex}
% ID: L12C3B14-43-TN
% Mức: NB
Phát biểu nào sau đây đúng về từ trường của dòng điện thẳng, dòng điện tròn và ống dây?
\choice
{\True Đường sức quanh dây dẫn thẳng là các đường tròn đồng tâm; từ trường trong lòng ống dây dài gần đều.}
{Từ trường trong lòng ống dây dài luôn bằng không.}
{Đại lượng đang xét không phụ thuộc các điều kiện vật lí nêu trong bài.}
{Có thể bỏ qua phương, chiều và đơn vị của các đại lượng trong mọi trường hợp.}
\loigiai{
Đường sức quanh dây dẫn thẳng là các đường tròn đồng tâm; từ trường trong lòng ống dây dài gần đều. Vì vậy phương án $A$ đúng; các phương án còn lại trái với định nghĩa, định luật hoặc điều kiện áp dụng.
}
\end{ex}

% ===== Câu 154 =====
% ID: L12C3B14-44-DS
% Mức: VD
\dangbt{Xác định chiều từ trường bằng kim nam châm và quy tắc nắm tay phải}
\begin{ex}
% ID: L12C3B14-44-DS
% Mức: VD
Xét xác định chiều từ trường bằng kim nam châm và quy tắc nắm tay phải (Tình huống 4). Hãy xác định tính đúng, sai của bốn phát biểu sau.
\choiceTF
{Trong dạng bài này, kết quả không phụ thuộc dữ kiện và có thể bỏ qua điều kiện áp dụng.}
{\True Khi giải bài thuộc dạng này phải xác định đúng phương, chiều, điều kiện áp dụng và đơn vị của các đại lượng.}
{Quy tắc bàn tay trái dùng để xác định chiều đường sức từ quanh dây dẫn thẳng.}
{\True Đầu Bắc của kim nam châm nhỏ chỉ theo chiều của vectơ cảm ứng từ; quy tắc nắm tay phải xác định chiều đường sức quanh dòng điện.}
\loigiai{
\begin{itemchoice}
\item (Sai) Trong dạng bài này, kết quả không phụ thuộc dữ kiện và có thể bỏ qua điều kiện áp dụng. (Vì): Trong dạng bài này, kết quả không phụ thuộc dữ kiện và có thể bỏ qua điều kiện áp dụng. b) Đúng: Khi giải bài thuộc dạng này phải xác định đúng phương, chiều, điều kiện áp dụng và đơn vị của các đại lượng. c) Sai: Quy tắc bàn tay trái dùng để xác định chiều đường sức từ quanh dây dẫn thẳng. d) Đúng: Đầu Bắc của kim nam châm nhỏ chỉ theo chiều của vectơ cảm ứng từ; quy tắc nắm tay phải xác định chiều đường sức quanh dòng điện. Kiến thức cốt lõi: Đầu Bắc của kim nam châm nhỏ chỉ theo chiều của vectơ cảm ứng từ; quy tắc nắm tay phải xác định chiều đường sức quanh dòng điện
\item (Đúng) Khi giải bài thuộc dạng này phải xác định đúng phương, chiều, điều kiện áp dụng và đơn vị của các đại lượng.
\item (Sai) Quy tắc bàn tay trái dùng để xác định chiều đường sức từ quanh dây dẫn thẳng.
\item (Đúng) Đầu Bắc của kim nam châm nhỏ chỉ theo chiều của vectơ cảm ứng từ; quy tắc nắm tay phải xác định chiều đường sức quanh dòng điện.
\end{itemchoice}
}
\end{ex}

% ===== Câu 155 =====
% ID: L12C3B14-44-TL
% Mức: TH
\dangbt{Mô tả đường sức từ và xác định chiều đường sức từ}
\begin{bt}
% ID: L12C3B14-44-TL
% Mức: TH
Giải thích vì sao nhận định sau đúng: “Tiếp tuyến với đường sức từ tại một điểm cho phương của vectơ cảm ứng từ; bên ngoài nam châm, đường sức đi từ cực Bắc đến cực Nam.”
\loigiai{
Cần nêu được: Tiếp tuyến với đường sức từ tại một điểm cho phương của vectơ cảm ứng từ; bên ngoài nam châm, đường sức đi từ cực Bắc đến cực Nam. Khi vận dụng phải xác định đúng dữ kiện, phương chiều nếu có, điều kiện áp dụng, hệ đơn vị và kiểm tra tính hợp lí của kết quả. Nhận định “Các đường sức từ có thể cắt nhau tại một điểm.” sai vì trái với kiến thức nêu trên.
}
\end{bt}

% ===== Câu 156 =====
% ID: L12C3B14-44-TLN
% Mức: TH
\begin{ex}
% ID: L12C3B14-44-TLN
% Mức: TH
Trong bốn phát biểu sau về mô tả đường sức từ và xác định chiều đường sức từ, có bao nhiêu phát biểu đúng? (Chỉ ghi một số nguyên.) (1) Các đường sức từ có thể cắt nhau tại một điểm. (2) Tiếp tuyến với đường sức từ tại một điểm cho phương của vectơ cảm ứng từ; bên ngoài nam châm, đường sức đi từ cực Bắc đến cực Nam. (3) Kết quả phải phù hợp ý nghĩa vật lí. (4) Mọi đại lượng đều là vô hướng.
\shortans{1}
\loigiai{
Đối chiếu từng phát biểu với kiến thức: Tiếp tuyến với đường sức từ tại một điểm cho phương của vectơ cảm ứng từ; bên ngoài nam châm, đường sức đi từ cực Bắc đến cực Nam. Các phát biểu đúng có số lượng là 1.
}
\end{ex}

% ===== Câu 157 =====
% ID: L12C3B14-44-TN
% Mức: NB
\dangbt{Từ trường của dòng điện thẳng, dòng điện tròn và ống dây}
\begin{ex}
% ID: L12C3B14-44-TN
% Mức: NB
Trong từ trường của dòng điện thẳng, dòng điện tròn và ống dây, nhận định nào phù hợp với kiến thức Vật lí?
\choice
{Từ trường trong lòng ống dây dài luôn bằng không.}
{\True Đường sức quanh dây dẫn thẳng là các đường tròn đồng tâm; từ trường trong lòng ống dây dài gần đều.}
{Có thể bỏ qua phương, chiều và đơn vị của các đại lượng trong mọi trường hợp.}
{Đại lượng đang xét không phụ thuộc các điều kiện vật lí nêu trong bài.}
\loigiai{
Đường sức quanh dây dẫn thẳng là các đường tròn đồng tâm; từ trường trong lòng ống dây dài gần đều. Vì vậy phương án $B$ đúng; các phương án còn lại trái với định nghĩa, định luật hoặc điều kiện áp dụng.
}
\end{ex}

% ===== Câu 158 =====
% ID: L12C3B14-45-DS
% Mức: VD
\dangbt{Xác định chiều từ trường bằng kim nam châm và quy tắc nắm tay phải}
\begin{ex}
% ID: L12C3B14-45-DS
% Mức: VD
Xét xác định chiều từ trường bằng kim nam châm và quy tắc nắm tay phải (Tình huống 5). Hãy xác định tính đúng, sai của bốn phát biểu sau.
\choiceTF
{\True Đầu Bắc của kim nam châm nhỏ chỉ theo chiều của vectơ cảm ứng từ; quy tắc nắm tay phải xác định chiều đường sức quanh dòng điện.}
{Trong dạng bài này, kết quả không phụ thuộc dữ kiện và có thể bỏ qua điều kiện áp dụng.}
{Quy tắc bàn tay trái dùng để xác định chiều đường sức từ quanh dây dẫn thẳng.}
{\True Khi giải bài thuộc dạng này phải xác định đúng phương, chiều, điều kiện áp dụng và đơn vị của các đại lượng.}
\loigiai{
\begin{itemchoice}
\item (Đúng) Đầu Bắc của kim nam châm nhỏ chỉ theo chiều của vectơ cảm ứng từ; quy tắc nắm tay phải xác định chiều đường sức quanh dòng điện. (Vì): ầu Bắc của kim nam châm nhỏ chỉ theo chiều của vectơ cảm ứng từ; quy tắc nắm tay phải xác định chiều đường sức quanh dòng điện. b) Sai: Trong dạng bài này, kết quả không phụ thuộc dữ kiện và có thể bỏ qua điều kiện áp dụng. c) Sai: Quy tắc bàn tay trái dùng để xác định chiều đường sức từ quanh dây dẫn thẳng. d) Đúng: Khi giải bài thuộc dạng này phải xác định đúng phương, chiều, điều kiện áp dụng và đơn vị của các đại lượng. Kiến thức cốt lõi: Đầu Bắc của kim nam châm nhỏ chỉ theo chiều của vectơ cảm ứng từ; quy tắc nắm tay phải xác định chiều đường sức quanh dòng điện
\item (Sai) Trong dạng bài này, kết quả không phụ thuộc dữ kiện và có thể bỏ qua điều kiện áp dụng.
\item (Sai) Quy tắc bàn tay trái dùng để xác định chiều đường sức từ quanh dây dẫn thẳng.
\item (Đúng) Khi giải bài thuộc dạng này phải xác định đúng phương, chiều, điều kiện áp dụng và đơn vị của các đại lượng.
\end{itemchoice}
}
\end{ex}

% ===== Câu 159 =====
% ID: L12C3B14-45-TL
% Mức: VD
\dangbt{Mô tả đường sức từ và xác định chiều đường sức từ}
\begin{bt}
% ID: L12C3B14-45-TL
% Mức: VD
Phân tích sai lầm trong nhận định: “Các đường sức từ có thể cắt nhau tại một điểm.”
\loigiai{
Cần nêu được: Tiếp tuyến với đường sức từ tại một điểm cho phương của vectơ cảm ứng từ; bên ngoài nam châm, đường sức đi từ cực Bắc đến cực Nam. Khi vận dụng phải xác định đúng dữ kiện, phương chiều nếu có, điều kiện áp dụng, hệ đơn vị và kiểm tra tính hợp lí của kết quả. Nhận định “Các đường sức từ có thể cắt nhau tại một điểm.” sai vì trái với kiến thức nêu trên.
}
\end{bt}

% ===== Câu 160 =====
% ID: L12C3B14-45-TLN
% Mức: TH
\begin{ex}
% ID: L12C3B14-45-TLN
% Mức: TH
Trong bốn phát biểu sau về mô tả đường sức từ và xác định chiều đường sức từ, có bao nhiêu phát biểu đúng? (Chỉ ghi một số nguyên.) (1) Cần đổi các đại lượng về hệ đơn vị thống nhất. (2) Tiếp tuyến với đường sức từ tại một điểm cho phương của vectơ cảm ứng từ; bên ngoài nam châm, đường sức đi từ cực Bắc đến cực Nam. (3) Các đường sức từ có thể cắt nhau tại một điểm. (4) Không cần xét chiều của đại lượng vectơ.
\shortans{2}
\loigiai{
Đối chiếu từng phát biểu với kiến thức: Tiếp tuyến với đường sức từ tại một điểm cho phương của vectơ cảm ứng từ; bên ngoài nam châm, đường sức đi từ cực Bắc đến cực Nam. Các phát biểu đúng có số lượng là 2.
}
\end{ex}

% ===== Câu 161 =====
% ID: L12C3B14-45-TN
% Mức: NB
\dangbt{Từ trường của dòng điện thẳng, dòng điện tròn và ống dây}
\begin{ex}
% ID: L12C3B14-45-TN
% Mức: NB
Tình huống 3: chọn kết luận chính xác về từ trường của dòng điện thẳng, dòng điện tròn và ống dây.
\choice
{Đại lượng đang xét không phụ thuộc các điều kiện vật lí nêu trong bài.}
{Có thể bỏ qua phương, chiều và đơn vị của các đại lượng trong mọi trường hợp.}
{\True Đường sức quanh dây dẫn thẳng là các đường tròn đồng tâm; từ trường trong lòng ống dây dài gần đều.}
{Từ trường trong lòng ống dây dài luôn bằng không.}
\loigiai{
Đường sức quanh dây dẫn thẳng là các đường tròn đồng tâm; từ trường trong lòng ống dây dài gần đều. Vì vậy phương án $C$ đúng; các phương án còn lại trái với định nghĩa, định luật hoặc điều kiện áp dụng.
}
\end{ex}

% ===== Câu 162 =====
% ID: L12C3B14-46-DS
% Mức: TH
\begin{ex}
% ID: L12C3B14-46-DS
% Mức: TH
Xét từ trường của dòng điện thẳng, dòng điện tròn và ống dây (Tình huống 1). Hãy xác định tính đúng, sai của bốn phát biểu sau.
\choiceTF
{\True Đường sức quanh dây dẫn thẳng là các đường tròn đồng tâm; từ trường trong lòng ống dây dài gần đều.}
{Từ trường trong lòng ống dây dài luôn bằng không.}
{\True Khi giải bài thuộc dạng này phải xác định đúng phương, chiều, điều kiện áp dụng và đơn vị của các đại lượng.}
{Trong dạng bài này, kết quả không phụ thuộc dữ kiện và có thể bỏ qua điều kiện áp dụng.}
\loigiai{
\begin{itemchoice}
\item (Đúng) Đường sức quanh dây dẫn thẳng là các đường tròn đồng tâm; từ trường trong lòng ống dây dài gần đều. (Vì): ường sức quanh dây dẫn thẳng là các đường tròn đồng tâm; từ trường trong lòng ống dây dài gần đều. b) Sai: Từ trường trong lòng ống dây dài luôn bằng không. c) Đúng: Khi giải bài thuộc dạng này phải xác định đúng phương, chiều, điều kiện áp dụng và đơn vị của các đại lượng. d) Sai: Trong dạng bài này, kết quả không phụ thuộc dữ kiện và có thể bỏ qua điều kiện áp dụng. Kiến thức cốt lõi: Đường sức quanh dây dẫn thẳng là các đường tròn đồng tâm; từ trường trong lòng ống dây dài gần đều
\item (Sai) Từ trường trong lòng ống dây dài luôn bằng không.
\item (Đúng) Khi giải bài thuộc dạng này phải xác định đúng phương, chiều, điều kiện áp dụng và đơn vị của các đại lượng.
\item (Sai) Trong dạng bài này, kết quả không phụ thuộc dữ kiện và có thể bỏ qua điều kiện áp dụng.
\end{itemchoice}
}
\end{ex}

% ===== Câu 163 =====
% ID: L12C3B14-46-TL
% Mức: VD
\dangbt{Mô tả đường sức từ và xác định chiều đường sức từ}
\begin{bt}
% ID: L12C3B14-46-TL
% Mức: VD
Đề xuất các bước giải một bài toán thuộc dạng mô tả đường sức từ và xác định chiều đường sức từ.
\loigiai{
Cần nêu được: Tiếp tuyến với đường sức từ tại một điểm cho phương của vectơ cảm ứng từ; bên ngoài nam châm, đường sức đi từ cực Bắc đến cực Nam. Khi vận dụng phải xác định đúng dữ kiện, phương chiều nếu có, điều kiện áp dụng, hệ đơn vị và kiểm tra tính hợp lí của kết quả. Nhận định “Các đường sức từ có thể cắt nhau tại một điểm.” sai vì trái với kiến thức nêu trên.
}
\end{bt}

% ===== Câu 164 =====
% ID: L12C3B14-46-TLN
% Mức: VD
\begin{ex}
% ID: L12C3B14-46-TLN
% Mức: VD
Trong bốn phát biểu sau về mô tả đường sức từ và xác định chiều đường sức từ, có bao nhiêu phát biểu đúng? (Chỉ ghi một số nguyên.) (1) Tiếp tuyến với đường sức từ tại một điểm cho phương của vectơ cảm ứng từ; bên ngoài nam châm, đường sức đi từ cực Bắc đến cực Nam. (2) Cần đối chiếu kết quả với điều kiện bài toán. (3) Các đường sức từ có thể cắt nhau tại một điểm. (4) Mọi trường hợp đều cho cùng một kết quả.
\shortans{2}
\loigiai{
Đối chiếu từng phát biểu với kiến thức: Tiếp tuyến với đường sức từ tại một điểm cho phương của vectơ cảm ứng từ; bên ngoài nam châm, đường sức đi từ cực Bắc đến cực Nam. Các phát biểu đúng có số lượng là 2.
}
\end{ex}

% ===== Câu 165 =====
% ID: L12C3B14-46-TN
% Mức: TH
\dangbt{Từ trường của dòng điện thẳng, dòng điện tròn và ống dây}
\begin{ex}
% ID: L12C3B14-46-TN
% Mức: TH
Nội dung nào mô tả đúng nhất từ trường của dòng điện thẳng, dòng điện tròn và ống dây?
\choice
{Có thể bỏ qua phương, chiều và đơn vị của các đại lượng trong mọi trường hợp.}
{Đại lượng đang xét không phụ thuộc các điều kiện vật lí nêu trong bài.}
{Từ trường trong lòng ống dây dài luôn bằng không.}
{\True Đường sức quanh dây dẫn thẳng là các đường tròn đồng tâm; từ trường trong lòng ống dây dài gần đều.}
\loigiai{
Đường sức quanh dây dẫn thẳng là các đường tròn đồng tâm; từ trường trong lòng ống dây dài gần đều. Vì vậy phương án $D$ đúng; các phương án còn lại trái với định nghĩa, định luật hoặc điều kiện áp dụng.
}
\end{ex}

% ===== Câu 166 =====
% ID: L12C3B14-47-DS
% Mức: TH
\begin{ex}
% ID: L12C3B14-47-DS
% Mức: TH
Xét từ trường của dòng điện thẳng, dòng điện tròn và ống dây (Tình huống 2). Hãy xác định tính đúng, sai của bốn phát biểu sau.
\choiceTF
{Từ trường trong lòng ống dây dài luôn bằng không.}
{\True Đường sức quanh dây dẫn thẳng là các đường tròn đồng tâm; từ trường trong lòng ống dây dài gần đều.}
{Trong dạng bài này, kết quả không phụ thuộc dữ kiện và có thể bỏ qua điều kiện áp dụng.}
{\True Khi giải bài thuộc dạng này phải xác định đúng phương, chiều, điều kiện áp dụng và đơn vị của các đại lượng.}
\loigiai{
\begin{itemchoice}
\item (Sai) Từ trường trong lòng ống dây dài luôn bằng không. (Vì): Từ trường trong lòng ống dây dài luôn bằng không. b) Đúng: Đường sức quanh dây dẫn thẳng là các đường tròn đồng tâm; từ trường trong lòng ống dây dài gần đều. c) Sai: Trong dạng bài này, kết quả không phụ thuộc dữ kiện và có thể bỏ qua điều kiện áp dụng. d) Đúng: Khi giải bài thuộc dạng này phải xác định đúng phương, chiều, điều kiện áp dụng và đơn vị của các đại lượng. Kiến thức cốt lõi: Đường sức quanh dây dẫn thẳng là các đường tròn đồng tâm; từ trường trong lòng ống dây dài gần đều
\item (Đúng) Đường sức quanh dây dẫn thẳng là các đường tròn đồng tâm; từ trường trong lòng ống dây dài gần đều.
\item (Sai) Trong dạng bài này, kết quả không phụ thuộc dữ kiện và có thể bỏ qua điều kiện áp dụng.
\item (Đúng) Khi giải bài thuộc dạng này phải xác định đúng phương, chiều, điều kiện áp dụng và đơn vị của các đại lượng.
\end{itemchoice}
}
\end{ex}

% ===== Câu 167 =====
% ID: L12C3B14-47-TL
% Mức: VD
\dangbt{Mô tả đường sức từ và xác định chiều đường sức từ}
\begin{bt}
% ID: L12C3B14-47-TL
% Mức: VD
Nêu một tình huống thực tiễn liên quan đến mô tả đường sức từ và xác định chiều đường sức từ và giải thích bằng kiến thức Vật lí.
\loigiai{
Cần nêu được: Tiếp tuyến với đường sức từ tại một điểm cho phương của vectơ cảm ứng từ; bên ngoài nam châm, đường sức đi từ cực Bắc đến cực Nam. Khi vận dụng phải xác định đúng dữ kiện, phương chiều nếu có, điều kiện áp dụng, hệ đơn vị và kiểm tra tính hợp lí của kết quả. Nhận định “Các đường sức từ có thể cắt nhau tại một điểm.” sai vì trái với kiến thức nêu trên.
}
\end{bt}

% ===== Câu 168 =====
% ID: L12C3B14-47-TLN
% Mức: VD
\begin{ex}
% ID: L12C3B14-47-TLN
% Mức: VD
Trong bốn phát biểu sau về mô tả đường sức từ và xác định chiều đường sức từ, có bao nhiêu phát biểu đúng? (Chỉ ghi một số nguyên.) (1) Các đường sức từ có thể cắt nhau tại một điểm. (2) Cần đọc đúng dữ kiện và giả thiết. (3) Tiếp tuyến với đường sức từ tại một điểm cho phương của vectơ cảm ứng từ; bên ngoài nam châm, đường sức đi từ cực Bắc đến cực Nam. (4) Có thể thay số trước khi xác định công thức.
\shortans{2}
\loigiai{
Đối chiếu từng phát biểu với kiến thức: Tiếp tuyến với đường sức từ tại một điểm cho phương của vectơ cảm ứng từ; bên ngoài nam châm, đường sức đi từ cực Bắc đến cực Nam. Các phát biểu đúng có số lượng là 2.
}
\end{ex}

% ===== Câu 169 =====
% ID: L12C3B14-47-TN
% Mức: TH
\dangbt{Từ trường của dòng điện thẳng, dòng điện tròn và ống dây}
\begin{ex}
% ID: L12C3B14-47-TN
% Mức: TH
Khi phân tích từ trường của dòng điện thẳng, dòng điện tròn và ống dây, phát biểu nào đúng?
\choice
{\True Đường sức quanh dây dẫn thẳng là các đường tròn đồng tâm; từ trường trong lòng ống dây dài gần đều.}
{Từ trường trong lòng ống dây dài luôn bằng không.}
{Đại lượng đang xét không phụ thuộc các điều kiện vật lí nêu trong bài.}
{Có thể bỏ qua phương, chiều và đơn vị của các đại lượng trong mọi trường hợp.}
\loigiai{
Đường sức quanh dây dẫn thẳng là các đường tròn đồng tâm; từ trường trong lòng ống dây dài gần đều. Vì vậy phương án $A$ đúng; các phương án còn lại trái với định nghĩa, định luật hoặc điều kiện áp dụng.
}
\end{ex}

% ===== Câu 170 =====
% ID: L12C3B14-48-DS
% Mức: VD
\begin{ex}
% ID: L12C3B14-48-DS
% Mức: VD
Xét từ trường của dòng điện thẳng, dòng điện tròn và ống dây (Tình huống 3). Hãy xác định tính đúng, sai của bốn phát biểu sau.
\choiceTF
{\True Khi giải bài thuộc dạng này phải xác định đúng phương, chiều, điều kiện áp dụng và đơn vị của các đại lượng.}
{Trong dạng bài này, kết quả không phụ thuộc dữ kiện và có thể bỏ qua điều kiện áp dụng.}
{\True Đường sức quanh dây dẫn thẳng là các đường tròn đồng tâm; từ trường trong lòng ống dây dài gần đều.}
{Từ trường trong lòng ống dây dài luôn bằng không.}
\loigiai{
\begin{itemchoice}
\item (Đúng) Khi giải bài thuộc dạng này phải xác định đúng phương, chiều, điều kiện áp dụng và đơn vị của các đại lượng. (Vì): Khi giải bài thuộc dạng này phải xác định đúng phương, chiều, điều kiện áp dụng và đơn vị của các đại lượng. b) Sai: Trong dạng bài này, kết quả không phụ thuộc dữ kiện và có thể bỏ qua điều kiện áp dụng. c) Đúng: Đường sức quanh dây dẫn thẳng là các đường tròn đồng tâm; từ trường trong lòng ống dây dài gần đều. d) Sai: Từ trường trong lòng ống dây dài luôn bằng không. Kiến thức cốt lõi: Đường sức quanh dây dẫn thẳng là các đường tròn đồng tâm; từ trường trong lòng ống dây dài gần đều
\item (Sai) Trong dạng bài này, kết quả không phụ thuộc dữ kiện và có thể bỏ qua điều kiện áp dụng.
\item (Đúng) Đường sức quanh dây dẫn thẳng là các đường tròn đồng tâm; từ trường trong lòng ống dây dài gần đều.
\item (Sai) Từ trường trong lòng ống dây dài luôn bằng không.
\end{itemchoice}
}
\end{ex}

% ===== Câu 171 =====
% ID: L12C3B14-48-TL
% Mức: VDC
\dangbt{Mô tả đường sức từ và xác định chiều đường sức từ}
\begin{bt}
% ID: L12C3B14-48-TL
% Mức: VDC
So sánh nhận định đúng “Tiếp tuyến với đường sức từ tại một điểm cho phương của vectơ cảm ứng từ; bên ngoài nam châm, đường sức đi từ cực Bắc đến cực Nam.” với nhận định sai “Các đường sức từ có thể cắt nhau tại một điểm.”; chỉ rõ căn cứ Vật lí.
\loigiai{
Cần nêu được: Tiếp tuyến với đường sức từ tại một điểm cho phương của vectơ cảm ứng từ; bên ngoài nam châm, đường sức đi từ cực Bắc đến cực Nam. Khi vận dụng phải xác định đúng dữ kiện, phương chiều nếu có, điều kiện áp dụng, hệ đơn vị và kiểm tra tính hợp lí của kết quả. Nhận định “Các đường sức từ có thể cắt nhau tại một điểm.” sai vì trái với kiến thức nêu trên.
}
\end{bt}

% ===== Câu 172 =====
% ID: L12C3B14-48-TLN
% Mức: VDC
\begin{ex}
% ID: L12C3B14-48-TLN
% Mức: VDC
Trong bốn phát biểu sau về mô tả đường sức từ và xác định chiều đường sức từ, có bao nhiêu phát biểu đúng? (Chỉ ghi một số nguyên.) (1) Kết quả cần có ý nghĩa vật lí. (2) Các đường sức từ có thể cắt nhau tại một điểm. (3) Phải dùng đúng định luật cho tình huống. (4) Tiếp tuyến với đường sức từ tại một điểm cho phương của vectơ cảm ứng từ; bên ngoài nam châm, đường sức đi từ cực Bắc đến cực Nam.
\shortans{3}
\loigiai{
Đối chiếu từng phát biểu với kiến thức: Tiếp tuyến với đường sức từ tại một điểm cho phương của vectơ cảm ứng từ; bên ngoài nam châm, đường sức đi từ cực Bắc đến cực Nam. Các phát biểu đúng có số lượng là 3.
}
\end{ex}

% ===== Câu 173 =====
% ID: L12C3B14-48-TN
% Mức: TH
\dangbt{Từ trường của dòng điện thẳng, dòng điện tròn và ống dây}
\begin{ex}
% ID: L12C3B14-48-TN
% Mức: TH
Một học sinh ôn tập từ trường của dòng điện thẳng, dòng điện tròn và ống dây. Kết luận nào của học sinh là đúng?
\choice
{Từ trường trong lòng ống dây dài luôn bằng không.}
{\True Đường sức quanh dây dẫn thẳng là các đường tròn đồng tâm; từ trường trong lòng ống dây dài gần đều.}
{Có thể bỏ qua phương, chiều và đơn vị của các đại lượng trong mọi trường hợp.}
{Đại lượng đang xét không phụ thuộc các điều kiện vật lí nêu trong bài.}
\loigiai{
Đường sức quanh dây dẫn thẳng là các đường tròn đồng tâm; từ trường trong lòng ống dây dài gần đều. Vì vậy phương án $B$ đúng; các phương án còn lại trái với định nghĩa, định luật hoặc điều kiện áp dụng.
}
\end{ex}

% ===== Câu 174 =====
% ID: L12C3B14-49-DS
% Mức: VD
\begin{ex}
% ID: L12C3B14-49-DS
% Mức: VD
Xét từ trường của dòng điện thẳng, dòng điện tròn và ống dây (Tình huống 4). Hãy xác định tính đúng, sai của bốn phát biểu sau.
\choiceTF
{Trong dạng bài này, kết quả không phụ thuộc dữ kiện và có thể bỏ qua điều kiện áp dụng.}
{\True Khi giải bài thuộc dạng này phải xác định đúng phương, chiều, điều kiện áp dụng và đơn vị của các đại lượng.}
{Từ trường trong lòng ống dây dài luôn bằng không.}
{\True Đường sức quanh dây dẫn thẳng là các đường tròn đồng tâm; từ trường trong lòng ống dây dài gần đều.}
\loigiai{
\begin{itemchoice}
\item (Sai) Trong dạng bài này, kết quả không phụ thuộc dữ kiện và có thể bỏ qua điều kiện áp dụng. (Vì): Trong dạng bài này, kết quả không phụ thuộc dữ kiện và có thể bỏ qua điều kiện áp dụng. b) Đúng: Khi giải bài thuộc dạng này phải xác định đúng phương, chiều, điều kiện áp dụng và đơn vị của các đại lượng. c) Sai: Từ trường trong lòng ống dây dài luôn bằng không. d) Đúng: Đường sức quanh dây dẫn thẳng là các đường tròn đồng tâm; từ trường trong lòng ống dây dài gần đều. Kiến thức cốt lõi: Đường sức quanh dây dẫn thẳng là các đường tròn đồng tâm; từ trường trong lòng ống dây dài gần đều
\item (Đúng) Khi giải bài thuộc dạng này phải xác định đúng phương, chiều, điều kiện áp dụng và đơn vị của các đại lượng.
\item (Sai) Từ trường trong lòng ống dây dài luôn bằng không.
\item (Đúng) Đường sức quanh dây dẫn thẳng là các đường tròn đồng tâm; từ trường trong lòng ống dây dài gần đều.
\end{itemchoice}
}
\end{ex}

% ===== Câu 175 =====
% ID: L12C3B14-49-TL
% Mức: TH
\dangbt{Xác định chiều từ trường bằng kim nam châm và quy tắc nắm tay phải}
\begin{bt}
% ID: L12C3B14-49-TL
% Mức: TH
Trình bày kiến thức cốt lõi của xác định chiều từ trường bằng kim nam châm và quy tắc nắm tay phải và nêu điều kiện áp dụng.
\loigiai{
Cần nêu được: Đầu Bắc của kim nam châm nhỏ chỉ theo chiều của vectơ cảm ứng từ; quy tắc nắm tay phải xác định chiều đường sức quanh dòng điện. Khi vận dụng phải xác định đúng dữ kiện, phương chiều nếu có, điều kiện áp dụng, hệ đơn vị và kiểm tra tính hợp lí của kết quả. Nhận định “Quy tắc bàn tay trái dùng để xác định chiều đường sức từ quanh dây dẫn thẳng.” sai vì trái với kiến thức nêu trên.
}
\end{bt}

% ===== Câu 176 =====
% ID: L12C3B14-49-TLN
% Mức: TH
\begin{ex}
% ID: L12C3B14-49-TLN
% Mức: TH
Trong bốn phát biểu sau về xác định chiều từ trường bằng kim nam châm và quy tắc nắm tay phải, có bao nhiêu phát biểu đúng? (Chỉ ghi một số nguyên.) (1) Đầu Bắc của kim nam châm nhỏ chỉ theo chiều của vectơ cảm ứng từ; quy tắc nắm tay phải xác định chiều đường sức quanh dòng điện. (2) Quy tắc bàn tay trái dùng để xác định chiều đường sức từ quanh dây dẫn thẳng. (3) Cần kiểm tra điều kiện áp dụng trước khi tính toán. (4) Có thể bỏ qua đơn vị của kết quả.
\shortans{2}
\loigiai{
Đối chiếu từng phát biểu với kiến thức: Đầu Bắc của kim nam châm nhỏ chỉ theo chiều của vectơ cảm ứng từ; quy tắc nắm tay phải xác định chiều đường sức quanh dòng điện. Các phát biểu đúng có số lượng là 2.
}
\end{ex}

% ===== Câu 177 =====
% ID: L12C3B14-49-TN
% Mức: TH
\dangbt{Từ trường của dòng điện thẳng, dòng điện tròn và ống dây}
\begin{ex}
% ID: L12C3B14-49-TN
% Mức: TH
Chọn phát biểu không trái với kiến thức về từ trường của dòng điện thẳng, dòng điện tròn và ống dây.
\choice
{Đại lượng đang xét không phụ thuộc các điều kiện vật lí nêu trong bài.}
{Có thể bỏ qua phương, chiều và đơn vị của các đại lượng trong mọi trường hợp.}
{\True Đường sức quanh dây dẫn thẳng là các đường tròn đồng tâm; từ trường trong lòng ống dây dài gần đều.}
{Từ trường trong lòng ống dây dài luôn bằng không.}
\loigiai{
Đường sức quanh dây dẫn thẳng là các đường tròn đồng tâm; từ trường trong lòng ống dây dài gần đều. Vì vậy phương án $C$ đúng; các phương án còn lại trái với định nghĩa, định luật hoặc điều kiện áp dụng.
}
\end{ex}

% ===== Câu 178 =====
% ID: L12C3B14-50-DS
% Mức: VD
\begin{ex}
% ID: L12C3B14-50-DS
% Mức: VD
Xét từ trường của dòng điện thẳng, dòng điện tròn và ống dây (Tình huống 5). Hãy xác định tính đúng, sai của bốn phát biểu sau.
\choiceTF
{\True Đường sức quanh dây dẫn thẳng là các đường tròn đồng tâm; từ trường trong lòng ống dây dài gần đều.}
{Trong dạng bài này, kết quả không phụ thuộc dữ kiện và có thể bỏ qua điều kiện áp dụng.}
{Từ trường trong lòng ống dây dài luôn bằng không.}
{\True Khi giải bài thuộc dạng này phải xác định đúng phương, chiều, điều kiện áp dụng và đơn vị của các đại lượng.}
\loigiai{
\begin{itemchoice}
\item (Đúng) Đường sức quanh dây dẫn thẳng là các đường tròn đồng tâm; từ trường trong lòng ống dây dài gần đều. (Vì): ường sức quanh dây dẫn thẳng là các đường tròn đồng tâm; từ trường trong lòng ống dây dài gần đều. b) Sai: Trong dạng bài này, kết quả không phụ thuộc dữ kiện và có thể bỏ qua điều kiện áp dụng. c) Sai: Từ trường trong lòng ống dây dài luôn bằng không. d) Đúng: Khi giải bài thuộc dạng này phải xác định đúng phương, chiều, điều kiện áp dụng và đơn vị của các đại lượng. Kiến thức cốt lõi: Đường sức quanh dây dẫn thẳng là các đường tròn đồng tâm; từ trường trong lòng ống dây dài gần đều
\item (Sai) Trong dạng bài này, kết quả không phụ thuộc dữ kiện và có thể bỏ qua điều kiện áp dụng.
\item (Sai) Từ trường trong lòng ống dây dài luôn bằng không.
\item (Đúng) Khi giải bài thuộc dạng này phải xác định đúng phương, chiều, điều kiện áp dụng và đơn vị của các đại lượng.
\end{itemchoice}
}
\end{ex}

% ===== Câu 179 =====
% ID: L12C3B14-50-TL
% Mức: TH
\dangbt{Xác định chiều từ trường bằng kim nam châm và quy tắc nắm tay phải}
\begin{bt}
% ID: L12C3B14-50-TL
% Mức: TH
Giải thích vì sao nhận định sau đúng: “Đầu Bắc của kim nam châm nhỏ chỉ theo chiều của vectơ cảm ứng từ; quy tắc nắm tay phải xác định chiều đường sức quanh dòng điện.”
\loigiai{
Cần nêu được: Đầu Bắc của kim nam châm nhỏ chỉ theo chiều của vectơ cảm ứng từ; quy tắc nắm tay phải xác định chiều đường sức quanh dòng điện. Khi vận dụng phải xác định đúng dữ kiện, phương chiều nếu có, điều kiện áp dụng, hệ đơn vị và kiểm tra tính hợp lí của kết quả. Nhận định “Quy tắc bàn tay trái dùng để xác định chiều đường sức từ quanh dây dẫn thẳng.” sai vì trái với kiến thức nêu trên.
}
\end{bt}

% ===== Câu 180 =====
% ID: L12C3B14-50-TLN
% Mức: TH
\begin{ex}
% ID: L12C3B14-50-TLN
% Mức: TH
Trong bốn phát biểu sau về xác định chiều từ trường bằng kim nam châm và quy tắc nắm tay phải, có bao nhiêu phát biểu đúng? (Chỉ ghi một số nguyên.) (1) Quy tắc bàn tay trái dùng để xác định chiều đường sức từ quanh dây dẫn thẳng. (2) Đầu Bắc của kim nam châm nhỏ chỉ theo chiều của vectơ cảm ứng từ; quy tắc nắm tay phải xác định chiều đường sức quanh dòng điện. (3) Kết quả phải phù hợp ý nghĩa vật lí. (4) Mọi đại lượng đều là vô hướng.
\shortans{1}
\loigiai{
Đối chiếu từng phát biểu với kiến thức: Đầu Bắc của kim nam châm nhỏ chỉ theo chiều của vectơ cảm ứng từ; quy tắc nắm tay phải xác định chiều đường sức quanh dòng điện. Các phát biểu đúng có số lượng là 1.
}
\end{ex}

% ===== Câu 181 =====
% ID: L12C3B14-50-TN
% Mức: VD
\dangbt{Từ trường của dòng điện thẳng, dòng điện tròn và ống dây}
\begin{ex}
% ID: L12C3B14-50-TN
% Mức: VD
Cơ sở Vật lí đúng dùng để giải từ trường của dòng điện thẳng, dòng điện tròn và ống dây là gì?
\choice
{Có thể bỏ qua phương, chiều và đơn vị của các đại lượng trong mọi trường hợp.}
{Đại lượng đang xét không phụ thuộc các điều kiện vật lí nêu trong bài.}
{Từ trường trong lòng ống dây dài luôn bằng không.}
{\True Đường sức quanh dây dẫn thẳng là các đường tròn đồng tâm; từ trường trong lòng ống dây dài gần đều.}
\loigiai{
Đường sức quanh dây dẫn thẳng là các đường tròn đồng tâm; từ trường trong lòng ống dây dài gần đều. Vì vậy phương án $D$ đúng; các phương án còn lại trái với định nghĩa, định luật hoặc điều kiện áp dụng.
}
\end{ex}

% ===== Câu 182 =====
% ID: L12C3B14-51-DS
% Mức: TH
\dangbt{So sánh, tổng hợp và biểu diễn các từ trường đơn giản}
\begin{ex}
% ID: L12C3B14-51-DS
% Mức: TH
Xét so sánh, tổng hợp và biểu diễn các từ trường đơn giản (Tình huống 1). Hãy xác định tính đúng, sai của bốn phát biểu sau.
\choiceTF
{\True Vectơ cảm ứng từ tổng hợp tại một điểm bằng tổng vectơ các cảm ứng từ thành phần.}
{Độ lớn cảm ứng từ tổng hợp luôn bằng tổng độ lớn các cảm ứng từ thành phần.}
{\True Khi giải bài thuộc dạng này phải xác định đúng phương, chiều, điều kiện áp dụng và đơn vị của các đại lượng.}
{Trong dạng bài này, kết quả không phụ thuộc dữ kiện và có thể bỏ qua điều kiện áp dụng.}
\loigiai{
\begin{itemchoice}
\item (Đúng) Vectơ cảm ứng từ tổng hợp tại một điểm bằng tổng vectơ các cảm ứng từ thành phần. (Vì): Vectơ cảm ứng từ tổng hợp tại một điểm bằng tổng vectơ các cảm ứng từ thành phần. b) Sai: Độ lớn cảm ứng từ tổng hợp luôn bằng tổng độ lớn các cảm ứng từ thành phần. c) Đúng: Khi giải bài thuộc dạng này phải xác định đúng phương, chiều, điều kiện áp dụng và đơn vị của các đại lượng. d) Sai: Trong dạng bài này, kết quả không phụ thuộc dữ kiện và có thể bỏ qua điều kiện áp dụng. Kiến thức cốt lõi: Vectơ cảm ứng từ tổng hợp tại một điểm bằng tổng vectơ các cảm ứng từ thành phần
\item (Sai) Độ lớn cảm ứng từ tổng hợp luôn bằng tổng độ lớn các cảm ứng từ thành phần.
\item (Đúng) Khi giải bài thuộc dạng này phải xác định đúng phương, chiều, điều kiện áp dụng và đơn vị của các đại lượng.
\item (Sai) Trong dạng bài này, kết quả không phụ thuộc dữ kiện và có thể bỏ qua điều kiện áp dụng.
\end{itemchoice}
}
\end{ex}

% ===== Câu 183 =====
% ID: L12C3B14-51-TL
% Mức: VD
\dangbt{Xác định chiều từ trường bằng kim nam châm và quy tắc nắm tay phải}
\begin{bt}
% ID: L12C3B14-51-TL
% Mức: VD
Phân tích sai lầm trong nhận định: “Quy tắc bàn tay trái dùng để xác định chiều đường sức từ quanh dây dẫn thẳng.”
\loigiai{
Cần nêu được: Đầu Bắc của kim nam châm nhỏ chỉ theo chiều của vectơ cảm ứng từ; quy tắc nắm tay phải xác định chiều đường sức quanh dòng điện. Khi vận dụng phải xác định đúng dữ kiện, phương chiều nếu có, điều kiện áp dụng, hệ đơn vị và kiểm tra tính hợp lí của kết quả. Nhận định “Quy tắc bàn tay trái dùng để xác định chiều đường sức từ quanh dây dẫn thẳng.” sai vì trái với kiến thức nêu trên.
}
\end{bt}

% ===== Câu 184 =====
% ID: L12C3B14-51-TLN
% Mức: TH
\begin{ex}
% ID: L12C3B14-51-TLN
% Mức: TH
Trong bốn phát biểu sau về xác định chiều từ trường bằng kim nam châm và quy tắc nắm tay phải, có bao nhiêu phát biểu đúng? (Chỉ ghi một số nguyên.) (1) Cần đổi các đại lượng về hệ đơn vị thống nhất. (2) Đầu Bắc của kim nam châm nhỏ chỉ theo chiều của vectơ cảm ứng từ; quy tắc nắm tay phải xác định chiều đường sức quanh dòng điện. (3) Quy tắc bàn tay trái dùng để xác định chiều đường sức từ quanh dây dẫn thẳng. (4) Không cần xét chiều của đại lượng vectơ.
\shortans{2}
\loigiai{
Đối chiếu từng phát biểu với kiến thức: Đầu Bắc của kim nam châm nhỏ chỉ theo chiều của vectơ cảm ứng từ; quy tắc nắm tay phải xác định chiều đường sức quanh dòng điện. Các phát biểu đúng có số lượng là 2.
}
\end{ex}

% ===== Câu 185 =====
% ID: L12C3B14-51-TN
% Mức: VD
\dangbt{Từ trường của dòng điện thẳng, dòng điện tròn và ống dây}
\begin{ex}
% ID: L12C3B14-51-TN
% Mức: VD
Trong một bài thực hành về từ trường của dòng điện thẳng, dòng điện tròn và ống dây, nhận xét nào đúng?
\choice
{\True Đường sức quanh dây dẫn thẳng là các đường tròn đồng tâm; từ trường trong lòng ống dây dài gần đều.}
{Từ trường trong lòng ống dây dài luôn bằng không.}
{Đại lượng đang xét không phụ thuộc các điều kiện vật lí nêu trong bài.}
{Có thể bỏ qua phương, chiều và đơn vị của các đại lượng trong mọi trường hợp.}
\loigiai{
Đường sức quanh dây dẫn thẳng là các đường tròn đồng tâm; từ trường trong lòng ống dây dài gần đều. Vì vậy phương án $A$ đúng; các phương án còn lại trái với định nghĩa, định luật hoặc điều kiện áp dụng.
}
\end{ex}

% ===== Câu 186 =====
% ID: L12C3B14-52-DS
% Mức: TH
\dangbt{So sánh, tổng hợp và biểu diễn các từ trường đơn giản}
\begin{ex}
% ID: L12C3B14-52-DS
% Mức: TH
Xét so sánh, tổng hợp và biểu diễn các từ trường đơn giản (Tình huống 2). Hãy xác định tính đúng, sai của bốn phát biểu sau.
\choiceTF
{Độ lớn cảm ứng từ tổng hợp luôn bằng tổng độ lớn các cảm ứng từ thành phần.}
{\True Vectơ cảm ứng từ tổng hợp tại một điểm bằng tổng vectơ các cảm ứng từ thành phần.}
{Trong dạng bài này, kết quả không phụ thuộc dữ kiện và có thể bỏ qua điều kiện áp dụng.}
{\True Khi giải bài thuộc dạng này phải xác định đúng phương, chiều, điều kiện áp dụng và đơn vị của các đại lượng.}
\loigiai{
\begin{itemchoice}
\item (Sai) Độ lớn cảm ứng từ tổng hợp luôn bằng tổng độ lớn các cảm ứng từ thành phần. (Vì): Độ lớn cảm ứng từ tổng hợp luôn bằng tổng độ lớn các cảm ứng từ thành phần. b) Đúng: Vectơ cảm ứng từ tổng hợp tại một điểm bằng tổng vectơ các cảm ứng từ thành phần. c) Sai: Trong dạng bài này, kết quả không phụ thuộc dữ kiện và có thể bỏ qua điều kiện áp dụng. d) Đúng: Khi giải bài thuộc dạng này phải xác định đúng phương, chiều, điều kiện áp dụng và đơn vị của các đại lượng. Kiến thức cốt lõi: Vectơ cảm ứng từ tổng hợp tại một điểm bằng tổng vectơ các cảm ứng từ thành phần
\item (Đúng) Vectơ cảm ứng từ tổng hợp tại một điểm bằng tổng vectơ các cảm ứng từ thành phần.
\item (Sai) Trong dạng bài này, kết quả không phụ thuộc dữ kiện và có thể bỏ qua điều kiện áp dụng.
\item (Đúng) Khi giải bài thuộc dạng này phải xác định đúng phương, chiều, điều kiện áp dụng và đơn vị của các đại lượng.
\end{itemchoice}
}
\end{ex}

% ===== Câu 187 =====
% ID: L12C3B14-52-TL
% Mức: VD
\dangbt{Xác định chiều từ trường bằng kim nam châm và quy tắc nắm tay phải}
\begin{bt}
% ID: L12C3B14-52-TL
% Mức: VD
Đề xuất các bước giải một bài toán thuộc dạng xác định chiều từ trường bằng kim nam châm và quy tắc nắm tay phải.
\loigiai{
Cần nêu được: Đầu Bắc của kim nam châm nhỏ chỉ theo chiều của vectơ cảm ứng từ; quy tắc nắm tay phải xác định chiều đường sức quanh dòng điện. Khi vận dụng phải xác định đúng dữ kiện, phương chiều nếu có, điều kiện áp dụng, hệ đơn vị và kiểm tra tính hợp lí của kết quả. Nhận định “Quy tắc bàn tay trái dùng để xác định chiều đường sức từ quanh dây dẫn thẳng.” sai vì trái với kiến thức nêu trên.
}
\end{bt}

% ===== Câu 188 =====
% ID: L12C3B14-52-TLN
% Mức: VD
\begin{ex}
% ID: L12C3B14-52-TLN
% Mức: VD
Trong bốn phát biểu sau về xác định chiều từ trường bằng kim nam châm và quy tắc nắm tay phải, có bao nhiêu phát biểu đúng? (Chỉ ghi một số nguyên.) (1) Đầu Bắc của kim nam châm nhỏ chỉ theo chiều của vectơ cảm ứng từ; quy tắc nắm tay phải xác định chiều đường sức quanh dòng điện. (2) Cần đối chiếu kết quả với điều kiện bài toán. (3) Quy tắc bàn tay trái dùng để xác định chiều đường sức từ quanh dây dẫn thẳng. (4) Mọi trường hợp đều cho cùng một kết quả.
\shortans{2}
\loigiai{
Đối chiếu từng phát biểu với kiến thức: Đầu Bắc của kim nam châm nhỏ chỉ theo chiều của vectơ cảm ứng từ; quy tắc nắm tay phải xác định chiều đường sức quanh dòng điện. Các phát biểu đúng có số lượng là 2.
}
\end{ex}

% ===== Câu 189 =====
% ID: L12C3B14-52-TN
% Mức: VD
\dangbt{Từ trường của dòng điện thẳng, dòng điện tròn và ống dây}
\begin{ex}
% ID: L12C3B14-52-TN
% Mức: VD
Kết luận đúng khi vận dụng từ trường của dòng điện thẳng, dòng điện tròn và ống dây là
\choice
{Từ trường trong lòng ống dây dài luôn bằng không.}
{\True Đường sức quanh dây dẫn thẳng là các đường tròn đồng tâm; từ trường trong lòng ống dây dài gần đều.}
{Có thể bỏ qua phương, chiều và đơn vị của các đại lượng trong mọi trường hợp.}
{Đại lượng đang xét không phụ thuộc các điều kiện vật lí nêu trong bài.}
\loigiai{
Đường sức quanh dây dẫn thẳng là các đường tròn đồng tâm; từ trường trong lòng ống dây dài gần đều. Vì vậy phương án $B$ đúng; các phương án còn lại trái với định nghĩa, định luật hoặc điều kiện áp dụng.
}
\end{ex}

% ===== Câu 190 =====
% ID: L12C3B14-53-DS
% Mức: VD
\dangbt{So sánh, tổng hợp và biểu diễn các từ trường đơn giản}
\begin{ex}
% ID: L12C3B14-53-DS
% Mức: VD
Xét so sánh, tổng hợp và biểu diễn các từ trường đơn giản (Tình huống 3). Hãy xác định tính đúng, sai của bốn phát biểu sau.
\choiceTF
{\True Khi giải bài thuộc dạng này phải xác định đúng phương, chiều, điều kiện áp dụng và đơn vị của các đại lượng.}
{Trong dạng bài này, kết quả không phụ thuộc dữ kiện và có thể bỏ qua điều kiện áp dụng.}
{\True Vectơ cảm ứng từ tổng hợp tại một điểm bằng tổng vectơ các cảm ứng từ thành phần.}
{Độ lớn cảm ứng từ tổng hợp luôn bằng tổng độ lớn các cảm ứng từ thành phần.}
\loigiai{
\begin{itemchoice}
\item (Đúng) Khi giải bài thuộc dạng này phải xác định đúng phương, chiều, điều kiện áp dụng và đơn vị của các đại lượng. (Vì): Khi giải bài thuộc dạng này phải xác định đúng phương, chiều, điều kiện áp dụng và đơn vị của các đại lượng. b) Sai: Trong dạng bài này, kết quả không phụ thuộc dữ kiện và có thể bỏ qua điều kiện áp dụng. c) Đúng: Vectơ cảm ứng từ tổng hợp tại một điểm bằng tổng vectơ các cảm ứng từ thành phần. d) Sai: Độ lớn cảm ứng từ tổng hợp luôn bằng tổng độ lớn các cảm ứng từ thành phần. Kiến thức cốt lõi: Vectơ cảm ứng từ tổng hợp tại một điểm bằng tổng vectơ các cảm ứng từ thành phần
\item (Sai) Trong dạng bài này, kết quả không phụ thuộc dữ kiện và có thể bỏ qua điều kiện áp dụng.
\item (Đúng) Vectơ cảm ứng từ tổng hợp tại một điểm bằng tổng vectơ các cảm ứng từ thành phần.
\item (Sai) Độ lớn cảm ứng từ tổng hợp luôn bằng tổng độ lớn các cảm ứng từ thành phần.
\end{itemchoice}
}
\end{ex}

% ===== Câu 191 =====
% ID: L12C3B14-53-TL
% Mức: VD
\dangbt{Xác định chiều từ trường bằng kim nam châm và quy tắc nắm tay phải}
\begin{bt}
% ID: L12C3B14-53-TL
% Mức: VD
Nêu một tình huống thực tiễn liên quan đến xác định chiều từ trường bằng kim nam châm và quy tắc nắm tay phải và giải thích bằng kiến thức Vật lí.
\loigiai{
Cần nêu được: Đầu Bắc của kim nam châm nhỏ chỉ theo chiều của vectơ cảm ứng từ; quy tắc nắm tay phải xác định chiều đường sức quanh dòng điện. Khi vận dụng phải xác định đúng dữ kiện, phương chiều nếu có, điều kiện áp dụng, hệ đơn vị và kiểm tra tính hợp lí của kết quả. Nhận định “Quy tắc bàn tay trái dùng để xác định chiều đường sức từ quanh dây dẫn thẳng.” sai vì trái với kiến thức nêu trên.
}
\end{bt}

% ===== Câu 192 =====
% ID: L12C3B14-53-TLN
% Mức: VD
\begin{ex}
% ID: L12C3B14-53-TLN
% Mức: VD
Trong bốn phát biểu sau về xác định chiều từ trường bằng kim nam châm và quy tắc nắm tay phải, có bao nhiêu phát biểu đúng? (Chỉ ghi một số nguyên.) (1) Quy tắc bàn tay trái dùng để xác định chiều đường sức từ quanh dây dẫn thẳng. (2) Cần đọc đúng dữ kiện và giả thiết. (3) Đầu Bắc của kim nam châm nhỏ chỉ theo chiều của vectơ cảm ứng từ; quy tắc nắm tay phải xác định chiều đường sức quanh dòng điện. (4) Có thể thay số trước khi xác định công thức.
\shortans{2}
\loigiai{
Đối chiếu từng phát biểu với kiến thức: Đầu Bắc của kim nam châm nhỏ chỉ theo chiều của vectơ cảm ứng từ; quy tắc nắm tay phải xác định chiều đường sức quanh dòng điện. Các phát biểu đúng có số lượng là 2.
}
\end{ex}

% ===== Câu 193 =====
% ID: L12C3B14-53-TN
% Mức: NB
\begin{ex}
% ID: L12C3B14-53-TN
% Mức: NB
Phát biểu nào sau đây đúng về xác định chiều từ trường bằng kim nam châm và quy tắc nắm tay phải?
\choice
{\True Đầu Bắc của kim nam châm nhỏ chỉ theo chiều của vectơ cảm ứng từ; quy tắc nắm tay phải xác định chiều đường sức quanh dòng điện.}
{Quy tắc bàn tay trái dùng để xác định chiều đường sức từ quanh dây dẫn thẳng.}
{Đại lượng đang xét không phụ thuộc các điều kiện vật lí nêu trong bài.}
{Có thể bỏ qua phương, chiều và đơn vị của các đại lượng trong mọi trường hợp.}
\loigiai{
Đầu Bắc của kim nam châm nhỏ chỉ theo chiều của vectơ cảm ứng từ; quy tắc nắm tay phải xác định chiều đường sức quanh dòng điện. Vì vậy phương án $A$ đúng; các phương án còn lại trái với định nghĩa, định luật hoặc điều kiện áp dụng.
}
\end{ex}

% ===== Câu 194 =====
% ID: L12C3B14-54-DS
% Mức: VD
\dangbt{So sánh, tổng hợp và biểu diễn các từ trường đơn giản}
\begin{ex}
% ID: L12C3B14-54-DS
% Mức: VD
Xét so sánh, tổng hợp và biểu diễn các từ trường đơn giản (Tình huống 4). Hãy xác định tính đúng, sai của bốn phát biểu sau.
\choiceTF
{Trong dạng bài này, kết quả không phụ thuộc dữ kiện và có thể bỏ qua điều kiện áp dụng.}
{\True Khi giải bài thuộc dạng này phải xác định đúng phương, chiều, điều kiện áp dụng và đơn vị của các đại lượng.}
{Độ lớn cảm ứng từ tổng hợp luôn bằng tổng độ lớn các cảm ứng từ thành phần.}
{\True Vectơ cảm ứng từ tổng hợp tại một điểm bằng tổng vectơ các cảm ứng từ thành phần.}
\loigiai{
\begin{itemchoice}
\item (Sai) Trong dạng bài này, kết quả không phụ thuộc dữ kiện và có thể bỏ qua điều kiện áp dụng. (Vì): Trong dạng bài này, kết quả không phụ thuộc dữ kiện và có thể bỏ qua điều kiện áp dụng. b) Đúng: Khi giải bài thuộc dạng này phải xác định đúng phương, chiều, điều kiện áp dụng và đơn vị của các đại lượng. c) Sai: Độ lớn cảm ứng từ tổng hợp luôn bằng tổng độ lớn các cảm ứng từ thành phần. d) Đúng: Vectơ cảm ứng từ tổng hợp tại một điểm bằng tổng vectơ các cảm ứng từ thành phần. Kiến thức cốt lõi: Vectơ cảm ứng từ tổng hợp tại một điểm bằng tổng vectơ các cảm ứng từ thành phần
\item (Đúng) Khi giải bài thuộc dạng này phải xác định đúng phương, chiều, điều kiện áp dụng và đơn vị của các đại lượng.
\item (Sai) Độ lớn cảm ứng từ tổng hợp luôn bằng tổng độ lớn các cảm ứng từ thành phần.
\item (Đúng) Vectơ cảm ứng từ tổng hợp tại một điểm bằng tổng vectơ các cảm ứng từ thành phần.
\end{itemchoice}
}
\end{ex}

% ===== Câu 195 =====
% ID: L12C3B14-54-TL
% Mức: VDC
\dangbt{Xác định chiều từ trường bằng kim nam châm và quy tắc nắm tay phải}
\begin{bt}
% ID: L12C3B14-54-TL
% Mức: VDC
So sánh nhận định đúng “Đầu Bắc của kim nam châm nhỏ chỉ theo chiều của vectơ cảm ứng từ; quy tắc nắm tay phải xác định chiều đường sức quanh dòng điện.” với nhận định sai “Quy tắc bàn tay trái dùng để xác định chiều đường sức từ quanh dây dẫn thẳng.”; chỉ rõ căn cứ Vật lí.
\loigiai{
Cần nêu được: Đầu Bắc của kim nam châm nhỏ chỉ theo chiều của vectơ cảm ứng từ; quy tắc nắm tay phải xác định chiều đường sức quanh dòng điện. Khi vận dụng phải xác định đúng dữ kiện, phương chiều nếu có, điều kiện áp dụng, hệ đơn vị và kiểm tra tính hợp lí của kết quả. Nhận định “Quy tắc bàn tay trái dùng để xác định chiều đường sức từ quanh dây dẫn thẳng.” sai vì trái với kiến thức nêu trên.
}
\end{bt}

% ===== Câu 196 =====
% ID: L12C3B14-54-TLN
% Mức: VDC
\begin{ex}
% ID: L12C3B14-54-TLN
% Mức: VDC
Trong bốn phát biểu sau về xác định chiều từ trường bằng kim nam châm và quy tắc nắm tay phải, có bao nhiêu phát biểu đúng? (Chỉ ghi một số nguyên.) (1) Kết quả cần có ý nghĩa vật lí. (2) Quy tắc bàn tay trái dùng để xác định chiều đường sức từ quanh dây dẫn thẳng. (3) Phải dùng đúng định luật cho tình huống. (4) Đầu Bắc của kim nam châm nhỏ chỉ theo chiều của vectơ cảm ứng từ; quy tắc nắm tay phải xác định chiều đường sức quanh dòng điện.
\shortans{3}
\loigiai{
Đối chiếu từng phát biểu với kiến thức: Đầu Bắc của kim nam châm nhỏ chỉ theo chiều của vectơ cảm ứng từ; quy tắc nắm tay phải xác định chiều đường sức quanh dòng điện. Các phát biểu đúng có số lượng là 3.
}
\end{ex}

% ===== Câu 197 =====
% ID: L12C3B14-54-TN
% Mức: NB
\begin{ex}
% ID: L12C3B14-54-TN
% Mức: NB
Trong xác định chiều từ trường bằng kim nam châm và quy tắc nắm tay phải, nhận định nào phù hợp với kiến thức Vật lí?
\choice
{Quy tắc bàn tay trái dùng để xác định chiều đường sức từ quanh dây dẫn thẳng.}
{\True Đầu Bắc của kim nam châm nhỏ chỉ theo chiều của vectơ cảm ứng từ; quy tắc nắm tay phải xác định chiều đường sức quanh dòng điện.}
{Có thể bỏ qua phương, chiều và đơn vị của các đại lượng trong mọi trường hợp.}
{Đại lượng đang xét không phụ thuộc các điều kiện vật lí nêu trong bài.}
\loigiai{
Đầu Bắc của kim nam châm nhỏ chỉ theo chiều của vectơ cảm ứng từ; quy tắc nắm tay phải xác định chiều đường sức quanh dòng điện. Vì vậy phương án $B$ đúng; các phương án còn lại trái với định nghĩa, định luật hoặc điều kiện áp dụng.
}
\end{ex}

% ===== Câu 198 =====
% ID: L12C3B14-55-DS
% Mức: VD
\dangbt{So sánh, tổng hợp và biểu diễn các từ trường đơn giản}
\begin{ex}
% ID: L12C3B14-55-DS
% Mức: VD
Xét so sánh, tổng hợp và biểu diễn các từ trường đơn giản (Tình huống 5). Hãy xác định tính đúng, sai của bốn phát biểu sau.
\choiceTF
{\True Vectơ cảm ứng từ tổng hợp tại một điểm bằng tổng vectơ các cảm ứng từ thành phần.}
{Trong dạng bài này, kết quả không phụ thuộc dữ kiện và có thể bỏ qua điều kiện áp dụng.}
{Độ lớn cảm ứng từ tổng hợp luôn bằng tổng độ lớn các cảm ứng từ thành phần.}
{\True Khi giải bài thuộc dạng này phải xác định đúng phương, chiều, điều kiện áp dụng và đơn vị của các đại lượng.}
\loigiai{
\begin{itemchoice}
\item (Đúng) Vectơ cảm ứng từ tổng hợp tại một điểm bằng tổng vectơ các cảm ứng từ thành phần. (Vì): Vectơ cảm ứng từ tổng hợp tại một điểm bằng tổng vectơ các cảm ứng từ thành phần. b) Sai: Trong dạng bài này, kết quả không phụ thuộc dữ kiện và có thể bỏ qua điều kiện áp dụng. c) Sai: Độ lớn cảm ứng từ tổng hợp luôn bằng tổng độ lớn các cảm ứng từ thành phần. d) Đúng: Khi giải bài thuộc dạng này phải xác định đúng phương, chiều, điều kiện áp dụng và đơn vị của các đại lượng. Kiến thức cốt lõi: Vectơ cảm ứng từ tổng hợp tại một điểm bằng tổng vectơ các cảm ứng từ thành phần
\item (Sai) Trong dạng bài này, kết quả không phụ thuộc dữ kiện và có thể bỏ qua điều kiện áp dụng.
\item (Sai) Độ lớn cảm ứng từ tổng hợp luôn bằng tổng độ lớn các cảm ứng từ thành phần.
\item (Đúng) Khi giải bài thuộc dạng này phải xác định đúng phương, chiều, điều kiện áp dụng và đơn vị của các đại lượng.
\end{itemchoice}
}
\end{ex}

% ===== Câu 199 =====
% ID: L12C3B14-55-TL
% Mức: TH
\dangbt{Từ trường của dòng điện thẳng, dòng điện tròn và ống dây}
\begin{bt}
% ID: L12C3B14-55-TL
% Mức: TH
Trình bày kiến thức cốt lõi của từ trường của dòng điện thẳng, dòng điện tròn và ống dây và nêu điều kiện áp dụng.
\loigiai{
Cần nêu được: Đường sức quanh dây dẫn thẳng là các đường tròn đồng tâm; từ trường trong lòng ống dây dài gần đều. Khi vận dụng phải xác định đúng dữ kiện, phương chiều nếu có, điều kiện áp dụng, hệ đơn vị và kiểm tra tính hợp lí của kết quả. Nhận định “Từ trường trong lòng ống dây dài luôn bằng không.” sai vì trái với kiến thức nêu trên.
}
\end{bt}

% ===== Câu 200 =====
% ID: L12C3B14-55-TLN
% Mức: TH
\begin{ex}
% ID: L12C3B14-55-TLN
% Mức: TH
Trong bốn phát biểu sau về từ trường của dòng điện thẳng, dòng điện tròn và ống dây, có bao nhiêu phát biểu đúng? (Chỉ ghi một số nguyên.) (1) Đường sức quanh dây dẫn thẳng là các đường tròn đồng tâm; từ trường trong lòng ống dây dài gần đều. (2) Từ trường trong lòng ống dây dài luôn bằng không. (3) Cần kiểm tra điều kiện áp dụng trước khi tính toán. (4) Có thể bỏ qua đơn vị của kết quả.
\shortans{2}
\loigiai{
Đối chiếu từng phát biểu với kiến thức: Đường sức quanh dây dẫn thẳng là các đường tròn đồng tâm; từ trường trong lòng ống dây dài gần đều. Các phát biểu đúng có số lượng là 2.
}
\end{ex}

% ===== Câu 201 =====
% ID: L12C3B14-55-TN
% Mức: NB
\dangbt{Xác định chiều từ trường bằng kim nam châm và quy tắc nắm tay phải}
\begin{ex}
% ID: L12C3B14-55-TN
% Mức: NB
Tình huống 3: chọn kết luận chính xác về xác định chiều từ trường bằng kim nam châm và quy tắc nắm tay phải.
\choice
{Đại lượng đang xét không phụ thuộc các điều kiện vật lí nêu trong bài.}
{Có thể bỏ qua phương, chiều và đơn vị của các đại lượng trong mọi trường hợp.}
{\True Đầu Bắc của kim nam châm nhỏ chỉ theo chiều của vectơ cảm ứng từ; quy tắc nắm tay phải xác định chiều đường sức quanh dòng điện.}
{Quy tắc bàn tay trái dùng để xác định chiều đường sức từ quanh dây dẫn thẳng.}
\loigiai{
Đầu Bắc của kim nam châm nhỏ chỉ theo chiều của vectơ cảm ứng từ; quy tắc nắm tay phải xác định chiều đường sức quanh dòng điện. Vì vậy phương án $C$ đúng; các phương án còn lại trái với định nghĩa, định luật hoặc điều kiện áp dụng.
}
\end{ex}

% ===== Câu 202 =====
% ID: L12C3B14-56-DS
% Mức: TH
\dangbt{Giải thích hiện tượng và ứng dụng của từ trường trong thực tiễn}
\begin{ex}
% ID: L12C3B14-56-DS
% Mức: TH
Xét giải thích hiện tượng và ứng dụng của từ trường trong thực tiễn (Tình huống 1). Hãy xác định tính đúng, sai của bốn phát biểu sau.
\choiceTF
{\True La bàn định hướng nhờ kim nam châm tương tác với từ trường Trái Đất.}
{La bàn hoạt động chủ yếu nhờ điện trường của Trái Đất.}
{\True Khi giải bài thuộc dạng này phải xác định đúng phương, chiều, điều kiện áp dụng và đơn vị của các đại lượng.}
{Trong dạng bài này, kết quả không phụ thuộc dữ kiện và có thể bỏ qua điều kiện áp dụng.}
\loigiai{
\begin{itemchoice}
\item (Đúng) La bàn định hướng nhờ kim nam châm tương tác với từ trường Trái Đất. (Vì): La bàn định hướng nhờ kim nam châm tương tác với từ trường Trái Đất. b) Sai: La bàn hoạt động chủ yếu nhờ điện trường của Trái Đất. c) Đúng: Khi giải bài thuộc dạng này phải xác định đúng phương, chiều, điều kiện áp dụng và đơn vị của các đại lượng. d) Sai: Trong dạng bài này, kết quả không phụ thuộc dữ kiện và có thể bỏ qua điều kiện áp dụng. Kiến thức cốt lõi: La bàn định hướng nhờ kim nam châm tương tác với từ trường Trái Đất
\item (Sai) La bàn hoạt động chủ yếu nhờ điện trường của Trái Đất.
\item (Đúng) Khi giải bài thuộc dạng này phải xác định đúng phương, chiều, điều kiện áp dụng và đơn vị của các đại lượng.
\item (Sai) Trong dạng bài này, kết quả không phụ thuộc dữ kiện và có thể bỏ qua điều kiện áp dụng.
\end{itemchoice}
}
\end{ex}

% ===== Câu 203 =====
% ID: L12C3B14-56-TL
% Mức: TH
\dangbt{Từ trường của dòng điện thẳng, dòng điện tròn và ống dây}
\begin{bt}
% ID: L12C3B14-56-TL
% Mức: TH
Giải thích vì sao nhận định sau đúng: “Đường sức quanh dây dẫn thẳng là các đường tròn đồng tâm; từ trường trong lòng ống dây dài gần đều.”
\loigiai{
Cần nêu được: Đường sức quanh dây dẫn thẳng là các đường tròn đồng tâm; từ trường trong lòng ống dây dài gần đều. Khi vận dụng phải xác định đúng dữ kiện, phương chiều nếu có, điều kiện áp dụng, hệ đơn vị và kiểm tra tính hợp lí của kết quả. Nhận định “Từ trường trong lòng ống dây dài luôn bằng không.” sai vì trái với kiến thức nêu trên.
}
\end{bt}

% ===== Câu 204 =====
% ID: L12C3B14-56-TLN
% Mức: TH
\begin{ex}
% ID: L12C3B14-56-TLN
% Mức: TH
Trong bốn phát biểu sau về từ trường của dòng điện thẳng, dòng điện tròn và ống dây, có bao nhiêu phát biểu đúng? (Chỉ ghi một số nguyên.) (1) Từ trường trong lòng ống dây dài luôn bằng không. (2) Đường sức quanh dây dẫn thẳng là các đường tròn đồng tâm; từ trường trong lòng ống dây dài gần đều. (3) Kết quả phải phù hợp ý nghĩa vật lí. (4) Mọi đại lượng đều là vô hướng.
\shortans{1}
\loigiai{
Đối chiếu từng phát biểu với kiến thức: Đường sức quanh dây dẫn thẳng là các đường tròn đồng tâm; từ trường trong lòng ống dây dài gần đều. Các phát biểu đúng có số lượng là 1.
}
\end{ex}

% ===== Câu 205 =====
% ID: L12C3B14-56-TN
% Mức: NB
\dangbt{Xác định chiều từ trường bằng kim nam châm và quy tắc nắm tay phải}
\begin{ex}
% ID: L12C3B14-56-TN
% Mức: NB
Một dây dẫn thẳng dài vô hạn có phương vuông góc với mặt phẳng trang giấy. Cho dòng điện chạy qua dây dẫn theo chiều từ trong ra ngoài. Hình nào dưới đây mô tả đúng đường sức từ trên mặt phẳng trang giấy của từ trường của dòng điện chạy trong dây dẫn?
\choice
{Hình 2}
{\True Hình 3}
{Hình 4}
{Hình 1}
\loigiai{
:Sử dụng quy tắc bàn tay phải về tìm chiều của dòng điện chạy trong dây dẫn thẳng dài vô hạn đặt trong từ trường “Giơ ngón cái của bàn tay phải hướng theo chiều dòng điện, khum bốn ngón kia xung quanh dây dẫn thì chiều từ cổ tay đến các ngón tay là chiều của đường sức từ”.
}
\end{ex}

% ===== Câu 206 =====
% ID: L12C3B14-57-DS
% Mức: TH
\dangbt{Giải thích hiện tượng và ứng dụng của từ trường trong thực tiễn}
\begin{ex}
% ID: L12C3B14-57-DS
% Mức: TH
Xét giải thích hiện tượng và ứng dụng của từ trường trong thực tiễn (Tình huống 2). Hãy xác định tính đúng, sai của bốn phát biểu sau.
\choiceTF
{La bàn hoạt động chủ yếu nhờ điện trường của Trái Đất.}
{\True La bàn định hướng nhờ kim nam châm tương tác với từ trường Trái Đất.}
{Trong dạng bài này, kết quả không phụ thuộc dữ kiện và có thể bỏ qua điều kiện áp dụng.}
{\True Khi giải bài thuộc dạng này phải xác định đúng phương, chiều, điều kiện áp dụng và đơn vị của các đại lượng.}
\loigiai{
\begin{itemchoice}
\item (Sai) La bàn hoạt động chủ yếu nhờ điện trường của Trái Đất. (Vì): La bàn hoạt động chủ yếu nhờ điện trường của Trái Đất. b) Đúng: La bàn định hướng nhờ kim nam châm tương tác với từ trường Trái Đất. c) Sai: Trong dạng bài này, kết quả không phụ thuộc dữ kiện và có thể bỏ qua điều kiện áp dụng. d) Đúng: Khi giải bài thuộc dạng này phải xác định đúng phương, chiều, điều kiện áp dụng và đơn vị của các đại lượng. Kiến thức cốt lõi: La bàn định hướng nhờ kim nam châm tương tác với từ trường Trái Đất
\item (Đúng) La bàn định hướng nhờ kim nam châm tương tác với từ trường Trái Đất.
\item (Sai) Trong dạng bài này, kết quả không phụ thuộc dữ kiện và có thể bỏ qua điều kiện áp dụng.
\item (Đúng) Khi giải bài thuộc dạng này phải xác định đúng phương, chiều, điều kiện áp dụng và đơn vị của các đại lượng.
\end{itemchoice}
}
\end{ex}

% ===== Câu 207 =====
% ID: L12C3B14-57-TL
% Mức: VD
\dangbt{Từ trường của dòng điện thẳng, dòng điện tròn và ống dây}
\begin{bt}
% ID: L12C3B14-57-TL
% Mức: VD
Phân tích sai lầm trong nhận định: “Từ trường trong lòng ống dây dài luôn bằng không.”
\loigiai{
Cần nêu được: Đường sức quanh dây dẫn thẳng là các đường tròn đồng tâm; từ trường trong lòng ống dây dài gần đều. Khi vận dụng phải xác định đúng dữ kiện, phương chiều nếu có, điều kiện áp dụng, hệ đơn vị và kiểm tra tính hợp lí của kết quả. Nhận định “Từ trường trong lòng ống dây dài luôn bằng không.” sai vì trái với kiến thức nêu trên.
}
\end{bt}

% ===== Câu 208 =====
% ID: L12C3B14-57-TLN
% Mức: TH
\begin{ex}
% ID: L12C3B14-57-TLN
% Mức: TH
Trong bốn phát biểu sau về từ trường của dòng điện thẳng, dòng điện tròn và ống dây, có bao nhiêu phát biểu đúng? (Chỉ ghi một số nguyên.) (1) Cần đổi các đại lượng về hệ đơn vị thống nhất. (2) Đường sức quanh dây dẫn thẳng là các đường tròn đồng tâm; từ trường trong lòng ống dây dài gần đều. (3) Từ trường trong lòng ống dây dài luôn bằng không. (4) Không cần xét chiều của đại lượng vectơ.
\shortans{2}
\loigiai{
Đối chiếu từng phát biểu với kiến thức: Đường sức quanh dây dẫn thẳng là các đường tròn đồng tâm; từ trường trong lòng ống dây dài gần đều. Các phát biểu đúng có số lượng là 2.
}
\end{ex}

% ===== Câu 209 =====
% ID: L12C3B14-57-TN
% Mức: TH
\dangbt{Xác định chiều từ trường bằng kim nam châm và quy tắc nắm tay phải}
\begin{ex}
% ID: L12C3B14-57-TN
% Mức: TH
Nội dung nào mô tả đúng nhất xác định chiều từ trường bằng kim nam châm và quy tắc nắm tay phải?
\choice
{Có thể bỏ qua phương, chiều và đơn vị của các đại lượng trong mọi trường hợp.}
{Đại lượng đang xét không phụ thuộc các điều kiện vật lí nêu trong bài.}
{Quy tắc bàn tay trái dùng để xác định chiều đường sức từ quanh dây dẫn thẳng.}
{\True Đầu Bắc của kim nam châm nhỏ chỉ theo chiều của vectơ cảm ứng từ; quy tắc nắm tay phải xác định chiều đường sức quanh dòng điện.}
\loigiai{
Đầu Bắc của kim nam châm nhỏ chỉ theo chiều của vectơ cảm ứng từ; quy tắc nắm tay phải xác định chiều đường sức quanh dòng điện. Vì vậy phương án $D$ đúng; các phương án còn lại trái với định nghĩa, định luật hoặc điều kiện áp dụng.
}
\end{ex}

% ===== Câu 210 =====
% ID: L12C3B14-58-DS
% Mức: VD
\dangbt{Giải thích hiện tượng và ứng dụng của từ trường trong thực tiễn}
\begin{ex}
% ID: L12C3B14-58-DS
% Mức: VD
Xét giải thích hiện tượng và ứng dụng của từ trường trong thực tiễn (Tình huống 3). Hãy xác định tính đúng, sai của bốn phát biểu sau.
\choiceTF
{\True Khi giải bài thuộc dạng này phải xác định đúng phương, chiều, điều kiện áp dụng và đơn vị của các đại lượng.}
{Trong dạng bài này, kết quả không phụ thuộc dữ kiện và có thể bỏ qua điều kiện áp dụng.}
{\True La bàn định hướng nhờ kim nam châm tương tác với từ trường Trái Đất.}
{La bàn hoạt động chủ yếu nhờ điện trường của Trái Đất.}
\loigiai{
\begin{itemchoice}
\item (Đúng) Khi giải bài thuộc dạng này phải xác định đúng phương, chiều, điều kiện áp dụng và đơn vị của các đại lượng. (Vì): Khi giải bài thuộc dạng này phải xác định đúng phương, chiều, điều kiện áp dụng và đơn vị của các đại lượng. b) Sai: Trong dạng bài này, kết quả không phụ thuộc dữ kiện và có thể bỏ qua điều kiện áp dụng. c) Đúng: La bàn định hướng nhờ kim nam châm tương tác với từ trường Trái Đất. d) Sai: La bàn hoạt động chủ yếu nhờ điện trường của Trái Đất. Kiến thức cốt lõi: La bàn định hướng nhờ kim nam châm tương tác với từ trường Trái Đất
\item (Sai) Trong dạng bài này, kết quả không phụ thuộc dữ kiện và có thể bỏ qua điều kiện áp dụng.
\item (Đúng) La bàn định hướng nhờ kim nam châm tương tác với từ trường Trái Đất.
\item (Sai) La bàn hoạt động chủ yếu nhờ điện trường của Trái Đất.
\end{itemchoice}
}
\end{ex}

% ===== Câu 211 =====
% ID: L12C3B14-58-TL
% Mức: VD
\dangbt{Từ trường của dòng điện thẳng, dòng điện tròn và ống dây}
\begin{bt}
% ID: L12C3B14-58-TL
% Mức: VD
Đề xuất các bước giải một bài toán thuộc dạng từ trường của dòng điện thẳng, dòng điện tròn và ống dây.
\loigiai{
Cần nêu được: Đường sức quanh dây dẫn thẳng là các đường tròn đồng tâm; từ trường trong lòng ống dây dài gần đều. Khi vận dụng phải xác định đúng dữ kiện, phương chiều nếu có, điều kiện áp dụng, hệ đơn vị và kiểm tra tính hợp lí của kết quả. Nhận định “Từ trường trong lòng ống dây dài luôn bằng không.” sai vì trái với kiến thức nêu trên.
}
\end{bt}

% ===== Câu 212 =====
% ID: L12C3B14-58-TLN
% Mức: VD
\begin{ex}
% ID: L12C3B14-58-TLN
% Mức: VD
Trong bốn phát biểu sau về từ trường của dòng điện thẳng, dòng điện tròn và ống dây, có bao nhiêu phát biểu đúng? (Chỉ ghi một số nguyên.) (1) Đường sức quanh dây dẫn thẳng là các đường tròn đồng tâm; từ trường trong lòng ống dây dài gần đều. (2) Cần đối chiếu kết quả với điều kiện bài toán. (3) Từ trường trong lòng ống dây dài luôn bằng không. (4) Mọi trường hợp đều cho cùng một kết quả.
\shortans{2}
\loigiai{
Đối chiếu từng phát biểu với kiến thức: Đường sức quanh dây dẫn thẳng là các đường tròn đồng tâm; từ trường trong lòng ống dây dài gần đều. Các phát biểu đúng có số lượng là 2.
}
\end{ex}

% ===== Câu 213 =====
% ID: L12C3B14-58-TN
% Mức: TH
\dangbt{Xác định chiều từ trường bằng kim nam châm và quy tắc nắm tay phải}
\begin{ex}
% ID: L12C3B14-58-TN
% Mức: TH
Khi phân tích xác định chiều từ trường bằng kim nam châm và quy tắc nắm tay phải, phát biểu nào đúng?
\choice
{\True Đầu Bắc của kim nam châm nhỏ chỉ theo chiều của vectơ cảm ứng từ; quy tắc nắm tay phải xác định chiều đường sức quanh dòng điện.}
{Quy tắc bàn tay trái dùng để xác định chiều đường sức từ quanh dây dẫn thẳng.}
{Đại lượng đang xét không phụ thuộc các điều kiện vật lí nêu trong bài.}
{Có thể bỏ qua phương, chiều và đơn vị của các đại lượng trong mọi trường hợp.}
\loigiai{
Đầu Bắc của kim nam châm nhỏ chỉ theo chiều của vectơ cảm ứng từ; quy tắc nắm tay phải xác định chiều đường sức quanh dòng điện. Vì vậy phương án $A$ đúng; các phương án còn lại trái với định nghĩa, định luật hoặc điều kiện áp dụng.
}
\end{ex}

% ===== Câu 214 =====
% ID: L12C3B14-59-DS
% Mức: VD
\dangbt{Giải thích hiện tượng và ứng dụng của từ trường trong thực tiễn}
\begin{ex}
% ID: L12C3B14-59-DS
% Mức: VD
Xét giải thích hiện tượng và ứng dụng của từ trường trong thực tiễn (Tình huống 4). Hãy xác định tính đúng, sai của bốn phát biểu sau.
\choiceTF
{Trong dạng bài này, kết quả không phụ thuộc dữ kiện và có thể bỏ qua điều kiện áp dụng.}
{\True Khi giải bài thuộc dạng này phải xác định đúng phương, chiều, điều kiện áp dụng và đơn vị của các đại lượng.}
{La bàn hoạt động chủ yếu nhờ điện trường của Trái Đất.}
{\True La bàn định hướng nhờ kim nam châm tương tác với từ trường Trái Đất.}
\loigiai{
\begin{itemchoice}
\item (Sai) Trong dạng bài này, kết quả không phụ thuộc dữ kiện và có thể bỏ qua điều kiện áp dụng. (Vì): Trong dạng bài này, kết quả không phụ thuộc dữ kiện và có thể bỏ qua điều kiện áp dụng. b) Đúng: Khi giải bài thuộc dạng này phải xác định đúng phương, chiều, điều kiện áp dụng và đơn vị của các đại lượng. c) Sai: La bàn hoạt động chủ yếu nhờ điện trường của Trái Đất. d) Đúng: La bàn định hướng nhờ kim nam châm tương tác với từ trường Trái Đất. Kiến thức cốt lõi: La bàn định hướng nhờ kim nam châm tương tác với từ trường Trái Đất
\item (Đúng) Khi giải bài thuộc dạng này phải xác định đúng phương, chiều, điều kiện áp dụng và đơn vị của các đại lượng.
\item (Sai) La bàn hoạt động chủ yếu nhờ điện trường của Trái Đất.
\item (Đúng) La bàn định hướng nhờ kim nam châm tương tác với từ trường Trái Đất.
\end{itemchoice}
}
\end{ex}

% ===== Câu 215 =====
% ID: L12C3B14-59-TL
% Mức: VD
\dangbt{Từ trường của dòng điện thẳng, dòng điện tròn và ống dây}
\begin{bt}
% ID: L12C3B14-59-TL
% Mức: VD
Nêu một tình huống thực tiễn liên quan đến từ trường của dòng điện thẳng, dòng điện tròn và ống dây và giải thích bằng kiến thức Vật lí.
\loigiai{
Cần nêu được: Đường sức quanh dây dẫn thẳng là các đường tròn đồng tâm; từ trường trong lòng ống dây dài gần đều. Khi vận dụng phải xác định đúng dữ kiện, phương chiều nếu có, điều kiện áp dụng, hệ đơn vị và kiểm tra tính hợp lí của kết quả. Nhận định “Từ trường trong lòng ống dây dài luôn bằng không.” sai vì trái với kiến thức nêu trên.
}
\end{bt}

% ===== Câu 216 =====
% ID: L12C3B14-59-TLN
% Mức: VD
\begin{ex}
% ID: L12C3B14-59-TLN
% Mức: VD
Trong bốn phát biểu sau về từ trường của dòng điện thẳng, dòng điện tròn và ống dây, có bao nhiêu phát biểu đúng? (Chỉ ghi một số nguyên.) (1) Từ trường trong lòng ống dây dài luôn bằng không. (2) Cần đọc đúng dữ kiện và giả thiết. (3) Đường sức quanh dây dẫn thẳng là các đường tròn đồng tâm; từ trường trong lòng ống dây dài gần đều. (4) Có thể thay số trước khi xác định công thức.
\shortans{2}
\loigiai{
Đối chiếu từng phát biểu với kiến thức: Đường sức quanh dây dẫn thẳng là các đường tròn đồng tâm; từ trường trong lòng ống dây dài gần đều. Các phát biểu đúng có số lượng là 2.
}
\end{ex}

% ===== Câu 217 =====
% ID: L12C3B14-59-TN
% Mức: TH
\dangbt{Xác định chiều từ trường bằng kim nam châm và quy tắc nắm tay phải}
\begin{ex}
% ID: L12C3B14-59-TN
% Mức: TH
Một học sinh ôn tập xác định chiều từ trường bằng kim nam châm và quy tắc nắm tay phải. Kết luận nào của học sinh là đúng?
\choice
{Quy tắc bàn tay trái dùng để xác định chiều đường sức từ quanh dây dẫn thẳng.}
{\True Đầu Bắc của kim nam châm nhỏ chỉ theo chiều của vectơ cảm ứng từ; quy tắc nắm tay phải xác định chiều đường sức quanh dòng điện.}
{Có thể bỏ qua phương, chiều và đơn vị của các đại lượng trong mọi trường hợp.}
{Đại lượng đang xét không phụ thuộc các điều kiện vật lí nêu trong bài.}
\loigiai{
Đầu Bắc của kim nam châm nhỏ chỉ theo chiều của vectơ cảm ứng từ; quy tắc nắm tay phải xác định chiều đường sức quanh dòng điện. Vì vậy phương án $B$ đúng; các phương án còn lại trái với định nghĩa, định luật hoặc điều kiện áp dụng.
}
\end{ex}

% ===== Câu 218 =====
% ID: L12C3B14-60-DS
% Mức: VD
\dangbt{Giải thích hiện tượng và ứng dụng của từ trường trong thực tiễn}
\begin{ex}
% ID: L12C3B14-60-DS
% Mức: VD
Xét giải thích hiện tượng và ứng dụng của từ trường trong thực tiễn (Tình huống 5). Hãy xác định tính đúng, sai của bốn phát biểu sau.
\choiceTF
{\True La bàn định hướng nhờ kim nam châm tương tác với từ trường Trái Đất.}
{Trong dạng bài này, kết quả không phụ thuộc dữ kiện và có thể bỏ qua điều kiện áp dụng.}
{La bàn hoạt động chủ yếu nhờ điện trường của Trái Đất.}
{\True Khi giải bài thuộc dạng này phải xác định đúng phương, chiều, điều kiện áp dụng và đơn vị của các đại lượng.}
\loigiai{
\begin{itemchoice}
\item (Đúng) La bàn định hướng nhờ kim nam châm tương tác với từ trường Trái Đất. (Vì): La bàn định hướng nhờ kim nam châm tương tác với từ trường Trái Đất. b) Sai: Trong dạng bài này, kết quả không phụ thuộc dữ kiện và có thể bỏ qua điều kiện áp dụng. c) Sai: La bàn hoạt động chủ yếu nhờ điện trường của Trái Đất. d) Đúng: Khi giải bài thuộc dạng này phải xác định đúng phương, chiều, điều kiện áp dụng và đơn vị của các đại lượng. Kiến thức cốt lõi: La bàn định hướng nhờ kim nam châm tương tác với từ trường Trái Đất
\item (Sai) Trong dạng bài này, kết quả không phụ thuộc dữ kiện và có thể bỏ qua điều kiện áp dụng.
\item (Sai) La bàn hoạt động chủ yếu nhờ điện trường của Trái Đất.
\item (Đúng) Khi giải bài thuộc dạng này phải xác định đúng phương, chiều, điều kiện áp dụng và đơn vị của các đại lượng.
\end{itemchoice}
}
\end{ex}

% ===== Câu 219 =====
% ID: L12C3B14-60-TL
% Mức: VDC
\dangbt{Từ trường của dòng điện thẳng, dòng điện tròn và ống dây}
\begin{bt}
% ID: L12C3B14-60-TL
% Mức: VDC
So sánh nhận định đúng “Đường sức quanh dây dẫn thẳng là các đường tròn đồng tâm; từ trường trong lòng ống dây dài gần đều.” với nhận định sai “Từ trường trong lòng ống dây dài luôn bằng không.”; chỉ rõ căn cứ Vật lí.
\loigiai{
Cần nêu được: Đường sức quanh dây dẫn thẳng là các đường tròn đồng tâm; từ trường trong lòng ống dây dài gần đều. Khi vận dụng phải xác định đúng dữ kiện, phương chiều nếu có, điều kiện áp dụng, hệ đơn vị và kiểm tra tính hợp lí của kết quả. Nhận định “Từ trường trong lòng ống dây dài luôn bằng không.” sai vì trái với kiến thức nêu trên.
}
\end{bt}

% ===== Câu 220 =====
% ID: L12C3B14-60-TLN
% Mức: VDC
\begin{ex}
% ID: L12C3B14-60-TLN
% Mức: VDC
Trong bốn phát biểu sau về từ trường của dòng điện thẳng, dòng điện tròn và ống dây, có bao nhiêu phát biểu đúng? (Chỉ ghi một số nguyên.) (1) Kết quả cần có ý nghĩa vật lí. (2) Từ trường trong lòng ống dây dài luôn bằng không. (3) Phải dùng đúng định luật cho tình huống. (4) Đường sức quanh dây dẫn thẳng là các đường tròn đồng tâm; từ trường trong lòng ống dây dài gần đều.
\shortans{3}
\loigiai{
Đối chiếu từng phát biểu với kiến thức: Đường sức quanh dây dẫn thẳng là các đường tròn đồng tâm; từ trường trong lòng ống dây dài gần đều. Các phát biểu đúng có số lượng là 3.
}
\end{ex}

% ===== Câu 221 =====
% ID: L12C3B14-60-TN
% Mức: TH
\dangbt{Xác định chiều từ trường bằng kim nam châm và quy tắc nắm tay phải}
\begin{ex}
% ID: L12C3B14-60-TN
% Mức: TH
Chọn phát biểu không trái với kiến thức về xác định chiều từ trường bằng kim nam châm và quy tắc nắm tay phải.
\choice
{Đại lượng đang xét không phụ thuộc các điều kiện vật lí nêu trong bài.}
{Có thể bỏ qua phương, chiều và đơn vị của các đại lượng trong mọi trường hợp.}
{\True Đầu Bắc của kim nam châm nhỏ chỉ theo chiều của vectơ cảm ứng từ; quy tắc nắm tay phải xác định chiều đường sức quanh dòng điện.}
{Quy tắc bàn tay trái dùng để xác định chiều đường sức từ quanh dây dẫn thẳng.}
\loigiai{
Đầu Bắc của kim nam châm nhỏ chỉ theo chiều của vectơ cảm ứng từ; quy tắc nắm tay phải xác định chiều đường sức quanh dòng điện. Vì vậy phương án $C$ đúng; các phương án còn lại trái với định nghĩa, định luật hoặc điều kiện áp dụng.
}
\end{ex}

% ===== Câu 222 =====
% ID: L12C3B14-61-TL
% Mức: TH
\dangbt{So sánh, tổng hợp và biểu diễn các từ trường đơn giản}
\begin{bt}
% ID: L12C3B14-61-TL
% Mức: TH
Trình bày kiến thức cốt lõi của so sánh, tổng hợp và biểu diễn các từ trường đơn giản và nêu điều kiện áp dụng.
\loigiai{
Cần nêu được: Vectơ cảm ứng từ tổng hợp tại một điểm bằng tổng vectơ các cảm ứng từ thành phần. Khi vận dụng phải xác định đúng dữ kiện, phương chiều nếu có, điều kiện áp dụng, hệ đơn vị và kiểm tra tính hợp lí của kết quả. Nhận định “Độ lớn cảm ứng từ tổng hợp luôn bằng tổng độ lớn các cảm ứng từ thành phần.” sai vì trái với kiến thức nêu trên.
}
\end{bt}

% ===== Câu 223 =====
% ID: L12C3B14-61-TLN
% Mức: TH
\begin{ex}
% ID: L12C3B14-61-TLN
% Mức: TH
Trong bốn phát biểu sau về so sánh, tổng hợp và biểu diễn các từ trường đơn giản, có bao nhiêu phát biểu đúng? (Chỉ ghi một số nguyên.) (1) Vectơ cảm ứng từ tổng hợp tại một điểm bằng tổng vectơ các cảm ứng từ thành phần. (2) Độ lớn cảm ứng từ tổng hợp luôn bằng tổng độ lớn các cảm ứng từ thành phần. (3) Cần kiểm tra điều kiện áp dụng trước khi tính toán. (4) Có thể bỏ qua đơn vị của kết quả.
\shortans{2}
\loigiai{
Đối chiếu từng phát biểu với kiến thức: Vectơ cảm ứng từ tổng hợp tại một điểm bằng tổng vectơ các cảm ứng từ thành phần. Các phát biểu đúng có số lượng là 2.
}
\end{ex}

% ===== Câu 224 =====
% ID: L12C3B14-61-TN
% Mức: VD
\dangbt{Xác định chiều từ trường bằng kim nam châm và quy tắc nắm tay phải}
\begin{ex}
% ID: L12C3B14-61-TN
% Mức: VD
Cơ sở Vật lí đúng dùng để giải xác định chiều từ trường bằng kim nam châm và quy tắc nắm tay phải là gì?
\choice
{Có thể bỏ qua phương, chiều và đơn vị của các đại lượng trong mọi trường hợp.}
{Đại lượng đang xét không phụ thuộc các điều kiện vật lí nêu trong bài.}
{Quy tắc bàn tay trái dùng để xác định chiều đường sức từ quanh dây dẫn thẳng.}
{\True Đầu Bắc của kim nam châm nhỏ chỉ theo chiều của vectơ cảm ứng từ; quy tắc nắm tay phải xác định chiều đường sức quanh dòng điện.}
\loigiai{
Đầu Bắc của kim nam châm nhỏ chỉ theo chiều của vectơ cảm ứng từ; quy tắc nắm tay phải xác định chiều đường sức quanh dòng điện. Vì vậy phương án $D$ đúng; các phương án còn lại trái với định nghĩa, định luật hoặc điều kiện áp dụng.
}
\end{ex}

% ===== Câu 225 =====
% ID: L12C3B14-62-TL
% Mức: TH
\dangbt{So sánh, tổng hợp và biểu diễn các từ trường đơn giản}
\begin{bt}
% ID: L12C3B14-62-TL
% Mức: TH
Giải thích vì sao nhận định sau đúng: “Vectơ cảm ứng từ tổng hợp tại một điểm bằng tổng vectơ các cảm ứng từ thành phần.”
\loigiai{
Cần nêu được: Vectơ cảm ứng từ tổng hợp tại một điểm bằng tổng vectơ các cảm ứng từ thành phần. Khi vận dụng phải xác định đúng dữ kiện, phương chiều nếu có, điều kiện áp dụng, hệ đơn vị và kiểm tra tính hợp lí của kết quả. Nhận định “Độ lớn cảm ứng từ tổng hợp luôn bằng tổng độ lớn các cảm ứng từ thành phần.” sai vì trái với kiến thức nêu trên.
}
\end{bt}

% ===== Câu 226 =====
% ID: L12C3B14-62-TLN
% Mức: TH
\begin{ex}
% ID: L12C3B14-62-TLN
% Mức: TH
Trong bốn phát biểu sau về so sánh, tổng hợp và biểu diễn các từ trường đơn giản, có bao nhiêu phát biểu đúng? (Chỉ ghi một số nguyên.) (1) Độ lớn cảm ứng từ tổng hợp luôn bằng tổng độ lớn các cảm ứng từ thành phần. (2) Vectơ cảm ứng từ tổng hợp tại một điểm bằng tổng vectơ các cảm ứng từ thành phần. (3) Kết quả phải phù hợp ý nghĩa vật lí. (4) Mọi đại lượng đều là vô hướng.
\shortans{1}
\loigiai{
Đối chiếu từng phát biểu với kiến thức: Vectơ cảm ứng từ tổng hợp tại một điểm bằng tổng vectơ các cảm ứng từ thành phần. Các phát biểu đúng có số lượng là 1.
}
\end{ex}

% ===== Câu 227 =====
% ID: L12C3B14-62-TN
% Mức: VD
\dangbt{Xác định chiều từ trường bằng kim nam châm và quy tắc nắm tay phải}
\begin{ex}
% ID: L12C3B14-62-TN
% Mức: VD
Trong một bài thực hành về xác định chiều từ trường bằng kim nam châm và quy tắc nắm tay phải, nhận xét nào đúng?
\choice
{\True Đầu Bắc của kim nam châm nhỏ chỉ theo chiều của vectơ cảm ứng từ; quy tắc nắm tay phải xác định chiều đường sức quanh dòng điện.}
{Quy tắc bàn tay trái dùng để xác định chiều đường sức từ quanh dây dẫn thẳng.}
{Đại lượng đang xét không phụ thuộc các điều kiện vật lí nêu trong bài.}
{Có thể bỏ qua phương, chiều và đơn vị của các đại lượng trong mọi trường hợp.}
\loigiai{
Đầu Bắc của kim nam châm nhỏ chỉ theo chiều của vectơ cảm ứng từ; quy tắc nắm tay phải xác định chiều đường sức quanh dòng điện. Vì vậy phương án $A$ đúng; các phương án còn lại trái với định nghĩa, định luật hoặc điều kiện áp dụng.
}
\end{ex}

% ===== Câu 228 =====
% ID: L12C3B14-63-TL
% Mức: VD
\dangbt{So sánh, tổng hợp và biểu diễn các từ trường đơn giản}
\begin{bt}
% ID: L12C3B14-63-TL
% Mức: VD
Phân tích sai lầm trong nhận định: “Độ lớn cảm ứng từ tổng hợp luôn bằng tổng độ lớn các cảm ứng từ thành phần.”
\loigiai{
Cần nêu được: Vectơ cảm ứng từ tổng hợp tại một điểm bằng tổng vectơ các cảm ứng từ thành phần. Khi vận dụng phải xác định đúng dữ kiện, phương chiều nếu có, điều kiện áp dụng, hệ đơn vị và kiểm tra tính hợp lí của kết quả. Nhận định “Độ lớn cảm ứng từ tổng hợp luôn bằng tổng độ lớn các cảm ứng từ thành phần.” sai vì trái với kiến thức nêu trên.
}
\end{bt}

% ===== Câu 229 =====
% ID: L12C3B14-63-TLN
% Mức: TH
\begin{ex}
% ID: L12C3B14-63-TLN
% Mức: TH
Trong bốn phát biểu sau về so sánh, tổng hợp và biểu diễn các từ trường đơn giản, có bao nhiêu phát biểu đúng? (Chỉ ghi một số nguyên.) (1) Cần đổi các đại lượng về hệ đơn vị thống nhất. (2) Vectơ cảm ứng từ tổng hợp tại một điểm bằng tổng vectơ các cảm ứng từ thành phần. (3) Độ lớn cảm ứng từ tổng hợp luôn bằng tổng độ lớn các cảm ứng từ thành phần. (4) Không cần xét chiều của đại lượng vectơ.
\shortans{2}
\loigiai{
Đối chiếu từng phát biểu với kiến thức: Vectơ cảm ứng từ tổng hợp tại một điểm bằng tổng vectơ các cảm ứng từ thành phần. Các phát biểu đúng có số lượng là 2.
}
\end{ex}

% ===== Câu 230 =====
% ID: L12C3B14-63-TN
% Mức: VD
\dangbt{Xác định chiều từ trường bằng kim nam châm và quy tắc nắm tay phải}
\begin{ex}
% ID: L12C3B14-63-TN
% Mức: VD
Kết luận đúng khi vận dụng xác định chiều từ trường bằng kim nam châm và quy tắc nắm tay phải là
\choice
{Quy tắc bàn tay trái dùng để xác định chiều đường sức từ quanh dây dẫn thẳng.}
{\True Đầu Bắc của kim nam châm nhỏ chỉ theo chiều của vectơ cảm ứng từ; quy tắc nắm tay phải xác định chiều đường sức quanh dòng điện.}
{Có thể bỏ qua phương, chiều và đơn vị của các đại lượng trong mọi trường hợp.}
{Đại lượng đang xét không phụ thuộc các điều kiện vật lí nêu trong bài.}
\loigiai{
Đầu Bắc của kim nam châm nhỏ chỉ theo chiều của vectơ cảm ứng từ; quy tắc nắm tay phải xác định chiều đường sức quanh dòng điện. Vì vậy phương án $B$ đúng; các phương án còn lại trái với định nghĩa, định luật hoặc điều kiện áp dụng.
}
\end{ex}

% ===== Câu 231 =====
% ID: L12C3B14-64-TL
% Mức: VD
\dangbt{So sánh, tổng hợp và biểu diễn các từ trường đơn giản}
\begin{bt}
% ID: L12C3B14-64-TL
% Mức: VD
Đề xuất các bước giải một bài toán thuộc dạng so sánh, tổng hợp và biểu diễn các từ trường đơn giản.
\loigiai{
Cần nêu được: Vectơ cảm ứng từ tổng hợp tại một điểm bằng tổng vectơ các cảm ứng từ thành phần. Khi vận dụng phải xác định đúng dữ kiện, phương chiều nếu có, điều kiện áp dụng, hệ đơn vị và kiểm tra tính hợp lí của kết quả. Nhận định “Độ lớn cảm ứng từ tổng hợp luôn bằng tổng độ lớn các cảm ứng từ thành phần.” sai vì trái với kiến thức nêu trên.
}
\end{bt}

% ===== Câu 232 =====
% ID: L12C3B14-64-TLN
% Mức: VD
\begin{ex}
% ID: L12C3B14-64-TLN
% Mức: VD
Trong bốn phát biểu sau về so sánh, tổng hợp và biểu diễn các từ trường đơn giản, có bao nhiêu phát biểu đúng? (Chỉ ghi một số nguyên.) (1) Vectơ cảm ứng từ tổng hợp tại một điểm bằng tổng vectơ các cảm ứng từ thành phần. (2) Cần đối chiếu kết quả với điều kiện bài toán. (3) Độ lớn cảm ứng từ tổng hợp luôn bằng tổng độ lớn các cảm ứng từ thành phần. (4) Mọi trường hợp đều cho cùng một kết quả.
\shortans{2}
\loigiai{
Đối chiếu từng phát biểu với kiến thức: Vectơ cảm ứng từ tổng hợp tại một điểm bằng tổng vectơ các cảm ứng từ thành phần. Các phát biểu đúng có số lượng là 2.
}
\end{ex}

% ===== Câu 233 =====
% ID: L12C3B14-64-TN
% Mức: NB
\dangbt{Nhận biết từ trường, nguồn gây ra từ trường và tương tác từ}
\begin{ex}
% ID: L12C3B14-64-TN
% Mức: NB
Trong nhận biết từ trường, nguồn gây ra từ trường và tương tác từ, nhận định nào phù hợp với kiến thức Vật lí?
\choice
{Mọi điện tích đứng yên đều tạo ra từ trường giống dòng điện.}
{\True Từ trường tồn tại xung quanh nam châm và dòng điện; nó gây lực từ lên nam châm hoặc dòng điện khác đặt trong đó.}
{Có thể bỏ qua phương, chiều và đơn vị của các đại lượng trong mọi trường hợp.}
{Đại lượng đang xét không phụ thuộc các điều kiện vật lí nêu trong bài.}
\loigiai{
Từ trường tồn tại xung quanh nam châm và dòng điện; nó gây lực từ lên nam châm hoặc dòng điện khác đặt trong đó. Vì vậy phương án $B$ đúng; các phương án còn lại trái với định nghĩa, định luật hoặc điều kiện áp dụng.
}
\end{ex}

% ===== Câu 234 =====
% ID: L12C3B14-65-TL
% Mức: VD
\dangbt{So sánh, tổng hợp và biểu diễn các từ trường đơn giản}
\begin{bt}
% ID: L12C3B14-65-TL
% Mức: VD
Nêu một tình huống thực tiễn liên quan đến so sánh, tổng hợp và biểu diễn các từ trường đơn giản và giải thích bằng kiến thức Vật lí.
\loigiai{
Cần nêu được: Vectơ cảm ứng từ tổng hợp tại một điểm bằng tổng vectơ các cảm ứng từ thành phần. Khi vận dụng phải xác định đúng dữ kiện, phương chiều nếu có, điều kiện áp dụng, hệ đơn vị và kiểm tra tính hợp lí của kết quả. Nhận định “Độ lớn cảm ứng từ tổng hợp luôn bằng tổng độ lớn các cảm ứng từ thành phần.” sai vì trái với kiến thức nêu trên.
}
\end{bt}

% ===== Câu 235 =====
% ID: L12C3B14-65-TLN
% Mức: VD
\begin{ex}
% ID: L12C3B14-65-TLN
% Mức: VD
Trong bốn phát biểu sau về so sánh, tổng hợp và biểu diễn các từ trường đơn giản, có bao nhiêu phát biểu đúng? (Chỉ ghi một số nguyên.) (1) Độ lớn cảm ứng từ tổng hợp luôn bằng tổng độ lớn các cảm ứng từ thành phần. (2) Cần đọc đúng dữ kiện và giả thiết. (3) Vectơ cảm ứng từ tổng hợp tại một điểm bằng tổng vectơ các cảm ứng từ thành phần. (4) Có thể thay số trước khi xác định công thức.
\shortans{2}
\loigiai{
Đối chiếu từng phát biểu với kiến thức: Vectơ cảm ứng từ tổng hợp tại một điểm bằng tổng vectơ các cảm ứng từ thành phần. Các phát biểu đúng có số lượng là 2.
}
\end{ex}

% ===== Câu 236 =====
% ID: L12C3B14-65-TN
% Mức: NB
\dangbt{Nhận biết từ trường, nguồn gây ra từ trường và tương tác từ}
\begin{ex}
% ID: L12C3B14-65-TN
% Mức: NB
Tình huống 3: chọn kết luận chính xác về nhận biết từ trường, nguồn gây ra từ trường và tương tác từ.
\choice
{Đại lượng đang xét không phụ thuộc các điều kiện vật lí nêu trong bài.}
{Có thể bỏ qua phương, chiều và đơn vị của các đại lượng trong mọi trường hợp.}
{\True Từ trường tồn tại xung quanh nam châm và dòng điện; nó gây lực từ lên nam châm hoặc dòng điện khác đặt trong đó.}
{Mọi điện tích đứng yên đều tạo ra từ trường giống dòng điện.}
\loigiai{
Từ trường tồn tại xung quanh nam châm và dòng điện; nó gây lực từ lên nam châm hoặc dòng điện khác đặt trong đó. Vì vậy phương án $C$ đúng; các phương án còn lại trái với định nghĩa, định luật hoặc điều kiện áp dụng.
}
\end{ex}

% ===== Câu 237 =====
% ID: L12C3B14-66-TL
% Mức: VDC
\dangbt{So sánh, tổng hợp và biểu diễn các từ trường đơn giản}
\begin{bt}
% ID: L12C3B14-66-TL
% Mức: VDC
So sánh nhận định đúng “Vectơ cảm ứng từ tổng hợp tại một điểm bằng tổng vectơ các cảm ứng từ thành phần.” với nhận định sai “Độ lớn cảm ứng từ tổng hợp luôn bằng tổng độ lớn các cảm ứng từ thành phần.”; chỉ rõ căn cứ Vật lí.
\loigiai{
Cần nêu được: Vectơ cảm ứng từ tổng hợp tại một điểm bằng tổng vectơ các cảm ứng từ thành phần. Khi vận dụng phải xác định đúng dữ kiện, phương chiều nếu có, điều kiện áp dụng, hệ đơn vị và kiểm tra tính hợp lí của kết quả. Nhận định “Độ lớn cảm ứng từ tổng hợp luôn bằng tổng độ lớn các cảm ứng từ thành phần.” sai vì trái với kiến thức nêu trên.
}
\end{bt}

% ===== Câu 238 =====
% ID: L12C3B14-66-TLN
% Mức: VDC
\begin{ex}
% ID: L12C3B14-66-TLN
% Mức: VDC
Trong bốn phát biểu sau về so sánh, tổng hợp và biểu diễn các từ trường đơn giản, có bao nhiêu phát biểu đúng? (Chỉ ghi một số nguyên.) (1) Kết quả cần có ý nghĩa vật lí. (2) Độ lớn cảm ứng từ tổng hợp luôn bằng tổng độ lớn các cảm ứng từ thành phần. (3) Phải dùng đúng định luật cho tình huống. (4) Vectơ cảm ứng từ tổng hợp tại một điểm bằng tổng vectơ các cảm ứng từ thành phần.
\shortans{3}
\loigiai{
Đối chiếu từng phát biểu với kiến thức: Vectơ cảm ứng từ tổng hợp tại một điểm bằng tổng vectơ các cảm ứng từ thành phần. Các phát biểu đúng có số lượng là 3.
}
\end{ex}

% ===== Câu 239 =====
% ID: L12C3B14-66-TN
% Mức: TH
\dangbt{Nhận biết từ trường, nguồn gây ra từ trường và tương tác từ}
\begin{ex}
% ID: L12C3B14-66-TN
% Mức: TH
Nội dung nào mô tả đúng nhất nhận biết từ trường, nguồn gây ra từ trường và tương tác từ?
\choice
{Có thể bỏ qua phương, chiều và đơn vị của các đại lượng trong mọi trường hợp.}
{Đại lượng đang xét không phụ thuộc các điều kiện vật lí nêu trong bài.}
{Mọi điện tích đứng yên đều tạo ra từ trường giống dòng điện.}
{\True Từ trường tồn tại xung quanh nam châm và dòng điện; nó gây lực từ lên nam châm hoặc dòng điện khác đặt trong đó.}
\loigiai{
Từ trường tồn tại xung quanh nam châm và dòng điện; nó gây lực từ lên nam châm hoặc dòng điện khác đặt trong đó. Vì vậy phương án $D$ đúng; các phương án còn lại trái với định nghĩa, định luật hoặc điều kiện áp dụng.
}
\end{ex}

% ===== Câu 240 =====
% ID: L12C3B14-67-TL
% Mức: TH
\dangbt{Giải thích hiện tượng và ứng dụng của từ trường trong thực tiễn}
\begin{bt}
% ID: L12C3B14-67-TL
% Mức: TH
Trình bày kiến thức cốt lõi của giải thích hiện tượng và ứng dụng của từ trường trong thực tiễn và nêu điều kiện áp dụng.
\loigiai{
Cần nêu được: La bàn định hướng nhờ kim nam châm tương tác với từ trường Trái Đất. Khi vận dụng phải xác định đúng dữ kiện, phương chiều nếu có, điều kiện áp dụng, hệ đơn vị và kiểm tra tính hợp lí của kết quả. Nhận định “La bàn hoạt động chủ yếu nhờ điện trường của Trái Đất.” sai vì trái với kiến thức nêu trên.
}
\end{bt}

% ===== Câu 241 =====
% ID: L12C3B14-67-TLN
% Mức: TH
\begin{ex}
% ID: L12C3B14-67-TLN
% Mức: TH
Trong bốn phát biểu sau về giải thích hiện tượng và ứng dụng của từ trường trong thực tiễn, có bao nhiêu phát biểu đúng? (Chỉ ghi một số nguyên.) (1) La bàn định hướng nhờ kim nam châm tương tác với từ trường Trái Đất. (2) La bàn hoạt động chủ yếu nhờ điện trường của Trái Đất. (3) Cần kiểm tra điều kiện áp dụng trước khi tính toán. (4) Có thể bỏ qua đơn vị của kết quả.
\shortans{2}
\loigiai{
Đối chiếu từng phát biểu với kiến thức: La bàn định hướng nhờ kim nam châm tương tác với từ trường Trái Đất. Các phát biểu đúng có số lượng là 2.
}
\end{ex}

% ===== Câu 242 =====
% ID: L12C3B14-67-TN
% Mức: TH
\dangbt{Nhận biết từ trường, nguồn gây ra từ trường và tương tác từ}
\begin{ex}
% ID: L12C3B14-67-TN
% Mức: TH
Khi phân tích nhận biết từ trường, nguồn gây ra từ trường và tương tác từ, phát biểu nào đúng?
\choice
{\True Từ trường tồn tại xung quanh nam châm và dòng điện; nó gây lực từ lên nam châm hoặc dòng điện khác đặt trong đó.}
{Mọi điện tích đứng yên đều tạo ra từ trường giống dòng điện.}
{Đại lượng đang xét không phụ thuộc các điều kiện vật lí nêu trong bài.}
{Có thể bỏ qua phương, chiều và đơn vị của các đại lượng trong mọi trường hợp.}
\loigiai{
Từ trường tồn tại xung quanh nam châm và dòng điện; nó gây lực từ lên nam châm hoặc dòng điện khác đặt trong đó. Vì vậy phương án $A$ đúng; các phương án còn lại trái với định nghĩa, định luật hoặc điều kiện áp dụng.
}
\end{ex}

% ===== Câu 243 =====
% ID: L12C3B14-68-TL
% Mức: TH
\dangbt{Giải thích hiện tượng và ứng dụng của từ trường trong thực tiễn}
\begin{bt}
% ID: L12C3B14-68-TL
% Mức: TH
Giải thích vì sao nhận định sau đúng: “La bàn định hướng nhờ kim nam châm tương tác với từ trường Trái Đất.”
\loigiai{
Cần nêu được: La bàn định hướng nhờ kim nam châm tương tác với từ trường Trái Đất. Khi vận dụng phải xác định đúng dữ kiện, phương chiều nếu có, điều kiện áp dụng, hệ đơn vị và kiểm tra tính hợp lí của kết quả. Nhận định “La bàn hoạt động chủ yếu nhờ điện trường của Trái Đất.” sai vì trái với kiến thức nêu trên.
}
\end{bt}

% ===== Câu 244 =====
% ID: L12C3B14-68-TLN
% Mức: TH
\begin{ex}
% ID: L12C3B14-68-TLN
% Mức: TH
Trong bốn phát biểu sau về giải thích hiện tượng và ứng dụng của từ trường trong thực tiễn, có bao nhiêu phát biểu đúng? (Chỉ ghi một số nguyên.) (1) La bàn hoạt động chủ yếu nhờ điện trường của Trái Đất. (2) La bàn định hướng nhờ kim nam châm tương tác với từ trường Trái Đất. (3) Kết quả phải phù hợp ý nghĩa vật lí. (4) Mọi đại lượng đều là vô hướng.
\shortans{1}
\loigiai{
Đối chiếu từng phát biểu với kiến thức: La bàn định hướng nhờ kim nam châm tương tác với từ trường Trái Đất. Các phát biểu đúng có số lượng là 1.
}
\end{ex}

% ===== Câu 245 =====
% ID: L12C3B14-68-TN
% Mức: TH
\dangbt{Nhận biết từ trường, nguồn gây ra từ trường và tương tác từ}
\begin{ex}
% ID: L12C3B14-68-TN
% Mức: TH
Một học sinh ôn tập nhận biết từ trường, nguồn gây ra từ trường và tương tác từ. Kết luận nào của học sinh là đúng?
\choice
{Mọi điện tích đứng yên đều tạo ra từ trường giống dòng điện.}
{\True Từ trường tồn tại xung quanh nam châm và dòng điện; nó gây lực từ lên nam châm hoặc dòng điện khác đặt trong đó.}
{Có thể bỏ qua phương, chiều và đơn vị của các đại lượng trong mọi trường hợp.}
{Đại lượng đang xét không phụ thuộc các điều kiện vật lí nêu trong bài.}
\loigiai{
Từ trường tồn tại xung quanh nam châm và dòng điện; nó gây lực từ lên nam châm hoặc dòng điện khác đặt trong đó. Vì vậy phương án $B$ đúng; các phương án còn lại trái với định nghĩa, định luật hoặc điều kiện áp dụng.
}
\end{ex}

% ===== Câu 246 =====
% ID: L12C3B14-69-TL
% Mức: VD
\dangbt{Giải thích hiện tượng và ứng dụng của từ trường trong thực tiễn}
\begin{bt}
% ID: L12C3B14-69-TL
% Mức: VD
Phân tích sai lầm trong nhận định: “La bàn hoạt động chủ yếu nhờ điện trường của Trái Đất.”
\loigiai{
Cần nêu được: La bàn định hướng nhờ kim nam châm tương tác với từ trường Trái Đất. Khi vận dụng phải xác định đúng dữ kiện, phương chiều nếu có, điều kiện áp dụng, hệ đơn vị và kiểm tra tính hợp lí của kết quả. Nhận định “La bàn hoạt động chủ yếu nhờ điện trường của Trái Đất.” sai vì trái với kiến thức nêu trên.
}
\end{bt}

% ===== Câu 247 =====
% ID: L12C3B14-69-TLN
% Mức: TH
\begin{ex}
% ID: L12C3B14-69-TLN
% Mức: TH
Trong bốn phát biểu sau về giải thích hiện tượng và ứng dụng của từ trường trong thực tiễn, có bao nhiêu phát biểu đúng? (Chỉ ghi một số nguyên.) (1) Cần đổi các đại lượng về hệ đơn vị thống nhất. (2) La bàn định hướng nhờ kim nam châm tương tác với từ trường Trái Đất. (3) La bàn hoạt động chủ yếu nhờ điện trường của Trái Đất. (4) Không cần xét chiều của đại lượng vectơ.
\shortans{2}
\loigiai{
Đối chiếu từng phát biểu với kiến thức: La bàn định hướng nhờ kim nam châm tương tác với từ trường Trái Đất. Các phát biểu đúng có số lượng là 2.
}
\end{ex}

% ===== Câu 248 =====
% ID: L12C3B14-69-TN
% Mức: TH
\dangbt{Nhận biết từ trường, nguồn gây ra từ trường và tương tác từ}
\begin{ex}
% ID: L12C3B14-69-TN
% Mức: TH
Chọn phát biểu không trái với kiến thức về nhận biết từ trường, nguồn gây ra từ trường và tương tác từ.
\choice
{Đại lượng đang xét không phụ thuộc các điều kiện vật lí nêu trong bài.}
{Có thể bỏ qua phương, chiều và đơn vị của các đại lượng trong mọi trường hợp.}
{\True Từ trường tồn tại xung quanh nam châm và dòng điện; nó gây lực từ lên nam châm hoặc dòng điện khác đặt trong đó.}
{Mọi điện tích đứng yên đều tạo ra từ trường giống dòng điện.}
\loigiai{
Từ trường tồn tại xung quanh nam châm và dòng điện; nó gây lực từ lên nam châm hoặc dòng điện khác đặt trong đó. Vì vậy phương án $C$ đúng; các phương án còn lại trái với định nghĩa, định luật hoặc điều kiện áp dụng.
}
\end{ex}

% ===== Câu 249 =====
% ID: L12C3B14-70-TL
% Mức: VD
\dangbt{Giải thích hiện tượng và ứng dụng của từ trường trong thực tiễn}
\begin{bt}
% ID: L12C3B14-70-TL
% Mức: VD
Đề xuất các bước giải một bài toán thuộc dạng giải thích hiện tượng và ứng dụng của từ trường trong thực tiễn.
\loigiai{
Cần nêu được: La bàn định hướng nhờ kim nam châm tương tác với từ trường Trái Đất. Khi vận dụng phải xác định đúng dữ kiện, phương chiều nếu có, điều kiện áp dụng, hệ đơn vị và kiểm tra tính hợp lí của kết quả. Nhận định “La bàn hoạt động chủ yếu nhờ điện trường của Trái Đất.” sai vì trái với kiến thức nêu trên.
}
\end{bt}

% ===== Câu 250 =====
% ID: L12C3B14-70-TLN
% Mức: VD
\begin{ex}
% ID: L12C3B14-70-TLN
% Mức: VD
Trong bốn phát biểu sau về giải thích hiện tượng và ứng dụng của từ trường trong thực tiễn, có bao nhiêu phát biểu đúng? (Chỉ ghi một số nguyên.) (1) La bàn định hướng nhờ kim nam châm tương tác với từ trường Trái Đất. (2) Cần đối chiếu kết quả với điều kiện bài toán. (3) La bàn hoạt động chủ yếu nhờ điện trường của Trái Đất. (4) Mọi trường hợp đều cho cùng một kết quả.
\shortans{2}
\loigiai{
Đối chiếu từng phát biểu với kiến thức: La bàn định hướng nhờ kim nam châm tương tác với từ trường Trái Đất. Các phát biểu đúng có số lượng là 2.
}
\end{ex}

% ===== Câu 251 =====
% ID: L12C3B14-70-TN
% Mức: VD
\dangbt{Nhận biết từ trường, nguồn gây ra từ trường và tương tác từ}
\begin{ex}
% ID: L12C3B14-70-TN
% Mức: VD
Cơ sở Vật lí đúng dùng để giải nhận biết từ trường, nguồn gây ra từ trường và tương tác từ là gì?
\choice
{Có thể bỏ qua phương, chiều và đơn vị của các đại lượng trong mọi trường hợp.}
{Đại lượng đang xét không phụ thuộc các điều kiện vật lí nêu trong bài.}
{Mọi điện tích đứng yên đều tạo ra từ trường giống dòng điện.}
{\True Từ trường tồn tại xung quanh nam châm và dòng điện; nó gây lực từ lên nam châm hoặc dòng điện khác đặt trong đó.}
\loigiai{
Từ trường tồn tại xung quanh nam châm và dòng điện; nó gây lực từ lên nam châm hoặc dòng điện khác đặt trong đó. Vì vậy phương án $D$ đúng; các phương án còn lại trái với định nghĩa, định luật hoặc điều kiện áp dụng.
}
\end{ex}

% ===== Câu 252 =====
% ID: L12C3B14-71-TL
% Mức: VD
\dangbt{Giải thích hiện tượng và ứng dụng của từ trường trong thực tiễn}
\begin{bt}
% ID: L12C3B14-71-TL
% Mức: VD
Nêu một tình huống thực tiễn liên quan đến giải thích hiện tượng và ứng dụng của từ trường trong thực tiễn và giải thích bằng kiến thức Vật lí.
\loigiai{
Cần nêu được: La bàn định hướng nhờ kim nam châm tương tác với từ trường Trái Đất. Khi vận dụng phải xác định đúng dữ kiện, phương chiều nếu có, điều kiện áp dụng, hệ đơn vị và kiểm tra tính hợp lí của kết quả. Nhận định “La bàn hoạt động chủ yếu nhờ điện trường của Trái Đất.” sai vì trái với kiến thức nêu trên.
}
\end{bt}

% ===== Câu 253 =====
% ID: L12C3B14-71-TLN
% Mức: VD
\begin{ex}
% ID: L12C3B14-71-TLN
% Mức: VD
Trong bốn phát biểu sau về giải thích hiện tượng và ứng dụng của từ trường trong thực tiễn, có bao nhiêu phát biểu đúng? (Chỉ ghi một số nguyên.) (1) La bàn hoạt động chủ yếu nhờ điện trường của Trái Đất. (2) Cần đọc đúng dữ kiện và giả thiết. (3) La bàn định hướng nhờ kim nam châm tương tác với từ trường Trái Đất. (4) Có thể thay số trước khi xác định công thức.
\shortans{2}
\loigiai{
Đối chiếu từng phát biểu với kiến thức: La bàn định hướng nhờ kim nam châm tương tác với từ trường Trái Đất. Các phát biểu đúng có số lượng là 2.
}
\end{ex}

% ===== Câu 254 =====
% ID: L12C3B14-71-TN
% Mức: VD
\dangbt{Nhận biết từ trường, nguồn gây ra từ trường và tương tác từ}
\begin{ex}
% ID: L12C3B14-71-TN
% Mức: VD
Trong một bài thực hành về nhận biết từ trường, nguồn gây ra từ trường và tương tác từ, nhận xét nào đúng?
\choice
{\True Từ trường tồn tại xung quanh nam châm và dòng điện; nó gây lực từ lên nam châm hoặc dòng điện khác đặt trong đó.}
{Mọi điện tích đứng yên đều tạo ra từ trường giống dòng điện.}
{Đại lượng đang xét không phụ thuộc các điều kiện vật lí nêu trong bài.}
{Có thể bỏ qua phương, chiều và đơn vị của các đại lượng trong mọi trường hợp.}
\loigiai{
Từ trường tồn tại xung quanh nam châm và dòng điện; nó gây lực từ lên nam châm hoặc dòng điện khác đặt trong đó. Vì vậy phương án $A$ đúng; các phương án còn lại trái với định nghĩa, định luật hoặc điều kiện áp dụng.
}
\end{ex}

% ===== Câu 255 =====
% ID: L12C3B14-72-TL
% Mức: VDC
\dangbt{Giải thích hiện tượng và ứng dụng của từ trường trong thực tiễn}
\begin{bt}
% ID: L12C3B14-72-TL
% Mức: VDC
So sánh nhận định đúng “La bàn định hướng nhờ kim nam châm tương tác với từ trường Trái Đất.” với nhận định sai “La bàn hoạt động chủ yếu nhờ điện trường của Trái Đất.”; chỉ rõ căn cứ Vật lí.
\loigiai{
Cần nêu được: La bàn định hướng nhờ kim nam châm tương tác với từ trường Trái Đất. Khi vận dụng phải xác định đúng dữ kiện, phương chiều nếu có, điều kiện áp dụng, hệ đơn vị và kiểm tra tính hợp lí của kết quả. Nhận định “La bàn hoạt động chủ yếu nhờ điện trường của Trái Đất.” sai vì trái với kiến thức nêu trên.
}
\end{bt}

% ===== Câu 256 =====
% ID: L12C3B14-72-TLN
% Mức: VDC
\begin{ex}
% ID: L12C3B14-72-TLN
% Mức: VDC
Trong bốn phát biểu sau về giải thích hiện tượng và ứng dụng của từ trường trong thực tiễn, có bao nhiêu phát biểu đúng? (Chỉ ghi một số nguyên.) (1) Kết quả cần có ý nghĩa vật lí. (2) La bàn hoạt động chủ yếu nhờ điện trường của Trái Đất. (3) Phải dùng đúng định luật cho tình huống. (4) La bàn định hướng nhờ kim nam châm tương tác với từ trường Trái Đất.
\shortans{3}
\loigiai{
Đối chiếu từng phát biểu với kiến thức: La bàn định hướng nhờ kim nam châm tương tác với từ trường Trái Đất. Các phát biểu đúng có số lượng là 3.
}
\end{ex}

% ===== Câu 257 =====
% ID: L12C3B14-72-TN
% Mức: VD
\dangbt{Nhận biết từ trường, nguồn gây ra từ trường và tương tác từ}
\begin{ex}
% ID: L12C3B14-72-TN
% Mức: VD
Kết luận đúng khi vận dụng nhận biết từ trường, nguồn gây ra từ trường và tương tác từ là
\choice
{Mọi điện tích đứng yên đều tạo ra từ trường giống dòng điện.}
{\True Từ trường tồn tại xung quanh nam châm và dòng điện; nó gây lực từ lên nam châm hoặc dòng điện khác đặt trong đó.}
{Có thể bỏ qua phương, chiều và đơn vị của các đại lượng trong mọi trường hợp.}
{Đại lượng đang xét không phụ thuộc các điều kiện vật lí nêu trong bài.}
\loigiai{
Từ trường tồn tại xung quanh nam châm và dòng điện; nó gây lực từ lên nam châm hoặc dòng điện khác đặt trong đó. Vì vậy phương án $B$ đúng; các phương án còn lại trái với định nghĩa, định luật hoặc điều kiện áp dụng.
}
\end{ex}

% ===== Câu 258 =====
% ID: L12C3B14-73-TN
% Mức: NB
\dangbt{Mô tả đường sức từ và xác định chiều đường sức từ}
\begin{ex}
% ID: L12C3B14-73-TN
% Mức: NB
Phát biểu nào sau đây đúng về mô tả đường sức từ và xác định chiều đường sức từ?
\choice
{\True Tiếp tuyến với đường sức từ tại một điểm cho phương của vectơ cảm ứng từ; bên ngoài nam châm, đường sức đi từ cực Bắc đến cực Nam.}
{Các đường sức từ có thể cắt nhau tại một điểm.}
{Đại lượng đang xét không phụ thuộc các điều kiện vật lí nêu trong bài.}
{Có thể bỏ qua phương, chiều và đơn vị của các đại lượng trong mọi trường hợp.}
\loigiai{
Tiếp tuyến với đường sức từ tại một điểm cho phương của vectơ cảm ứng từ; bên ngoài nam châm, đường sức đi từ cực Bắc đến cực Nam. Vì vậy phương án $A$ đúng; các phương án còn lại trái với định nghĩa, định luật hoặc điều kiện áp dụng.
}
\end{ex}

% ===== Câu 259 =====
% ID: L12C3B14-74-TN
% Mức: NB
\begin{ex}
% ID: L12C3B14-74-TN
% Mức: NB
Trong mô tả đường sức từ và xác định chiều đường sức từ, nhận định nào phù hợp với kiến thức Vật lí?
\choice
{Các đường sức từ có thể cắt nhau tại một điểm.}
{\True Tiếp tuyến với đường sức từ tại một điểm cho phương của vectơ cảm ứng từ; bên ngoài nam châm, đường sức đi từ cực Bắc đến cực Nam.}
{Có thể bỏ qua phương, chiều và đơn vị của các đại lượng trong mọi trường hợp.}
{Đại lượng đang xét không phụ thuộc các điều kiện vật lí nêu trong bài.}
\loigiai{
Tiếp tuyến với đường sức từ tại một điểm cho phương của vectơ cảm ứng từ; bên ngoài nam châm, đường sức đi từ cực Bắc đến cực Nam. Vì vậy phương án $B$ đúng; các phương án còn lại trái với định nghĩa, định luật hoặc điều kiện áp dụng.
}
\end{ex}

% ===== Câu 260 =====
% ID: L12C3B14-75-TN
% Mức: NB
\begin{ex}
% ID: L12C3B14-75-TN
% Mức: NB
Tình huống 3: chọn kết luận chính xác về mô tả đường sức từ và xác định chiều đường sức từ.
\choice
{Đại lượng đang xét không phụ thuộc các điều kiện vật lí nêu trong bài.}
{Có thể bỏ qua phương, chiều và đơn vị của các đại lượng trong mọi trường hợp.}
{\True Tiếp tuyến với đường sức từ tại một điểm cho phương của vectơ cảm ứng từ; bên ngoài nam châm, đường sức đi từ cực Bắc đến cực Nam.}
{Các đường sức từ có thể cắt nhau tại một điểm.}
\loigiai{
Tiếp tuyến với đường sức từ tại một điểm cho phương của vectơ cảm ứng từ; bên ngoài nam châm, đường sức đi từ cực Bắc đến cực Nam. Vì vậy phương án $C$ đúng; các phương án còn lại trái với định nghĩa, định luật hoặc điều kiện áp dụng.
}
\end{ex}

% ===== Câu 261 =====
% ID: L12C3B14-76-TN
% Mức: TH
\begin{ex}
% ID: L12C3B14-76-TN
% Mức: TH
Nội dung nào mô tả đúng nhất mô tả đường sức từ và xác định chiều đường sức từ?
\choice
{Có thể bỏ qua phương, chiều và đơn vị của các đại lượng trong mọi trường hợp.}
{Đại lượng đang xét không phụ thuộc các điều kiện vật lí nêu trong bài.}
{Các đường sức từ có thể cắt nhau tại một điểm.}
{\True Tiếp tuyến với đường sức từ tại một điểm cho phương của vectơ cảm ứng từ; bên ngoài nam châm, đường sức đi từ cực Bắc đến cực Nam.}
\loigiai{
Tiếp tuyến với đường sức từ tại một điểm cho phương của vectơ cảm ứng từ; bên ngoài nam châm, đường sức đi từ cực Bắc đến cực Nam. Vì vậy phương án $D$ đúng; các phương án còn lại trái với định nghĩa, định luật hoặc điều kiện áp dụng.
}
\end{ex}

% ===== Câu 262 =====
% ID: L12C3B14-77-TN
% Mức: TH
\begin{ex}
% ID: L12C3B14-77-TN
% Mức: TH
Khi phân tích mô tả đường sức từ và xác định chiều đường sức từ, phát biểu nào đúng?
\choice
{\True Tiếp tuyến với đường sức từ tại một điểm cho phương của vectơ cảm ứng từ; bên ngoài nam châm, đường sức đi từ cực Bắc đến cực Nam.}
{Các đường sức từ có thể cắt nhau tại một điểm.}
{Đại lượng đang xét không phụ thuộc các điều kiện vật lí nêu trong bài.}
{Có thể bỏ qua phương, chiều và đơn vị của các đại lượng trong mọi trường hợp.}
\loigiai{
Tiếp tuyến với đường sức từ tại một điểm cho phương của vectơ cảm ứng từ; bên ngoài nam châm, đường sức đi từ cực Bắc đến cực Nam. Vì vậy phương án $A$ đúng; các phương án còn lại trái với định nghĩa, định luật hoặc điều kiện áp dụng.
}
\end{ex}

% ===== Câu 263 =====
% ID: L12C3B14-78-TN
% Mức: TH
\begin{ex}
% ID: L12C3B14-78-TN
% Mức: TH
Một học sinh ôn tập mô tả đường sức từ và xác định chiều đường sức từ. Kết luận nào của học sinh là đúng?
\choice
{Các đường sức từ có thể cắt nhau tại một điểm.}
{\True Tiếp tuyến với đường sức từ tại một điểm cho phương của vectơ cảm ứng từ; bên ngoài nam châm, đường sức đi từ cực Bắc đến cực Nam.}
{Có thể bỏ qua phương, chiều và đơn vị của các đại lượng trong mọi trường hợp.}
{Đại lượng đang xét không phụ thuộc các điều kiện vật lí nêu trong bài.}
\loigiai{
Tiếp tuyến với đường sức từ tại một điểm cho phương của vectơ cảm ứng từ; bên ngoài nam châm, đường sức đi từ cực Bắc đến cực Nam. Vì vậy phương án $B$ đúng; các phương án còn lại trái với định nghĩa, định luật hoặc điều kiện áp dụng.
}
\end{ex}

% ===== Câu 264 =====
% ID: L12C3B14-79-TN
% Mức: TH
\begin{ex}
% ID: L12C3B14-79-TN
% Mức: TH
Chọn phát biểu không trái với kiến thức về mô tả đường sức từ và xác định chiều đường sức từ.
\choice
{Đại lượng đang xét không phụ thuộc các điều kiện vật lí nêu trong bài.}
{Có thể bỏ qua phương, chiều và đơn vị của các đại lượng trong mọi trường hợp.}
{\True Tiếp tuyến với đường sức từ tại một điểm cho phương của vectơ cảm ứng từ; bên ngoài nam châm, đường sức đi từ cực Bắc đến cực Nam.}
{Các đường sức từ có thể cắt nhau tại một điểm.}
\loigiai{
Tiếp tuyến với đường sức từ tại một điểm cho phương của vectơ cảm ứng từ; bên ngoài nam châm, đường sức đi từ cực Bắc đến cực Nam. Vì vậy phương án $C$ đúng; các phương án còn lại trái với định nghĩa, định luật hoặc điều kiện áp dụng.
}
\end{ex}

% ===== Câu 265 =====
% ID: L12C3B14-80-TN
% Mức: VD
\begin{ex}
% ID: L12C3B14-80-TN
% Mức: VD
Cơ sở Vật lí đúng dùng để giải mô tả đường sức từ và xác định chiều đường sức từ là gì?
\choice
{Có thể bỏ qua phương, chiều và đơn vị của các đại lượng trong mọi trường hợp.}
{Đại lượng đang xét không phụ thuộc các điều kiện vật lí nêu trong bài.}
{Các đường sức từ có thể cắt nhau tại một điểm.}
{\True Tiếp tuyến với đường sức từ tại một điểm cho phương của vectơ cảm ứng từ; bên ngoài nam châm, đường sức đi từ cực Bắc đến cực Nam.}
\loigiai{
Tiếp tuyến với đường sức từ tại một điểm cho phương của vectơ cảm ứng từ; bên ngoài nam châm, đường sức đi từ cực Bắc đến cực Nam. Vì vậy phương án $D$ đúng; các phương án còn lại trái với định nghĩa, định luật hoặc điều kiện áp dụng.
}
\end{ex}

% ===== Câu 266 =====
% ID: L12C3B14-81-TN
% Mức: VD
\begin{ex}
% ID: L12C3B14-81-TN
% Mức: VD
Trong một bài thực hành về mô tả đường sức từ và xác định chiều đường sức từ, nhận xét nào đúng?
\choice
{\True Tiếp tuyến với đường sức từ tại một điểm cho phương của vectơ cảm ứng từ; bên ngoài nam châm, đường sức đi từ cực Bắc đến cực Nam.}
{Các đường sức từ có thể cắt nhau tại một điểm.}
{Đại lượng đang xét không phụ thuộc các điều kiện vật lí nêu trong bài.}
{Có thể bỏ qua phương, chiều và đơn vị của các đại lượng trong mọi trường hợp.}
\loigiai{
Tiếp tuyến với đường sức từ tại một điểm cho phương của vectơ cảm ứng từ; bên ngoài nam châm, đường sức đi từ cực Bắc đến cực Nam. Vì vậy phương án $A$ đúng; các phương án còn lại trái với định nghĩa, định luật hoặc điều kiện áp dụng.
}
\end{ex}

% ===== Câu 267 =====
% ID: L12C3B14-82-TN
% Mức: VD
\begin{ex}
% ID: L12C3B14-82-TN
% Mức: VD
Kết luận đúng khi vận dụng mô tả đường sức từ và xác định chiều đường sức từ là
\choice
{Các đường sức từ có thể cắt nhau tại một điểm.}
{\True Tiếp tuyến với đường sức từ tại một điểm cho phương của vectơ cảm ứng từ; bên ngoài nam châm, đường sức đi từ cực Bắc đến cực Nam.}
{Có thể bỏ qua phương, chiều và đơn vị của các đại lượng trong mọi trường hợp.}
{Đại lượng đang xét không phụ thuộc các điều kiện vật lí nêu trong bài.}
\loigiai{
Tiếp tuyến với đường sức từ tại một điểm cho phương của vectơ cảm ứng từ; bên ngoài nam châm, đường sức đi từ cực Bắc đến cực Nam. Vì vậy phương án $B$ đúng; các phương án còn lại trái với định nghĩa, định luật hoặc điều kiện áp dụng.
}
\end{ex}

% ===== Câu 268 =====
% ID: L12C3B14-83-TN
% Mức: NB
\dangbt{Xác định chiều từ trường bằng kim nam châm và quy tắc nắm tay phải}
\begin{ex}
% ID: L12C3B14-83-TN
% Mức: NB
Phát biểu nào sau đây đúng về xác định chiều từ trường bằng kim nam châm và quy tắc nắm tay phải?
\choice
{\True Đầu Bắc của kim nam châm nhỏ chỉ theo chiều của vectơ cảm ứng từ; quy tắc nắm tay phải xác định chiều đường sức quanh dòng điện.}
{Quy tắc bàn tay trái dùng để xác định chiều đường sức từ quanh dây dẫn thẳng.}
{Đại lượng đang xét không phụ thuộc các điều kiện vật lí nêu trong bài.}
{Có thể bỏ qua phương, chiều và đơn vị của các đại lượng trong mọi trường hợp.}
\loigiai{
Đầu Bắc của kim nam châm nhỏ chỉ theo chiều của vectơ cảm ứng từ; quy tắc nắm tay phải xác định chiều đường sức quanh dòng điện. Vì vậy phương án $A$ đúng; các phương án còn lại trái với định nghĩa, định luật hoặc điều kiện áp dụng.
}
\end{ex}

% ===== Câu 269 =====
% ID: L12C3B14-84-TN
% Mức: NB
\begin{ex}
% ID: L12C3B14-84-TN
% Mức: NB
Trong xác định chiều từ trường bằng kim nam châm và quy tắc nắm tay phải, nhận định nào phù hợp với kiến thức Vật lí?
\choice
{Quy tắc bàn tay trái dùng để xác định chiều đường sức từ quanh dây dẫn thẳng.}
{\True Đầu Bắc của kim nam châm nhỏ chỉ theo chiều của vectơ cảm ứng từ; quy tắc nắm tay phải xác định chiều đường sức quanh dòng điện.}
{Có thể bỏ qua phương, chiều và đơn vị của các đại lượng trong mọi trường hợp.}
{Đại lượng đang xét không phụ thuộc các điều kiện vật lí nêu trong bài.}
\loigiai{
Đầu Bắc của kim nam châm nhỏ chỉ theo chiều của vectơ cảm ứng từ; quy tắc nắm tay phải xác định chiều đường sức quanh dòng điện. Vì vậy phương án $B$ đúng; các phương án còn lại trái với định nghĩa, định luật hoặc điều kiện áp dụng.
}
\end{ex}

% ===== Câu 270 =====
% ID: L12C3B14-85-TN
% Mức: NB
\begin{ex}
% ID: L12C3B14-85-TN
% Mức: NB
Tình huống 3: chọn kết luận chính xác về xác định chiều từ trường bằng kim nam châm và quy tắc nắm tay phải.
\choice
{Đại lượng đang xét không phụ thuộc các điều kiện vật lí nêu trong bài.}
{Có thể bỏ qua phương, chiều và đơn vị của các đại lượng trong mọi trường hợp.}
{\True Đầu Bắc của kim nam châm nhỏ chỉ theo chiều của vectơ cảm ứng từ; quy tắc nắm tay phải xác định chiều đường sức quanh dòng điện.}
{Quy tắc bàn tay trái dùng để xác định chiều đường sức từ quanh dây dẫn thẳng.}
\loigiai{
Đầu Bắc của kim nam châm nhỏ chỉ theo chiều của vectơ cảm ứng từ; quy tắc nắm tay phải xác định chiều đường sức quanh dòng điện. Vì vậy phương án $C$ đúng; các phương án còn lại trái với định nghĩa, định luật hoặc điều kiện áp dụng.
}
\end{ex}

% ===== Câu 271 =====
% ID: L12C3B14-86-TN
% Mức: TH
\begin{ex}
% ID: L12C3B14-86-TN
% Mức: TH
Nội dung nào mô tả đúng nhất xác định chiều từ trường bằng kim nam châm và quy tắc nắm tay phải?
\choice
{Có thể bỏ qua phương, chiều và đơn vị của các đại lượng trong mọi trường hợp.}
{Đại lượng đang xét không phụ thuộc các điều kiện vật lí nêu trong bài.}
{Quy tắc bàn tay trái dùng để xác định chiều đường sức từ quanh dây dẫn thẳng.}
{\True Đầu Bắc của kim nam châm nhỏ chỉ theo chiều của vectơ cảm ứng từ; quy tắc nắm tay phải xác định chiều đường sức quanh dòng điện.}
\loigiai{
Đầu Bắc của kim nam châm nhỏ chỉ theo chiều của vectơ cảm ứng từ; quy tắc nắm tay phải xác định chiều đường sức quanh dòng điện. Vì vậy phương án $D$ đúng; các phương án còn lại trái với định nghĩa, định luật hoặc điều kiện áp dụng.
}
\end{ex}

% ===== Câu 272 =====
% ID: L12C3B14-87-TN
% Mức: TH
\begin{ex}
% ID: L12C3B14-87-TN
% Mức: TH
Khi phân tích xác định chiều từ trường bằng kim nam châm và quy tắc nắm tay phải, phát biểu nào đúng?
\choice
{\True Đầu Bắc của kim nam châm nhỏ chỉ theo chiều của vectơ cảm ứng từ; quy tắc nắm tay phải xác định chiều đường sức quanh dòng điện.}
{Quy tắc bàn tay trái dùng để xác định chiều đường sức từ quanh dây dẫn thẳng.}
{Đại lượng đang xét không phụ thuộc các điều kiện vật lí nêu trong bài.}
{Có thể bỏ qua phương, chiều và đơn vị của các đại lượng trong mọi trường hợp.}
\loigiai{
Đầu Bắc của kim nam châm nhỏ chỉ theo chiều của vectơ cảm ứng từ; quy tắc nắm tay phải xác định chiều đường sức quanh dòng điện. Vì vậy phương án $A$ đúng; các phương án còn lại trái với định nghĩa, định luật hoặc điều kiện áp dụng.
}
\end{ex}

% ===== Câu 273 =====
% ID: L12C3B14-88-TN
% Mức: TH
\begin{ex}
% ID: L12C3B14-88-TN
% Mức: TH
Một học sinh ôn tập xác định chiều từ trường bằng kim nam châm và quy tắc nắm tay phải. Kết luận nào của học sinh là đúng?
\choice
{Quy tắc bàn tay trái dùng để xác định chiều đường sức từ quanh dây dẫn thẳng.}
{\True Đầu Bắc của kim nam châm nhỏ chỉ theo chiều của vectơ cảm ứng từ; quy tắc nắm tay phải xác định chiều đường sức quanh dòng điện.}
{Có thể bỏ qua phương, chiều và đơn vị của các đại lượng trong mọi trường hợp.}
{Đại lượng đang xét không phụ thuộc các điều kiện vật lí nêu trong bài.}
\loigiai{
Đầu Bắc của kim nam châm nhỏ chỉ theo chiều của vectơ cảm ứng từ; quy tắc nắm tay phải xác định chiều đường sức quanh dòng điện. Vì vậy phương án $B$ đúng; các phương án còn lại trái với định nghĩa, định luật hoặc điều kiện áp dụng.
}
\end{ex}

% ===== Câu 274 =====
% ID: L12C3B14-89-TN
% Mức: TH
\begin{ex}
% ID: L12C3B14-89-TN
% Mức: TH
Chọn phát biểu không trái với kiến thức về xác định chiều từ trường bằng kim nam châm và quy tắc nắm tay phải.
\choice
{Đại lượng đang xét không phụ thuộc các điều kiện vật lí nêu trong bài.}
{Có thể bỏ qua phương, chiều và đơn vị của các đại lượng trong mọi trường hợp.}
{\True Đầu Bắc của kim nam châm nhỏ chỉ theo chiều của vectơ cảm ứng từ; quy tắc nắm tay phải xác định chiều đường sức quanh dòng điện.}
{Quy tắc bàn tay trái dùng để xác định chiều đường sức từ quanh dây dẫn thẳng.}
\loigiai{
Đầu Bắc của kim nam châm nhỏ chỉ theo chiều của vectơ cảm ứng từ; quy tắc nắm tay phải xác định chiều đường sức quanh dòng điện. Vì vậy phương án $C$ đúng; các phương án còn lại trái với định nghĩa, định luật hoặc điều kiện áp dụng.
}
\end{ex}

% ===== Câu 275 =====
% ID: L12C3B14-90-TN
% Mức: VD
\begin{ex}
% ID: L12C3B14-90-TN
% Mức: VD
Cơ sở Vật lí đúng dùng để giải xác định chiều từ trường bằng kim nam châm và quy tắc nắm tay phải là gì?
\choice
{Có thể bỏ qua phương, chiều và đơn vị của các đại lượng trong mọi trường hợp.}
{Đại lượng đang xét không phụ thuộc các điều kiện vật lí nêu trong bài.}
{Quy tắc bàn tay trái dùng để xác định chiều đường sức từ quanh dây dẫn thẳng.}
{\True Đầu Bắc của kim nam châm nhỏ chỉ theo chiều của vectơ cảm ứng từ; quy tắc nắm tay phải xác định chiều đường sức quanh dòng điện.}
\loigiai{
Đầu Bắc của kim nam châm nhỏ chỉ theo chiều của vectơ cảm ứng từ; quy tắc nắm tay phải xác định chiều đường sức quanh dòng điện. Vì vậy phương án $D$ đúng; các phương án còn lại trái với định nghĩa, định luật hoặc điều kiện áp dụng.
}
\end{ex}

% ===== Câu 276 =====
% ID: L12C3B14-91-TN
% Mức: VD
\begin{ex}
% ID: L12C3B14-91-TN
% Mức: VD
Trong một bài thực hành về xác định chiều từ trường bằng kim nam châm và quy tắc nắm tay phải, nhận xét nào đúng?
\choice
{\True Đầu Bắc của kim nam châm nhỏ chỉ theo chiều của vectơ cảm ứng từ; quy tắc nắm tay phải xác định chiều đường sức quanh dòng điện.}
{Quy tắc bàn tay trái dùng để xác định chiều đường sức từ quanh dây dẫn thẳng.}
{Đại lượng đang xét không phụ thuộc các điều kiện vật lí nêu trong bài.}
{Có thể bỏ qua phương, chiều và đơn vị của các đại lượng trong mọi trường hợp.}
\loigiai{
Đầu Bắc của kim nam châm nhỏ chỉ theo chiều của vectơ cảm ứng từ; quy tắc nắm tay phải xác định chiều đường sức quanh dòng điện. Vì vậy phương án $A$ đúng; các phương án còn lại trái với định nghĩa, định luật hoặc điều kiện áp dụng.
}
\end{ex}

% ===== Câu 277 =====
% ID: L12C3B14-92-TN
% Mức: VD
\begin{ex}
% ID: L12C3B14-92-TN
% Mức: VD
Kết luận đúng khi vận dụng xác định chiều từ trường bằng kim nam châm và quy tắc nắm tay phải là
\choice
{Quy tắc bàn tay trái dùng để xác định chiều đường sức từ quanh dây dẫn thẳng.}
{\True Đầu Bắc của kim nam châm nhỏ chỉ theo chiều của vectơ cảm ứng từ; quy tắc nắm tay phải xác định chiều đường sức quanh dòng điện.}
{Có thể bỏ qua phương, chiều và đơn vị của các đại lượng trong mọi trường hợp.}
{Đại lượng đang xét không phụ thuộc các điều kiện vật lí nêu trong bài.}
\loigiai{
Đầu Bắc của kim nam châm nhỏ chỉ theo chiều của vectơ cảm ứng từ; quy tắc nắm tay phải xác định chiều đường sức quanh dòng điện. Vì vậy phương án $B$ đúng; các phương án còn lại trái với định nghĩa, định luật hoặc điều kiện áp dụng.
}
\end{ex}

% ===== Câu 278 =====
% ID: L12C3B14-93-TN
% Mức: NB
\dangbt{Từ trường của dòng điện thẳng, dòng điện tròn và ống dây}
\begin{ex}
% ID: L12C3B14-93-TN
% Mức: NB
Phát biểu nào sau đây đúng về từ trường của dòng điện thẳng, dòng điện tròn và ống dây?
\choice
{\True Đường sức quanh dây dẫn thẳng là các đường tròn đồng tâm; từ trường trong lòng ống dây dài gần đều.}
{Từ trường trong lòng ống dây dài luôn bằng không.}
{Đại lượng đang xét không phụ thuộc các điều kiện vật lí nêu trong bài.}
{Có thể bỏ qua phương, chiều và đơn vị của các đại lượng trong mọi trường hợp.}
\loigiai{
Đường sức quanh dây dẫn thẳng là các đường tròn đồng tâm; từ trường trong lòng ống dây dài gần đều. Vì vậy phương án $A$ đúng; các phương án còn lại trái với định nghĩa, định luật hoặc điều kiện áp dụng.
}
\end{ex}

% ===== Câu 279 =====
% ID: L12C3B14-94-TN
% Mức: NB
\begin{ex}
% ID: L12C3B14-94-TN
% Mức: NB
Trong từ trường của dòng điện thẳng, dòng điện tròn và ống dây, nhận định nào phù hợp với kiến thức Vật lí?
\choice
{Từ trường trong lòng ống dây dài luôn bằng không.}
{\True Đường sức quanh dây dẫn thẳng là các đường tròn đồng tâm; từ trường trong lòng ống dây dài gần đều.}
{Có thể bỏ qua phương, chiều và đơn vị của các đại lượng trong mọi trường hợp.}
{Đại lượng đang xét không phụ thuộc các điều kiện vật lí nêu trong bài.}
\loigiai{
Đường sức quanh dây dẫn thẳng là các đường tròn đồng tâm; từ trường trong lòng ống dây dài gần đều. Vì vậy phương án $B$ đúng; các phương án còn lại trái với định nghĩa, định luật hoặc điều kiện áp dụng.
}
\end{ex}

% ===== Câu 280 =====
% ID: L12C3B14-95-TN
% Mức: NB
\begin{ex}
% ID: L12C3B14-95-TN
% Mức: NB
Tình huống 3: chọn kết luận chính xác về từ trường của dòng điện thẳng, dòng điện tròn và ống dây.
\choice
{Đại lượng đang xét không phụ thuộc các điều kiện vật lí nêu trong bài.}
{Có thể bỏ qua phương, chiều và đơn vị của các đại lượng trong mọi trường hợp.}
{\True Đường sức quanh dây dẫn thẳng là các đường tròn đồng tâm; từ trường trong lòng ống dây dài gần đều.}
{Từ trường trong lòng ống dây dài luôn bằng không.}
\loigiai{
Đường sức quanh dây dẫn thẳng là các đường tròn đồng tâm; từ trường trong lòng ống dây dài gần đều. Vì vậy phương án $C$ đúng; các phương án còn lại trái với định nghĩa, định luật hoặc điều kiện áp dụng.
}
\end{ex}

% ===== Câu 281 =====
% ID: L12C3B14-96-TN
% Mức: TH
\begin{ex}
% ID: L12C3B14-96-TN
% Mức: TH
Nội dung nào mô tả đúng nhất từ trường của dòng điện thẳng, dòng điện tròn và ống dây?
\choice
{Có thể bỏ qua phương, chiều và đơn vị của các đại lượng trong mọi trường hợp.}
{Đại lượng đang xét không phụ thuộc các điều kiện vật lí nêu trong bài.}
{Từ trường trong lòng ống dây dài luôn bằng không.}
{\True Đường sức quanh dây dẫn thẳng là các đường tròn đồng tâm; từ trường trong lòng ống dây dài gần đều.}
\loigiai{
Đường sức quanh dây dẫn thẳng là các đường tròn đồng tâm; từ trường trong lòng ống dây dài gần đều. Vì vậy phương án $D$ đúng; các phương án còn lại trái với định nghĩa, định luật hoặc điều kiện áp dụng.
}
\end{ex}

% ===== Câu 282 =====
% ID: L12C3B14-97-TN
% Mức: TH
\begin{ex}
% ID: L12C3B14-97-TN
% Mức: TH
Khi phân tích từ trường của dòng điện thẳng, dòng điện tròn và ống dây, phát biểu nào đúng?
\choice
{\True Đường sức quanh dây dẫn thẳng là các đường tròn đồng tâm; từ trường trong lòng ống dây dài gần đều.}
{Từ trường trong lòng ống dây dài luôn bằng không.}
{Đại lượng đang xét không phụ thuộc các điều kiện vật lí nêu trong bài.}
{Có thể bỏ qua phương, chiều và đơn vị của các đại lượng trong mọi trường hợp.}
\loigiai{
Đường sức quanh dây dẫn thẳng là các đường tròn đồng tâm; từ trường trong lòng ống dây dài gần đều. Vì vậy phương án $A$ đúng; các phương án còn lại trái với định nghĩa, định luật hoặc điều kiện áp dụng.
}
\end{ex}

% ===== Câu 283 =====
% ID: L12C3B14-98-TN
% Mức: TH
\begin{ex}
% ID: L12C3B14-98-TN
% Mức: TH
Một học sinh ôn tập từ trường của dòng điện thẳng, dòng điện tròn và ống dây. Kết luận nào của học sinh là đúng?
\choice
{Từ trường trong lòng ống dây dài luôn bằng không.}
{\True Đường sức quanh dây dẫn thẳng là các đường tròn đồng tâm; từ trường trong lòng ống dây dài gần đều.}
{Có thể bỏ qua phương, chiều và đơn vị của các đại lượng trong mọi trường hợp.}
{Đại lượng đang xét không phụ thuộc các điều kiện vật lí nêu trong bài.}
\loigiai{
Đường sức quanh dây dẫn thẳng là các đường tròn đồng tâm; từ trường trong lòng ống dây dài gần đều. Vì vậy phương án $B$ đúng; các phương án còn lại trái với định nghĩa, định luật hoặc điều kiện áp dụng.
}
\end{ex}

% ===== Câu 284 =====
% ID: L12C3B14-99-TN
% Mức: TH
\begin{ex}
% ID: L12C3B14-99-TN
% Mức: TH
Chọn phát biểu không trái với kiến thức về từ trường của dòng điện thẳng, dòng điện tròn và ống dây.
\choice
{Đại lượng đang xét không phụ thuộc các điều kiện vật lí nêu trong bài.}
{Có thể bỏ qua phương, chiều và đơn vị của các đại lượng trong mọi trường hợp.}
{\True Đường sức quanh dây dẫn thẳng là các đường tròn đồng tâm; từ trường trong lòng ống dây dài gần đều.}
{Từ trường trong lòng ống dây dài luôn bằng không.}
\loigiai{
Đường sức quanh dây dẫn thẳng là các đường tròn đồng tâm; từ trường trong lòng ống dây dài gần đều. Vì vậy phương án $C$ đúng; các phương án còn lại trái với định nghĩa, định luật hoặc điều kiện áp dụng.
}
\end{ex}
